% Generated by scripts/book.py — do not edit by hand; re-run the script.
\documentclass[11pt,oneside]{book}
\usepackage{iftex}
\ifPDFTeX
  \usepackage[T1]{fontenc}
  \usepackage[utf8]{inputenc}
  \usepackage{lmodern}
\else
  \usepackage{fontspec}
  
\fi
\usepackage{amsmath,amssymb}
\usepackage{newunicodechar}
\usepackage[a4paper,margin=2cm]{geometry}
\usepackage{xcolor}
\usepackage{microtype}
\usepackage{multicol}
\usepackage{enumitem}
\usepackage{needspace}
\usepackage[most]{tcolorbox}
\usepackage{fancyhdr}
\usepackage{truncate}
\usepackage{etoolbox}
\usepackage{hyperref}
\usepackage{bookmark}

% --- portable Unicode: every glyph the corpus uses, mapped to engine-agnostic LaTeX ---
\newunicodechar{—}{---}
\newunicodechar{–}{--}
\newunicodechar{…}{\ldots{}}
\newunicodechar{§}{\S{}}
\newunicodechar{·}{\textperiodcentered{}}
\newunicodechar{×}{\ensuremath{\times}}
\newunicodechar{−}{\ensuremath{-}}
\newunicodechar{≈}{\ensuremath{\approx}}
\newunicodechar{≠}{\ensuremath{\neq}}
\newunicodechar{≥}{\ensuremath{\geq}}
\newunicodechar{≤}{\ensuremath{\leq}}
\newunicodechar{∴}{\ensuremath{\therefore}}
\newunicodechar{∧}{\ensuremath{\wedge}}
\newunicodechar{¬}{\ensuremath{\neg}}
\newunicodechar{→}{\ensuremath{\rightarrow}}
\newunicodechar{↦}{\ensuremath{\mapsto}}
\newunicodechar{⇒}{\ensuremath{\Rightarrow}}
\newunicodechar{⟹}{\ensuremath{\Longrightarrow}}
\newunicodechar{Φ}{\ensuremath{\Phi}}
\newunicodechar{Σ}{\ensuremath{\Sigma}}
\newunicodechar{ℛ}{\ensuremath{\mathcal{R}}}
\newunicodechar{𝒰}{\ensuremath{\mathcal{U}}}
\newunicodechar{′}{\ensuremath{{}^{\prime}}}
\newunicodechar{″}{\ensuremath{{}^{\prime\prime}}}
\newunicodechar{₂}{\textsubscript{2}}
\newunicodechar{ö}{\"o}
\newunicodechar{é}{\'e}
\newunicodechar{ł}{\l{}}

% --- palette ---
\definecolor{refcolor}{RGB}{120,28,32}
\definecolor{rulecolor}{RGB}{170,170,170}
\definecolor{asidebg}{RGB}{244,242,236}
\definecolor{asideline}{RGB}{193,158,103}
\hypersetup{colorlinks=true, linkcolor=refcolor, urlcolor=refcolor, citecolor=refcolor,
  breaklinks=true, pdftitle={An Encyclopedia of Philosophy}, pdfauthor={philosophy web}, bookmarksnumbered=false}
\setcounter{tocdepth}{1}

% --- two-column feel: a thin rule between the columns ---
\setlength{\columnseprule}{0.4pt}
\renewcommand{\columnseprulecolor}{\color{rulecolor}}
\setlength{\columnsep}{1.6em}
\raggedbottom
\setlist{leftmargin=1.3em, itemsep=2pt, topsep=3pt, parsep=0pt}

% --- running head: book title (verso) and current entry as a guide word (recto) ---
\pagestyle{fancy}
\fancyhf{}
\renewcommand{\headrulewidth}{0.3pt}
% long argument/claim headwords are truncated at a word boundary (…) so the guide word fits
\fancyhead[LE]{\small\scshape An Encyclopedia of Philosophy}
\fancyhead[RO]{\small\scshape\truncate{0.7\headwidth}{\rightmark}}
\fancyhead[LO,RE]{\small\thepage}
\fancypagestyle{plain}{\fancyhf{}\renewcommand{\headrulewidth}{0pt}\fancyfoot[C]{\small\thepage}}

% --- letter / group divider: full-width band on a fresh page ---
\newcommand{\bookgroup}[1]{%
  \clearpage
  \phantomsection\addcontentsline{toc}{chapter}{#1}%
  \begin{center}\rule{0.5\linewidth}{0.4pt}\\[4pt]
    {\Huge\bfseries #1}\\[3pt]\rule{0.5\linewidth}{0.4pt}\end{center}
  \vspace{0.6\baselineskip}\markboth{#1}{#1}}

% --- type glyph: a small, dimmed mark of the entry's kind (its LaTeX comes from glyphs.py) ---
\newcommand{\entryglyph}[1]{{\normalfont\footnotesize\color{black!45}#1}}

% --- an entry: anchor + PDF bookmark + headword + optional subtitle + etiquette line ---
% #1 anchor id · #2 rich headword (body + ToC) · #3 etiquette · #4 plain guide word (mark) ·
% #5 subtitle (the full title when it differs from the headword; empty = omitted) · #6 type glyph
\newcommand{\entryhead}[6]{%
  \par\addvspace{0.9\baselineskip}\needspace{4\baselineskip}%
  \phantomsection\label{entry:#1}\addcontentsline{toc}{section}{#2}%
  \markright{#4}%
  {\raggedright\large\bfseries\entryglyph{#6}\hspace{0.4em}#2\par}%
  \ifstrempty{#5}{}{{\nopagebreak\smallskip\raggedright\itshape\small\color{black!75} #5\par}}%
  {\nopagebreak\smallskip\footnotesize\itshape\color{black!65} #3\par}%
  \nopagebreak\vspace{2pt}}

% --- a `### subheading` within an entry: small-caps run-in ---
\newcommand{\webheading}[1]{\par\addvspace{5pt}\noindent{\normalfont\scshape #1}\par\nopagebreak\smallskip}

% --- `### In plain terms` -> shaded aside that can break across columns/pages ---
\newtcolorbox{plainterms}{breakable, enhanced, boxrule=0pt, frame hidden,
  colback=asidebg, borderline west={2pt}{0pt}{asideline},
  left=8pt, right=6pt, top=5pt, bottom=5pt, before skip=7pt, after skip=7pt,
  fontupper=\small}

% --- `### See also` list ---
\newlist{seealso}{itemize}{1}
\setlist[seealso]{label=\textcolor{refcolor}{\small$\triangleright$}, leftmargin=1.4em,
  itemsep=2pt, topsep=2pt, parsep=0pt}

% --- front matter: the Symbols key, a "list of symbols" two-column table (glyph -> meaning) ---
\newcommand{\symbolkeyheading}{\thispagestyle{plain}\markboth{}{}%
  \vspace*{1.4\baselineskip}{\centering\Large\bfseries Symbols\par}\vspace{1.3\baselineskip}}
\newenvironment{symbolkey}
  {\par\begin{center}\begin{tabular}{r@{\hskip 1.6em}l}}
  {\end{tabular}\end{center}\clearpage}
% #1 = glyph (right-aligned, full strength) · #2 = meaning (left-aligned)
\newcommand{\symbolrow}[2]{#1 & #2\\[8pt]}

% --- front matter: a List of Entries row (#1 type glyph · #2 anchor id · #3 name) ---
\newcommand{\entryline}[3]{\noindent\makebox[1.5em][l]{\entryglyph{#1}}%
  \hyperref[entry:#2]{#3}\,\dotfill\,\pageref{entry:#2}\par}
\newenvironment{listofentries}{\footnotesize\begin{multicols}{2}\raggedright\setlength{\parskip}{2pt}}
  {\end{multicols}\clearpage}

% --- back matter: Sources (bibliography of every cited source, hanging indent) ---
\newcommand{\sourcesband}{%
  \clearpage
  \phantomsection\addcontentsline{toc}{chapter}{Sources}%
  \begin{center}\rule{0.5\linewidth}{0.4pt}\\[4pt]
    {\Huge\bfseries Sources}\\[3pt]\rule{0.5\linewidth}{0.4pt}\end{center}
  \vspace{0.6\baselineskip}\markboth{Sources}{Sources}}
\newenvironment{sourcelist}{\small\setlength{\parindent}{0pt}\raggedright\sloppy}{\par}
% #1 = anchor id (empty for a second/third work under the same source node) · #2 = the reference
\newcommand{\sourceitem}[2]{%
  \par\addvspace{3pt}\ifstrempty{#1}{}{\phantomsection\label{source:#1}}%
  \hangindent=1.5em\hangafter=1 #2\par}

\begin{document}
\frontmatter
\thispagestyle{empty}
\begin{titlepage}\centering
  \vspace*{0.20\textheight}
  {\Huge\bfseries An Encyclopedia of Philosophy\par}
  \vspace{1.2\baselineskip}
  {\large\itshape Distilled from the philosophy web\par}
  \vfill
  {\small June 2026\par}
\end{titlepage}
\symbolkeyheading
\begin{symbolkey}
\symbolrow{$\blacklozenge$}{\emph{concept} — a term and its definition}
\symbolrow{$\odot$}{\emph{claim} — a single truth-apt proposition}
\symbolrow{$\vdash$}{\emph{argument} — premises leading to a conclusion}
\symbolrow{$\bigstar$}{\emph{position} — a named standpoint or view}
\symbolrow{?}{\emph{question} — an open issue under debate}
\end{symbolkey}
\markboth{}{}\thispagestyle{plain}
\vspace*{1.2\baselineskip}
{\centering\Large\bfseries List of Entries\par}
{\centering\small 72 entries \,\textperiodcentered\, June 2026\par}
\vspace{1\baselineskip}
\begin{listofentries}
\entryline{$\odot$}{claim-aposteriori-necessities-are-misdescribed-possibilities}{A-posteriori necessities as misdescribed possibilities}
\entryline{$\vdash$}{argument-conceivability-not-guide-to-possibility}{A-posteriori-necessity objection}
\entryline{$\vdash$}{argument-ability-hypothesis}{Ability hypothesis}
\entryline{$\odot$}{claim-no-nonarbitrary-threshold-for-experience}{Absence of an experiential threshold}
\entryline{$\odot$}{claim-strong-emergence-is-acceptable}{Acceptability of strong emergence}
\entryline{$\blacklozenge$}{concept-access-consciousness}{Access consciousness}
\entryline{$\blacklozenge$}{concept-access-vs-phenomenal-consciousness}{Access vs. phenomenal consciousness}
\entryline{$\vdash$}{argument-acquaintance-hypothesis}{Acquaintance hypothesis}
\entryline{$\vdash$}{argument-bike-rider-learns-facts}{Bike-rider rejoinder}
\entryline{$\odot$}{claim-consciousness-categorical-not-dispositional}{Categorical, not dispositional}
\entryline{$\vdash$}{argument-causal-argument-for-physicalism}{Causal argument for physicalism}
\entryline{$\blacklozenge$}{concept-causal-closure}{Causal closure of the physical}
\entryline{$\vdash$}{argument-causal-exclusion}{Causal exclusion argument}
\entryline{$\blacklozenge$}{concept-chinese-nation}{Chinese nation}
\entryline{$\blacklozenge$}{concept-chinese-room}{Chinese room}
\entryline{$\vdash$}{argument-cognitive-closure}{Cognitive closure}
\entryline{$\odot$}{claim-phenomenal-intensions-coincide}{Coincident phenomenal intensions}
\entryline{$\vdash$}{argument-combination-problem}{Combination problem}
\entryline{$\blacklozenge$}{concept-completeness-of-physics}{Completeness of physics}
\entryline{$\odot$}{claim-conceivability-entails-possibility}{Conceivability entails possibility}
\entryline{$\vdash$}{argument-no-arbitrary-line-panpsychism}{Continuity argument}
\entryline{$\bigstar$}{position-cosmopsychism}{Cosmopsychism}
\entryline{$\odot$}{claim-cosmos-is-fundamental-subject}{Cosmos as fundamental subject}
\entryline{$\vdash$}{argument-strong-emergence-acceptable}{Emergence reply}
\entryline{$\blacklozenge$}{concept-ai-hard-problem-etiologies}{Etiologies of AI hard-problem discourse}
\entryline{$\odot$}{claim-irreducible-mental-causally-excluded}{Exclusion of the irreducible mental}
\entryline{$\odot$}{claim-consciousness-fundamental-nonphysical}{Fundamental non-physical consciousness}
\entryline{$\odot$}{claim-consciousness-fundamental-ubiquitous}{Fundamental ubiquitous consciousness}
\entryline{$\vdash$}{argument-genetic-panpsychism}{Genetic argument}
\entryline{$\odot$}{claim-subject-unity-graded-by-integration}{Graded subject unity}
\entryline{$\blacklozenge$}{concept-hard-problem}{Hard problem}
\entryline{?}{question-hard-problem}{Hard problem (the question)}
\entryline{$\vdash$}{argument-indexical-reply}{Indexical reply}
\entryline{$\odot$}{claim-knowhow-is-knowledge-that}{Know-how as knowledge-that}
\entryline{$\blacklozenge$}{concept-knowledge-argument}{Knowledge argument}
\entryline{$\vdash$}{argument-knowledge-argument}{Knowledge argument (Mary's room)}
\entryline{$\odot$}{claim-mary-gains-abilities-not-facts}{Mary gains abilities}
\entryline{$\odot$}{claim-mary-gains-acquaintance}{Mary gains acquaintance}
\entryline{$\odot$}{claim-mary-gains-old-fact-new-guise}{Mary gains an old fact's new guise}
\entryline{$\odot$}{claim-mary-gains-self-locating-knowledge}{Mary gains self-locating knowledge}
\entryline{$\odot$}{claim-mary-knows-all-physical}{Mary knows all physical facts}
\entryline{$\odot$}{claim-mary-learns-nothing-new}{Mary learns nothing new}
\entryline{$\odot$}{claim-mary-learns-on-release}{Mary learns on release}
\entryline{$\vdash$}{argument-mary-misimagined}{Mary misimagined}
\entryline{$\vdash$}{argument-mary-omniscience-unrealizable}{Mary's omniscience unrealizable}
\entryline{$\odot$}{claim-micro-subjects-do-not-sum}{Micro-subjects do not sum}
\entryline{$\vdash$}{argument-modal-rationalism}{Modal rationalism}
\entryline{$\odot$}{claim-no-radical-emergence}{No radical emergence}
\entryline{$\odot$}{claim-some-truths-about-experience-not-physically-deducible}{Non-deducible experiential truths}
\entryline{$\vdash$}{argument-old-fact-new-guise}{Old fact, new guise}
\entryline{$\blacklozenge$}{concept-panpsychism}{Panpsychism}
\entryline{$\bigstar$}{position-panpsychism}{Panpsychism (position)}
\entryline{$\blacklozenge$}{concept-paradox-of-phenomenal-judgment}{Paradox of phenomenal judgment}
\entryline{$\blacklozenge$}{concept-phenomenal-consciousness}{Phenomenal consciousness}
\entryline{$\odot$}{claim-phenomenal-residue}{Phenomenal residue}
\entryline{$\blacklozenge$}{concept-philosophical-zombie}{Philosophical zombie}
\entryline{$\bigstar$}{position-physicalism}{Physicalism}
\entryline{$\blacklozenge$}{concept-primary-intension}{Primary intension}
\entryline{$\blacklozenge$}{concept-qualia}{Qualia}
\entryline{$\odot$}{claim-phenomenal-consciousness-is-real}{Reality of phenomenal consciousness}
\entryline{$\bigstar$}{position-russellian-monism}{Russellian monism}
\entryline{$\blacklozenge$}{concept-secondary-intension}{Secondary intension}
\entryline{$\odot$}{claim-functional-description-structural-dynamical-only}{Structure and dynamics only}
\entryline{$\odot$}{claim-subjects-actually-merge-and-divide}{Subjects merge and divide}
\entryline{$\blacklozenge$}{concept-materialism-types}{Type-A / Type-B / Type-C materialism}
\entryline{$\bigstar$}{position-type-a-materialism}{Type-A materialism}
\entryline{$\bigstar$}{position-type-b-materialism}{Type-B materialism}
\entryline{$\bigstar$}{position-type-c-materialism}{Type-C materialism}
\entryline{$\blacklozenge$}{concept-introspective-gap-varieties}{Varieties of introspective explanatory gap}
\entryline{?}{question-what-does-mary-learn}{What Mary learns}
\entryline{$\vdash$}{argument-zombie-conceivability}{Zombie argument}
\entryline{$\odot$}{claim-zombie-conceivable}{Zombie conceivability}
\end{listofentries}
\mainmatter
\bookgroup{A}
\begin{multicols}{2}
\entryhead{claim-aposteriori-necessities-are-misdescribed-possibilities}{A-posteriori necessities as misdescribed possibilities}{Claim \,\textperiodcentered\, after Chalmers}{A-posteriori necessities as misdescribed possibilities}{Kripkean a-posteriori necessities are misdescribed genuine possibilities, not conceived impossibilities}{$\odot$}
Kripkean a-posteriori necessities are not cases in which ideal conceiving reaches a genuine impossibility; they are genuine \emph{possibilities} described in a misleading way. This is the two-dimensional diagnosis of the standard counterexamples to ``conceivability entails possibility'' (\textsc{A-posteriori necessities}, from Kripke's \emph{Naming and Necessity} \hyperref[source:source-kripke-1980]{\textsc{Kripke — Naming and Necessity (1980)}}), defended by David Chalmers (\hyperref[source:source-chalmers-1996]{\textsc{Chalmers — The Conscious Mind (1996)}}; \hyperref[source:source-chalmers-2002]{\textsc{Chalmers — On Sense and Intension (2002)}}).

When one ``conceives water that is not H₂O,'' one conceives a \emph{genuinely possible} scenario — some watery stuff that is not H₂O — and merely \emph{misdescribes} it as ``water.'' Distinguish a term's \emph{primary} (epistemic) \hyperref[entry:concept-primary-intension]{\textsc{Primary intension}} from its \emph{secondary} (subjunctive) \hyperref[entry:concept-secondary-intension]{\textsc{Secondary intension}}: the conceived scenario is possible under the primary intension, while the necessity attaches only to the secondary. So the a-posteriori necessities are not impossibilities reached by ideal conceiving — they are possibilities wearing a misleading label.

This is a claim about the \emph{modal status} of the standard counterexamples, not a denial that they are necessary. Its payoff is to preserve the link between \textbf{ideal primary conceivability} and \textbf{primary possibility} by explaining away every apparent defeater as a description error — the engine of \hyperref[entry:argument-modal-rationalism]{\textsc{Modal rationalism}}.

The position that denies it is type-B physicalism (see \hyperref[entry:concept-materialism-types]{\textsc{Type-A / Type-B / Type-C materialism}}), which holds that some a-posteriori necessities are \emph{strong}: no possible world even verifies the primary intension, so there is nothing being misdescribed — just an impossibility conceived. The claim is distinct from \textsc{A-posteriori necessities}, which it reinterprets rather than denies: that claim says such necessities exist; this one says they do not break the conceivability→possibility link once the two intensions are distinguished.

\webheading{See also}
\begin{seealso}
  \item \textsc{A-posteriori necessities} — the Kripkean necessities this claim reinterprets.
  \item \hyperref[entry:concept-primary-intension]{\textsc{Primary intension}} / \hyperref[entry:concept-secondary-intension]{\textsc{Secondary intension}} — the distinction that turns an apparent impossibility into a misdescribed possibility.
  \item \hyperref[entry:argument-modal-rationalism]{\textsc{Modal rationalism}} — the argument this claim is the load-bearing premise of.
  \item \hyperref[entry:concept-materialism-types]{\textsc{Type-A / Type-B / Type-C materialism}} — type-B physicalism, which resists the diagnosis by positing \emph{strong} necessities.
\end{seealso}

\entryhead{argument-conceivability-not-guide-to-possibility}{A-posteriori-necessity objection}{Argument}{A-posteriori-necessity objection}{Conceivability is not a reliable guide to possibility (the a-posteriori-necessity objection)}{$\vdash$}
Conceivability is not, in general, a reliable guide to metaphysical possibility — so the bridge the zombie argument leans on cannot simply be assumed. The objection turns a celebrated result in the metaphysics of modality against the \hyperref[entry:argument-zombie-conceivability]{\textsc{Zombie argument}}: since Saul Kripke's work on naming, philosophers have widely accepted that some truths are \emph{necessary} yet knowable only \emph{a posteriori}, and in every such case a coherent-seeming scenario is conceivable even though it is metaphysically impossible.

The move is a counterexample to a general principle:

\begin{enumerate}
  \item (\textsc{A-posteriori necessities}, from Kripke's \emph{Naming and Necessity} \hyperref[source:source-kripke-1980]{\textsc{Kripke — Naming and Necessity (1980)}}) There are necessary truths knowable only a posteriori — water is H₂O, heat is molecular motion, Hesperus is Phosphorus — whose negations are conceivable yet metaphysically impossible.
\end{enumerate}

∴ Conceivability does not in general entail possibility: \hyperref[entry:claim-conceivability-entails-possibility]{\textsc{Conceivability entails possibility}} fails as a universal bridge.

Applied to consciousness, this is the \textbf{type-B physicalist} strategy (see \hyperref[entry:concept-materialism-types]{\textsc{Type-A / Type-B / Type-C materialism}}): a phenomenal–physical identity (``pain \emph{is} such-and-such neural state'') could itself be a \emph{strong} a posteriori necessity. If so, a \hyperref[entry:concept-philosophical-zombie]{\textsc{Philosophical zombie}} is conceivable but metaphysically impossible, and the \hyperref[entry:argument-zombie-conceivability]{\textsc{Zombie argument}} is cut at its bridge — without conceding any non-physical fact.

\webheading{The crux: is consciousness like water?}

The \textbf{weakest premise} is the tacit one: that the \emph{phenomenal} case is relevantly like the \emph{water} case. The two-dimensional rejoinder (\hyperref[entry:argument-modal-rationalism]{\textsc{Modal rationalism}}, built on \hyperref[entry:claim-aposteriori-necessities-are-misdescribed-possibilities]{\textsc{A-posteriori necessities as misdescribed possibilities}}) replies that every Kripkean necessity is a \emph{misdescribed possibility}, and that for phenomenal terms the primary and secondary intensions \emph{coincide} — there is no appearance/reality wedge to exploit — so the water-style escape is unavailable for zombies. Whether the phenomenal is genuinely disanalogous to water is where the dispute lives.

This should not be confused with \textsc{Phenomenal-concept deflation}: that explains away the \emph{appearance} of an explanatory gap through the architecture of phenomenal concepts; the present move makes the purely \emph{modal} point that a-posteriori necessities already break conceivability→possibility, independently of any account of phenomenal concepts.

\begin{plainterms}
The whole zombie argument needs the rule ``if you can coherently imagine it, it's really possible.'' This objection says that rule is just false. Take ``water is H₂O'': for centuries people could coherently imagine water turning out to be something other than H₂O — yet it is impossible, because water \emph{just is} H₂O, a fact we had to discover rather than deduce. So imaginability isn't proof of possibility. Maybe consciousness is the same: a zombie feels imaginable only because we haven't yet seen that experience \emph{just is} some physical process. The standard comeback, due to David Chalmers (see \hyperref[source:source-chalmers-1996]{\textsc{Chalmers — The Conscious Mind (1996)}}), is that the water trick works only because ``water'' hides a surface/essence gap — and feelings have no such gap, so the trick cannot be repeated for zombies.
\end{plainterms}

\webheading{See also}
\begin{seealso}
  \item \textsc{A-posteriori necessities} — the premise: Kripkean necessities whose negations are conceivable yet impossible.
  \item \hyperref[entry:claim-conceivability-entails-possibility]{\textsc{Conceivability entails possibility}} — the bridge this argument attacks.
  \item \hyperref[entry:argument-zombie-conceivability]{\textsc{Zombie argument}} — the argument severed if the objection holds.
  \item \hyperref[entry:argument-modal-rationalism]{\textsc{Modal rationalism}} — the two-dimensional rejoinder that the Kripke cases are misdescribed possibilities, not genuine counterexamples.
  \item \hyperref[entry:claim-aposteriori-necessities-are-misdescribed-possibilities]{\textsc{A-posteriori necessities as misdescribed possibilities}} — the claim that rejoinder rests on.
  \item \hyperref[entry:concept-materialism-types]{\textsc{Type-A / Type-B / Type-C materialism}} — type-B physicalism, the position this objection arms.
  \item \hyperref[entry:concept-philosophical-zombie]{\textsc{Philosophical zombie}} — the case alleged to be conceivable-but-impossible.
  \item \textsc{Phenomenal-concept deflation} — the allied but distinct phenomenal-concept strategy.
\end{seealso}

\entryhead{argument-ability-hypothesis}{Ability hypothesis}{Argument \,\textperiodcentered\, after Lewis}{Ability hypothesis}{The Ability Hypothesis — experience teaches know-how, so Mary learns no new fact}{$\vdash$}
The \emph{Ability Hypothesis} is David Lewis's reply to the \hyperref[entry:argument-knowledge-argument]{\textsc{Knowledge Argument}}, also defended by Laurence Nemirow. It holds that knowing what an experience is like is not a piece of factual knowledge at all, but a bundle of \emph{abilities}: to recognise the experience, to imagine it, and to remember it. If that is right, then when Mary the colour scientist, who knows every physical fact about colour vision but has never seen red, finally sees it, she gains know-how rather than a new fact, and the inference from her learning something to the falsity of physicalism is blocked.

\webheading{The argument}

The attack lands on the Knowledge Argument's verdict step: the critical question (CQ2 of the thought-experiment scheme) of how to \emph{read} the agreed intuition that Mary learns something. The intuition is robust, but reading it as \emph{acquiring a new propositional fact} is theory-laden, and the verdict survives under a deflationary reading that the original argument cannot use. Given an analysis of ``knowing what it is like'', the inference is deductive:

\begin{enumerate}
  \item Knowing what an experience is like consists in possessing certain abilities: to recognise the experience when it occurs, to imagine it, and to remember it (Lewis \hyperref[source:source-lewis-1988]{\textsc{1988}}; Nemirow \hyperref[source:source-nemirow-1990]{\textsc{1990}}).
  \item On first seeing red, what Mary gains is exactly this: knowing what seeing red is like.
  \item Abilities are not propositional knowledge; gaining one is acquiring a skill, not learning a truth.
\end{enumerate}

∴ \hyperref[entry:claim-mary-gains-abilities-not-facts]{\textsc{Mary gains no new fact}}. So the step from her epistemic progress to \hyperref[entry:claim-some-truths-about-experience-not-physically-deducible]{\textsc{truths about experience that no physical description yields}} is blocked: there is no \emph{truth} she lacked.

\webheading{What it concedes}

The move is notable for conceding both premises of the Knowledge Argument as written. It grants \hyperref[entry:claim-mary-knows-all-physical]{\textsc{that Mary knew all the physical facts}} outright, and \hyperref[entry:claim-mary-learns-on-release]{\textsc{that she learns something on release}} under the know-how reading. What it severs is the inference between them. ``Mary learns something'' is true; ``Mary learns a new fact'' is what the conclusion needs; and the case cannot fund the stronger reading without begging the question against physicalism.

\webheading{The weak point}

The most disputed premise is the first, the ability analysis itself. Three counters are standard. Mary seems able to \emph{use} her new knowledge in reasoning and inference (``seeing orange must be like \emph{this}, only more yellowish''), which bare abilities do not explain. One can apparently have the abilities without the knowledge (a distracted imaginer who can summon red but is not thereby informed), and arguably the knowledge while lacking the imaginative ability (someone paralysed of imagination). And the acquaintance theorist offers a rival non-factual analysis, \hyperref[entry:claim-mary-gains-acquaintance]{\textsc{that what Mary gains is acquaintance rather than a fact}}, so even granting that no new fact is learned, the ability analysis is not the only option.

\begin{plainterms}
This is the most famous escape from the Mary story (the \hyperref[entry:argument-knowledge-argument]{\textsc{Knowledge Argument}}: since Mary, who knew all the physics of colour, still learns something when she first sees red, physics must be incomplete). The escape turns on an ambiguity in ``learning''. Learning a \emph{fact} (that Paris is the capital of France) is one thing; picking up a \emph{skill} (juggling) is another. On this view Mary picks up skills, namely recognising, imagining, and remembering red. No fact was missing from her books, so nothing was missing from physics. The catch: her new grasp of red seems usable in actual reasoning, the way facts are and skills are not, so perhaps it is not \emph{just} a skill.
\end{plainterms}

\webheading{See also}
\begin{seealso}
  \item \hyperref[entry:argument-knowledge-argument]{\textsc{Knowledge argument (Mary's room)}} — the argument this one attacks: that Mary learns a fact the physical description omits.
  \item \hyperref[entry:concept-knowledge-argument]{\textsc{Knowledge argument}} — the Mary thought experiment itself.
  \item \hyperref[entry:claim-mary-gains-abilities-not-facts]{\textsc{Mary gains abilities}} — the conclusion: Mary acquires abilities, not a new fact.
  \item \hyperref[entry:claim-mary-gains-acquaintance]{\textsc{Mary gains acquaintance}} — the rival non-factual reply: acquaintance rather than know-how.
  \item \hyperref[entry:claim-some-truths-about-experience-not-physically-deducible]{\textsc{Non-deducible experiential truths}} — the anti-physicalist conclusion this attack is meant to block.
  \item \hyperref[source:source-lewis-1988]{\textsc{Lewis — What Experience Teaches (1988)}} — Lewis's ``What Experience Teaches'' (1988), the primary statement.
  \item \hyperref[source:source-nemirow-1990]{\textsc{Nemirow — Physicalism and the Cognitive Role of Acquaintance (1990)}} — Nemirow's independent defence of the ability analysis.
\end{seealso}

\entryhead{claim-no-nonarbitrary-threshold-for-experience}{Absence of an experiential threshold}{Claim \,\textperiodcentered\, after Nagel}{Absence of an experiential threshold}{There is no non-arbitrary point in nature at which experience first appears}{$\odot$}
If experience were confined to brains, or to life, or to systems above some threshold of complexity, there ought to be a principled place where it first switches on — a line such that below it there is nothing it is like to be the system and above it there is. No such non-arbitrary line is in view. The relevant differences across any proposed boundary look like differences of degree and organisation, not the kind of difference that could mark the abrupt entry of an entirely new ontological feature into the world.

This is the premise of the \hyperref[entry:argument-no-arbitrary-line-panpsychism]{\textsc{continuity argument}} for panpsychism: absent a principled threshold, the parsimonious option is that experience, in some attenuated form, reaches all the way down. It is denied by emergentists and physicalists who hold there \emph{is} a principled threshold — a level of integrated information or of functional organisation at which experience genuinely, if strongly, emerges (see \hyperref[entry:claim-strong-emergence-is-acceptable]{\textsc{Acceptability of strong emergence}}). It differs from \hyperref[entry:claim-no-radical-emergence]{\textsc{No radical emergence}}: that claim says experience could not brutely \emph{appear} at all from the non-experiential, whereas this one says that even granting it could, there is no non-arbitrary \emph{place} to locate its first appearance.

\webheading{See also}
\begin{seealso}
  \item \hyperref[entry:argument-no-arbitrary-line-panpsychism]{\textsc{Continuity argument}} — the argument this premise drives.
  \item \hyperref[entry:claim-strong-emergence-is-acceptable]{\textsc{Acceptability of strong emergence}} — its denial: there \emph{is} a principled emergence threshold.
  \item \hyperref[entry:claim-no-radical-emergence]{\textsc{No radical emergence}} — the neighbouring premise about \emph{whether}, not \emph{where}, experience could emerge.
\end{seealso}

\entryhead{claim-strong-emergence-is-acceptable}{Acceptability of strong emergence}{Claim}{Acceptability of strong emergence}{Strong emergence of experience from non-experiential matter is acceptable}{$\odot$}
Experience can arise as a genuinely novel, fundamental feature at some level of physical organisation — a brain, say — without being present, even in proto-form, in that organisation's parts. On this view \emph{strong} emergence (the appearance of properties not grounded in the powers of the base) is a coherent and even expected feature of a layered natural world, not the magic its critics allege. The dispute is specifically about strong emergence of the \emph{phenomenal}; \emph{weak} emergence, where higher-level patterns are fully grounded in and predictable from the base, is accepted by almost everyone and is not at issue here.

This is the standard physicalist denial of \hyperref[entry:claim-no-radical-emergence]{\textsc{No radical emergence}}, and so the load-bearing premise of the \hyperref[entry:argument-strong-emergence-acceptable]{\textsc{emergence reply}} that undercuts the genetic argument for panpsychism: if strong emergence is acceptable, the genetic argument's choice between fundamental experience and unintelligible magic is a false dilemma. It is denied by panpsychists, who hold that strong emergence of the phenomenal is unintelligible (Strawson), and by those who think calling experience ``strongly emergent'' merely relabels the explanatory gap rather than closing it.

\webheading{See also}
\begin{seealso}
  \item \hyperref[entry:claim-no-radical-emergence]{\textsc{No radical emergence}} — the claim this one denies.
  \item \hyperref[entry:argument-strong-emergence-acceptable]{\textsc{Emergence reply}} — the objection this premise drives against the genetic argument.
  \item \hyperref[entry:argument-genetic-panpsychism]{\textsc{Genetic argument}} — the argument that the acceptability of strong emergence is meant to undercut.
\end{seealso}

\entryhead{concept-access-consciousness}{Access consciousness}{Concept \,\textperiodcentered\, after Block}{Access consciousness}{}{$\blacklozenge$}
\emph{Access consciousness} is the functional sense of ``conscious'': a mental state is access-conscious when its content is \emph{available} to the rest of the mind, poised for use in reasoning, for the direct rational control of action, and for verbal report. The term is Ned Block's (he abbreviates it \emph{A-consciousness}). It is a matter of what a system can \emph{do with} a state's content, settled entirely by the state's role in the cognitive economy and not at all by how, or whether, it feels.

A simple test fixes the idea: when you read this sentence, its meaning is poised for you to report, question, and reason with, and that availability is the sentence's access-consciousness. Theories of consciousness in the \emph{global workspace} family (see \textsc{Global workspace theory}) effectively identify being conscious with being access-conscious in this sense, broadcast widely enough across the system to be used.

Access consciousness is defined by contrast with \hyperref[entry:concept-phenomenal-consciousness]{\textsc{Phenomenal consciousness}}, the felt or ``what-it-is-like'' side of experience. Block's central claim is that the two can come apart in both directions: a system might have the availability with no accompanying feel, and experience might even \emph{overflow} what gets globally broadcast. That dissociation is set out at \hyperref[entry:concept-access-vs-phenomenal-consciousness]{\textsc{Access vs. phenomenal consciousness}}. Keeping the two apart is the standard guard against treating a system's fluent self-report, an access-style capacity, as evidence that it \emph{feels} anything.

\begin{plainterms}
Calling a thought ``access-conscious'' just means the information is \emph{available} to the rest of your mind, so that you can act on it, talk about it, and reason with it. It says nothing about whether the thought \emph{feels} like anything. That is the whole point of the label: it lets us separate ``the system can use and report this'' from ``there's something it's like to have it'', so we don't slide from a machine describing its own inner states to assuming it must therefore feel them.
\end{plainterms}

\webheading{See also}
\begin{seealso}
  \item \hyperref[entry:concept-phenomenal-consciousness]{\textsc{Phenomenal consciousness}} — the felt sense of consciousness it is contrasted with.
  \item \hyperref[entry:concept-access-vs-phenomenal-consciousness]{\textsc{Access vs. phenomenal consciousness}} — Block's distinction and the case that the two dissociate in both directions.
  \item \textsc{Global workspace theory} — a theory that ties consciousness to access (global broadcast).
  \item \textsc{First-person access} — privileged first-person access, a distinct, epistemic notion.
\end{seealso}

\entryhead{concept-access-vs-phenomenal-consciousness}{Access vs. phenomenal consciousness}{Concept \,\textperiodcentered\, after Block}{Access vs. phenomenal consciousness}{Access consciousness vs. phenomenal consciousness}{$\blacklozenge$}
\emph{Access consciousness} and \emph{phenomenal consciousness} are two distinct things the single word ``consciousness'' can mean, separated by Ned Block (``On a Confusion about a Function of Consciousness'', 1995). A state is \emph{access-conscious} when its content is available to the rest of the mind — poised for use in reasoning, for the direct rational control of action, and for verbal report. A state is \emph{phenomenally conscious} when there is something it is like to be in it — when it carries a felt, qualitative character. The first is a matter of what a system can \emph{do with} a state's content; the second, of how the state \emph{feels} from the inside.

\webheading{The two senses}

\hyperref[entry:concept-access-consciousness]{\textsc{Access consciousness}} (\emph{A-consciousness}) is a functional, information-availability notion: a content is A-conscious to the degree that it is \emph{poised} to be drawn on by the systems that reason, decide, and speak. Whether a state qualifies is settled entirely by the role its content plays in the cognitive economy — what downstream processes can reach and use it — and by nothing about how it feels. When you read this sentence, its meaning is poised for report and inference; that availability is its access side.

\emph{Phenomenal consciousness} (P-consciousness) is the ``what-it-is-like'' aspect of a mental state: the redness of red, the hurt of a pain, the timbre of a heard note. It is the felt or qualitative character of experience, and it is the property at issue in the \hyperref[entry:concept-hard-problem]{\textsc{hard problem}} of consciousness — the question of why any physical process should be accompanied by felt experience at all. The particular look of these black marks on the page, over and above your ability to report and reason about the words, is their phenomenal side.

\webheading{Why the distinction matters}

The point of drawing the line is that the two can come apart, at least in principle, in \emph{both} directions. A system might carry richly reportable models of its own inner states — answering fluently about what it is ``noticing'' or ``feeling'' — while there is nothing it is like to be that system at all: access without any accompanying feel. Block argues the converse is possible too, that phenomenal consciousness may \emph{overflow} access, so that experience holds more than ever gets globally broadcast for report. The canonical illustration is a brief flash of a full array of letters: one seems to see the whole array vividly, yet can read out only a few before the image fades — as though the felt richness outruns what access captures. So ``has A'' neither entails nor is entailed by ``has P.''

This independence is the pivot on which several moves in the debate over machine introspection turn. It marks the central danger: mistaking a system's \emph{functional, access-style} capacity — its ability to report on its own states, as some current language models do to a limited and unreliable degree (see \textsc{LLM functional introspection}) — for evidence of a \emph{phenomenal} inner life. It also yields a subtler observation: \emph{posing} the hard problem from the first-person side is itself an access-dependent act, requiring access to a representation of one's own phenomenal-seeming states (see \textsc{Posing the gap requires access}), even though merely \emph{having} phenomenal states requires no such access.

\begin{plainterms}
``Conscious'' runs together two different things. One is \textbf{access}: the information is available to the rest of your mind — you can reason with it, act on it, and say it out loud. The other is \textbf{feel}: there's something it's like to undergo the state (the redness of red, the hurt of pain). They're not the same. You could imagine a machine that has all the \emph{access} — it reports its inner states fluently — with no \emph{feel} at all; and some philosophers think the reverse can happen too, that we feel more than we can ever report. Keeping the two apart is the main guard against a tempting mistake: seeing a system talk about its own insides and concluding it must therefore \emph{feel} something.
\end{plainterms}

\webheading{See also}
\begin{seealso}
  \item \hyperref[entry:concept-access-consciousness]{\textsc{Access consciousness}} — the functional half of the distinction, treated on its own.
  \item \hyperref[entry:concept-phenomenal-consciousness]{\textsc{Phenomenal consciousness}} — the felt half of the distinction, treated on its own.
  \item \hyperref[entry:concept-hard-problem]{\textsc{Hard problem}} — the felt, phenomenal side is exactly what this problem is about; access is precisely what it is \emph{not}.
  \item \textsc{LLM functional introspection} — a candidate case of access-style self-report in machines, and the prime spot where it must not be read as the phenomenal kind.
  \item \textsc{Posing the gap requires access} — that stating the hard problem is itself an access-dependent act, even though undergoing phenomenal states is not.
\end{seealso}

\entryhead{argument-acquaintance-hypothesis}{Acquaintance hypothesis}{Argument \,\textperiodcentered\, after Conee}{Acquaintance hypothesis}{The Acquaintance Hypothesis — Mary meets the experience rather than learning a truth}{$\vdash$}
The \emph{acquaintance hypothesis} answers the \hyperref[entry:argument-knowledge-argument]{\textsc{knowledge argument}} by reclassifying what Mary gains: the felt verdict that ``Mary learns something'' on first seeing red is genuine, but it tracks a change of \emph{relation}, not a change in \emph{information}. The move is deductive, given a taxonomy of knowledge. It raises thought-experiment critical question CQ2 (verdict) against the knowledge argument from a different angle than the \hyperref[entry:argument-ability-hypothesis]{\textsc{ability hypothesis}}, which reads the same verdict as a gain of skill.

\webheading{The inference}

\begin{enumerate}
  \item Knowledge comes in at least three irreducible kinds: knowledge of facts, know-how, and \emph{knowledge by acquaintance} — a maximally direct cognitive relation to a thing (Earl Conee, \hyperref[source:source-conee-1994]{\textsc{Conee — Phenomenal Knowledge (1994)}}; the category descends from Bertrand Russell's distinction between knowledge by acquaintance and knowledge by description, \hyperref[source:source-russell-1911]{\textsc{Russell — Knowledge by Acquaintance and Knowledge by Description (1911)}}).
  \item By stipulation (see \hyperref[entry:claim-mary-knows-all-physical]{\textsc{Mary knows all the physical facts}}) Mary lacks no physical fact; what she lacks before release is acquaintance with the experience of red.
  \item Coming to be acquainted with a thing is not coming to know a new truth about it — one can be fully informed about a thing one has never encountered.
\end{enumerate}

Therefore \hyperref[entry:claim-mary-gains-acquaintance]{\textsc{her post-release gain is acquaintance}}, and the inference to \hyperref[entry:claim-some-truths-about-experience-not-physically-deducible]{\textsc{truths about experience that physics cannot deduce}} fails: no truth was ever beyond her.

\webheading{Compared with its neighbours}

Against the ability theorist, the argument does not require that Mary gain any reusable skill — premise 1 keeps acquaintance distinct from know-how. Against the old-fact/new-guise theorist (see \hyperref[entry:argument-old-fact-new-guise]{\textsc{old fact, new guise}}) it does not even require that Mary acquire a new \emph{concept}: the change is in her relation to the experience, not in her conceptual repertoire.

\webheading{Weakest premise}

Premise 3. The proponent of the knowledge argument replies that acquaintance with the experience of red immediately puts Mary in a position to know truths she could not know before (``seeing red is like \emph{this}''). If acquaintance is knowledge-yielding, it is not factually innocent, and the missing truths reappear; if it yields no truths, calling it ``knowledge'' looks honorific, and the hypothesis risks relabelling the datum rather than explaining it.

\begin{plainterms}
The Mary story (see \hyperref[entry:concept-knowledge-argument]{\textsc{the knowledge argument}}) says: she knew all the facts, yet gained something on seeing red — so some facts aren't physical. This reply: what she gained wasn't a fact at all but an \emph{encounter}. Knowing every fact about a city isn't the same as standing in it; arriving doesn't add a line to the guidebook. Mary finally stands in front of red. The worry about this reply: the moment she's there, can't she now state truths she couldn't before (``so \emph{this} is what red is like'')? If yes, the missing facts sneak back in.
\end{plainterms}

\webheading{See also}
\begin{seealso}
  \item \hyperref[entry:argument-knowledge-argument]{\textsc{Knowledge argument (Mary's room)}} — the argument this reply rebuts, by reclassifying Mary's gain as a relation rather than a fact.
  \item \hyperref[entry:argument-ability-hypothesis]{\textsc{Ability hypothesis}} — the sibling reply that reads the same verdict as a gain of skill, not acquaintance.
  \item \hyperref[entry:argument-old-fact-new-guise]{\textsc{Old fact, new guise}} — the reply that grants a new \emph{concept} of an old fact; acquaintance needs not even that.
\end{seealso}

\end{multicols}
\bookgroup{B}
\begin{multicols}{2}
\entryhead{argument-bike-rider-learns-facts}{Bike-rider rejoinder}{Argument \,\textperiodcentered\, after Stanley}{Bike-rider rejoinder}{The bike rider learns facts — the ability hypothesis rests on a false dichotomy}{$\vdash$}
This reply targets the ability hypothesis (see \hyperref[entry:argument-ability-hypothesis]{\textsc{Ability hypothesis}}), the popular escape from the knowledge argument which holds that when Mary first sees red she gains only \emph{abilities} — recognising, imagining, remembering the colour — and not new factual knowledge. The objection grants that what Mary gains might be exactly the bike-rider's kind of learning, then denies the unspoken assumption the escape needs: that acquiring an ability is not itself the learning of a fact.

\webheading{The argument}

The inference is \emph{deductive}, drawn from an independently defended thesis in epistemology.

\begin{enumerate}
  \item Know-how is a species of knowledge-\emph{that} (see \hyperref[entry:claim-knowhow-is-knowledge-that]{\textsc{knowing-how is a kind of knowing-that}}): to know how to F is to know, of a way \emph{w}, that \emph{w} is a way for one to F, grasped under a \emph{practical mode of presentation} — a hands-on way of holding a truth that one cannot fully put into words (see \hyperref[source:source-stanley-williamson-2001]{\textsc{Stanley \& Williamson — Knowing How (2001)}}). Learning to ride a bike is not fact-free skill acquisition; one learns truths about how one's own body and the machine work together.
  \item The ability hypothesis blocks the knowledge argument only if Mary's gained abilities involve \emph{no} new knowledge-that — ``abilities are had, not true.''
  \item By (1), the abilities Mary gains are themselves constituted by propositional knowledge, held under practical or recognitional modes of presentation.
\end{enumerate}

It follows that the ability hypothesis fails \emph{on its own terms}. Even granting that Mary's gain is exactly the bike-rider's learning, that learning is fact-learning, so ``she gains abilities'' does not entail ``she gains no new fact,'' and the hypothesis's conclusion (see \hyperref[entry:claim-mary-gains-abilities-not-facts]{\textsc{that Mary gains abilities rather than facts}}) is left unsupported.

\webheading{What it does and does not establish}

Closing the ability-hypothesis route does \textbf{not} rescue the knowledge argument (see \hyperref[entry:argument-knowledge-argument]{\textsc{Knowledge argument (Mary's room)}}). The facts the bike-rider learns are facts about her own body-and-situation, grasped practically — which is grist for the indexical reply (see \hyperref[entry:argument-indexical-reply]{\textsc{Indexical reply}}): reclassifying Mary's gain as factual but \emph{self-locating} leaves physicalism untouched while abandoning the fact/ability dichotomy of Lewis's original proposal (see \hyperref[source:source-lewis-1988]{\textsc{Lewis — What Experience Teaches (1988)}}). The effect is to move the debate from ``fact versus ability'' to ``what kind of fact, under what mode of presentation.''

\webheading{The weakest premise}

Premise (1) is the contested one. The \emph{Rylean regress} presses that \emph{applying} knowledge-that itself requires know-how, which loops the analysis; dissociation cases, such as motor skills surviving amnesia, suggest two distinct cognitive systems rather than one piece of knowledge held under two modes; and ``practical modes of presentation'' are accused of doing the same unexplained work that \emph{phenomenal concepts} do in the old-fact/new-guise reply (see \hyperref[entry:argument-old-fact-new-guise]{\textsc{Old fact, new guise}}) — a label for the difference rather than an account of it. That worry takes its sharpest shape in the isolation/explanation trade-off facing such concepts (see \textsc{Phenomenal-concept trade-off}).

\begin{plainterms}
A famous escape from the Mary puzzle (see \hyperref[entry:argument-ability-hypothesis]{\textsc{Ability hypothesis}}) says that when Mary first sees red she only gains \emph{skills} — like learning to ride a bike — and skills are not facts, so no fact was missing from her science. This objection takes aim at the ``skills aren't facts'' step. When you learn to ride a bike you plainly do come to know something true: how your body and that bike work together, known in a hands-on way you cannot fully put into words. If gaining a skill is itself a way of learning facts, the escape route closes: calling Mary's gain a ``skill'' no longer shows she learned no fact. It only reopens the original question in a sharper form — what \emph{kind} of fact? A general one physics should have covered, or a personal, you-are-here fact no textbook could ever contain? The ``she learns where she is, not what the world contains'' reply (see \hyperref[entry:argument-indexical-reply]{\textsc{Indexical reply}}) picks up exactly there.
\end{plainterms}

\webheading{See also}
\begin{seealso}
  \item \hyperref[entry:argument-ability-hypothesis]{\textsc{Ability hypothesis}} — the target: the view that Mary gains abilities, not facts, which this reply undercuts on its own terms.
  \item \hyperref[entry:claim-knowhow-is-knowledge-that]{\textsc{Know-how as knowledge-that}} — the load-bearing premise that know-how is a kind of knowledge-that.
  \item \hyperref[entry:concept-knowledge-argument]{\textsc{Knowledge argument}} — the thought experiment in the background; this reply leaves it unrescued.
\end{seealso}

\end{multicols}
\bookgroup{C}
\begin{multicols}{2}
\entryhead{claim-consciousness-categorical-not-dispositional}{Categorical, not dispositional}{Claim}{Categorical, not dispositional}{Phenomenal consciousness is categorical, not merely dispositional}{$\odot$}
Phenomenal consciousness is an \emph{occurrent, categorical} matter — it is actually, presently \emph{there} — rather than a merely \emph{dispositional} affair, a matter of what a system \emph{would} do under inputs it may never actually receive. A disposition is an ``iffy'' property: glass counts as fragile in virtue of how it \emph{would} behave if struck, even if it is never struck. Experience does not look like that. It looks like the paradigm of a categorical fact, something occurrently happening, not a standing counterfactual profile.

The standard support for the contrast is that dispositions require a \emph{categorical basis} — fragility holds in virtue of some actual microstructure of the glass — a demand associated with D. M. Armstrong (see \hyperref[source:source-armstrong-1968]{\textsc{Armstrong — A Materialist Theory of the Mind (1968)}}); on the Russellian side, the intrinsic, categorical nature underlying structure is what Galen Strawson (see \hyperref[source:source-strawson-2006]{\textsc{Strawson — Realistic Monism: Why Physicalism Entails Panpsychism (2006)}}) presses as the home of experience. The claim itself stays modest about \emph{what} the categorical basis is: it does not say whether that basis is the system's current physical configuration or some intrinsic nature beneath it. It asserts only that the phenomenal is categorical, not dispositional.

A dispositionalist or functionalist about consciousness would deny it, holding that to be in a phenomenal state just \emph{is} to occupy a causal-dispositional role. The claim is kin to, but distinct from, \textsc{phenomenal determinacy}: determinacy says there is a \emph{definite} fact about one's current phenomenal state, whereas this says the fact is \emph{categorical/occurrent} rather than dispositional. The two function as companion premises behind \textsc{the decomposition-indeterminacy objection} and \textsc{the structure-and-dynamics argument} respectively.

\begin{plainterms}
A ``disposition'' is an iffy, would-do fact: saying glass is fragile means it \emph{would} break if struck, even if it never actually is. A ``categorical'' fact is how something actually, presently \emph{is}. This claim says your conscious experience is the second kind — it's really there right now, not just a bundle of ``what the system would do if poked.'' A pain isn't a standing tendency to react; it's an occurrent felt thing happening. (Anyone who defines a mental state as nothing but a causal role would push back.)
\end{plainterms}

\webheading{See also}
\begin{seealso}
  \item \textsc{phenomenal determinacy} — the companion claim that there is a \emph{definite} fact about one's current experience; this one adds that the fact is \emph{categorical}.
  \item \textsc{decomposition-indeterminacy objection} — relies on determinacy; pairs with this claim's categoricity in the case against role-based consciousness.
  \item \textsc{structure-and-dynamics argument} — the argument this claim feeds: function captures only structure and dynamics, never the categorical/intrinsic.
\end{seealso}

\entryhead{argument-causal-argument-for-physicalism}{Causal argument for physicalism}{Argument \,\textperiodcentered\, after Papineau}{Causal argument for physicalism}{The causal argument for physicalism}{$\vdash$}
The \textbf{causal argument} (also the \emph{overdetermination} or \emph{exclusion} argument) is the principal positive case for physicalism about the mind. It reasons from the causal structure of the physical world to the conclusion that mental events must be physical events, leaving the anti-physicalist a stark and unwelcome choice.

\textbf{Inference form:} deductive (an argument by elimination).

\begin{enumerate}
  \item \textbf{Completeness.} Every physical effect has a sufficient physical cause (\textsc{Completeness premise} — the \hyperref[entry:concept-completeness-of-physics]{\textsc{Completeness of physics}}).
  \item \textbf{Efficacy.} Some mental events cause physical effects (see \textsc{Mental causation}).
  \item \textbf{No overdetermination.} Those physical effects are not systematically overdetermined by a distinct physical cause and a distinct non-physical cause (see \textsc{No systematic overdetermination}).
  \item From (1) a bodily movement caused by a mental event already has a sufficient physical cause; from (3) the mental cause is not a \emph{second}, independent cause standing beside it. The only way for the mental event to be a genuine cause of that movement, rather than an idle accompaniment, is for it to \emph{be} the sufficient physical cause — i.e. to be physical.
  \item ∴ Mental events are physical events (see \textsc{Mental events are physical}, \hyperref[entry:position-physicalism]{\textsc{Physicalism}}).
\end{enumerate}

The force of the argument is best felt as a \textbf{trilemma}: holding completeness fixed, one must give up exactly one of the remaining options. Deny efficacy and you get \textsc{Epiphenomenalism} — the mind as causally inert by-product. Deny no-overdetermination and you posit a baroque, permanent coincidence in which every action is doubly caused. Accept all three and you get physicalism. Interactionist \textsc{Dualism} is forced to reject premise (1), which is why so much weight rests on whether the completeness of physics is genuinely established.

\textbf{Weakest premise:} the \textbf{no-overdetermination} premise (3). Completeness is defended as hard empirical fact, and efficacy is close to a datum of common sense, so the dualist's most promising move is to accept that mental and physical causes coincide — either by biting the overdetermination bullet (the two causes are distinct but the effect is harmlessly double-caused) or by arguing that mind and brain are not ``distinct'' enough for their joint causation to count as the objectionable kind of overdetermination. The exclusion literature (Kim, \hyperref[source:source-kim-2005]{\textsc{Kim — Physicalism, or Something Near Enough (2005)}}) is largely a fight over premise (3) and what ``distinct sufficient cause'' must mean.

\begin{plainterms}
This is the main argument \emph{for} the view that the mind is physical. It runs in three steps. First: every physical event has a complete physical cause — science never needs to invoke a soul to explain why a body moves (see \hyperref[entry:concept-completeness-of-physics]{\textsc{Completeness of physics}}). Second: your mind really does move your body — your decision to raise your arm is why it goes up. Third: your arm's rising isn't caused \emph{twice over}, once by physics and again, separately, by a non-physical mind. Put those together and there's only one place left for your decision to fit: it has to \emph{be} the physical cause — that is, it has to be something physical happening in your brain. The alternatives are both uncomfortable: either your decisions don't actually do anything (they just ride along — that's \textsc{Epiphenomenalism}), or every single action you take is mysteriously caused two independent times at once. Faced with that menu, the physicalist says: simplest to admit the mind is physical.
\end{plainterms}

\entryhead{concept-causal-closure}{Causal closure of the physical}{Concept}{Causal closure of the physical}{}{$\blacklozenge$}
The \textbf{causal closure of the physical} is the principle that the physical domain is causally self-enclosed: no non-physical event ever causes a physical one. Causal chains that arrive at a physical effect never originate, or pass, outside physics. Put as a slogan, \emph{physics is a closed system} — to trace the cause of any physical happening is to stay forever among physical happenings.

Closure is the negative, boundary-policing face of a thesis whose positive face is the \hyperref[entry:concept-completeness-of-physics]{\textsc{Completeness of physics}}: completeness says each physical effect \emph{has} a sufficient physical cause; closure says it has \emph{no} non-physical cause. The two are easy to conflate and are often used interchangeably, but they come apart on one possibility — \textbf{overdetermination}. Completeness alone is silent about whether a physical effect might \emph{also} have a distinct non-physical cause alongside its sufficient physical one (as a death may have two independently sufficient causes). It is the further premise that such systematic double-causation does not occur (see \textsc{No systematic overdetermination}) that turns completeness into full closure. Hence the dependency: closure is what you get from completeness once overdetermination is ruled out.

\webheading{Why it matters}

Closure is the premise that forces the central dilemma of philosophy of mind. Suppose mental events make a real difference to the physical world — that a decision moves a hand, a pain prompts a wince. If the physical domain is closed, then every such bodily effect already has its cause \emph{inside} physics, and the mental cause can earn its keep in only one way: by \emph{being} physical. The alternatives are stark — either the mental is identical to (or realised by) the physical (see \hyperref[entry:position-physicalism]{\textsc{Physicalism}}, \textsc{Mental events are physical}), or it is a causal idler that the physical world runs on without (see \textsc{Epiphenomenalism}). This is the backbone of Kim's \textbf{causal exclusion argument} (see \hyperref[source:source-kim-2005]{\textsc{Kim — Physicalism, or Something Near Enough (2005)}}): given closure, supervenience, and an exclusion principle barring redundant double-causes, an irreducible mental property would have no work to do. Interactionist \textsc{Dualism}, which insists the non-physical mind pushes the body around, is precisely the denial of closure.

It is worth marking what closure does \emph{not} say. It does not deny the existence of non-physical things, nor even that non-physical things have effects — only that they have no effects \emph{on the physical}. A closure-respecting dualist could still hold that the physical causes the mental (the body produces experiences) while denying any traffic the other way; that is one route to \textsc{Epiphenomenalism}. What closure forbids is causal influence crossing \emph{into} the physical from outside it.

\begin{plainterms}
Causal closure says the physical world is a sealed room as far as causes go: nothing from outside — no soul, no free-floating mind, no spirit — ever reaches in and nudges a physical thing. If you follow the chain of causes behind any physical event, you only ever meet more physical events; you never cross a border into the non-physical. This is the step that makes the mind-body problem bite. If your decisions really do move your body, and nothing non-physical is allowed to move physical stuff, then your decisions must themselves be physical goings-on. The only escape is to say your decisions don't actually \emph{do} anything — they just ride along — which is a hard pill (that is what ``epiphenomenalism'' amounts to). It is the close cousin of the \hyperref[entry:concept-completeness-of-physics]{\textsc{Completeness of physics}}: completeness says a physical cause is always \emph{enough}, closure adds that a non-physical cause is never \emph{there}.
\end{plainterms}

\webheading{See also}
\begin{seealso}
  \item \hyperref[entry:concept-completeness-of-physics]{\textsc{Completeness of physics}} — the positive companion principle; closure is completeness plus the exclusion of overdetermination.
  \item \textsc{No systematic overdetermination} — the premise that bridges completeness to closure.
  \item \hyperref[entry:argument-causal-argument-for-physicalism]{\textsc{Causal argument for physicalism}} — the inference closure powers, against a freely-causing non-physical mind.
  \item \hyperref[entry:position-physicalism]{\textsc{Physicalism}} — the view closure is wielded to support.
  \item \textsc{Dualism} — interactionist dualism is the denial of closure.
  \item \textsc{Epiphenomenalism} — the position left to a non-physical mind once closure is granted.
  \item \hyperref[source:source-kim-2005]{\textsc{Kim — Physicalism, or Something Near Enough (2005)}} — Kim's exclusion argument, the sharpest deployment of closure.
\end{seealso}

\entryhead{argument-causal-exclusion}{Causal exclusion argument}{Argument \,\textperiodcentered\, after Kim}{Causal exclusion argument}{Kim's causal exclusion (supervenience) argument}{$\vdash$}
The \textbf{causal exclusion argument} (Kim's \emph{supervenience argument}; \hyperref[source:source-kim-2005]{\textsc{Kim — Physicalism, or Something Near Enough (2005)}}) is an internal critique of \textbf{non-reductive physicalism} — the popular view that mental properties supervene on physical properties without being identical to them. Kim argues that this position cannot, in the end, secure mental causation: an irreducible mental property is squeezed out of the causal story by the very physical base on which it supervenes. The non-reductive physicalist is driven to a dilemma: reduce the mental, or give up its efficacy.

\textbf{Inference form:} deductive (a reductio of irreducible mental causation).

\begin{enumerate}
  \item \textbf{Closure.} Every physical effect has a sufficient physical cause (\textsc{Completeness premise} — the \hyperref[entry:concept-completeness-of-physics]{\textsc{Completeness of physics}}).
  \item \textbf{Supervenience without identity.} The non-reductive physicalist holds that a mental property \emph{M} supervenes on a physical base \emph{P} but is not identical to \emph{P}. (This is the position's own defining commitment, assumed here for reductio.)
  \item \textbf{Exclusion.} A single effect does not have two distinct, independently sufficient causes, barring genuine overdetermination (see \textsc{No systematic overdetermination}).
  \item Suppose \emph{M} causes a physical effect \emph{E}. By (1), \emph{E} already has a sufficient physical cause — paradigmatically \emph{P} itself, the supervenience base of \emph{M}. By (2), \emph{M} is distinct from \emph{P}. By (3), \emph{M} and \emph{P} cannot both be independently sufficient causes of \emph{E} (and the supervenience relation makes a coincidence of two \emph{independent} causes the wrong description). So \emph{M} is excluded: the causal work is done by \emph{P}.
  \item ∴ An irreducible mental property is causally excluded by its physical base (see \hyperref[entry:claim-irreducible-mental-causally-excluded]{\textsc{Exclusion of the irreducible mental}}).
\end{enumerate}

The result is the \textbf{reduce-or-eliminate dilemma}. If mental causation is to be preserved, \emph{M} must after all \emph{be} a physical property (see \textsc{Mental events are physical}) — reductive physicalism, the conclusion Kim ultimately embraces (``physicalism, or something near enough''). Refuse the identification and \emph{M} is left inert, an \textsc{epiphenomenon} (see \textsc{Causally inert qualia}). Either way, non-reductive physicalism as a stable middle position collapses.

\textbf{Distinct from} the \hyperref[entry:argument-causal-argument-for-physicalism]{\textsc{causal argument for physicalism}}: that argument runs the same exclusion machinery \emph{outward}, against the interactionist \textsc{dualist} whose mind is non-physical, to conclude that the mind is physical. Kim's exclusion argument runs \emph{inward}, against the fellow physicalist who thinks supervenience buys an autonomous mental level; its distinctive premise is supervenience-without-identity (2), and its distinctive target is the \emph{vertical} relation between a property and its own realiser, not a horizontal contest between physical and non-physical causes.

\textbf{Weakest premise:} the \textbf{exclusion} principle (3), as applied to a property and its supervenience base. Non-reductive physicalists reply that \emph{M} and \emph{P} are not ``distinct'' in the way exclusion requires — they are not two competing causes but one event described at two levels, so their joint relevance to \emph{E} is not the objectionable overdetermination the principle rules out. Others deny that closure plus exclusion settles which level is the \emph{real} cause, appealing to a more permissive (e.g. counterfactual or proportionality-based) account of causation on which \emph{M} can be the difference-maker even when \emph{P} suffices. The dispute thus turns on what ``distinct sufficient cause'' must mean.

\begin{plainterms}
Many philosophers want to say the mind is made of nothing but physical stuff, yet is not \emph{just} the brain — mental properties ``ride on'' brain properties without being identical to them. Kim's argument says this comfortable middle ground doesn't hold. Whenever a mental state seems to cause something — your decision makes your arm rise — the underlying brain state is already enough to make the arm rise. So either the decision \emph{is} that brain state (and we're back to plain ``the mind is the brain''), or the decision is a spare wheel that spins without driving anything (epiphenomenalism). There's no third option where the mind is both distinct from the brain and still does real work.
\end{plainterms}

\webheading{See also}
\begin{seealso}
  \item \hyperref[entry:claim-irreducible-mental-causally-excluded]{\textsc{Exclusion of the irreducible mental}} — the conclusion: a supervening-but-distinct mental property has no causal work left to do.
  \item \hyperref[entry:argument-causal-argument-for-physicalism]{\textsc{Causal argument for physicalism}} — the same exclusion machinery aimed outward at dualism rather than inward at non-reductive physicalism.
  \item \textsc{Epiphenomenalism} and \textsc{Causally inert qualia} — the horn of the dilemma that keeps the mental distinct at the cost of its efficacy.
  \item \textsc{Mental events are physical} — the reductive horn Kim recommends, the only way to keep the mental a cause.
  \item \hyperref[entry:concept-causal-closure]{\textsc{Causal closure of the physical}} and \hyperref[entry:concept-completeness-of-physics]{\textsc{Completeness of physics}} — the closure premise the argument leans on.
\end{seealso}

\entryhead{concept-chinese-nation}{Chinese nation}{Concept \,\textperiodcentered\, after Block}{Chinese nation}{Chinese Nation (homunculi-headed robot)}{$\blacklozenge$}
The \emph{Chinese Nation} is a thought experiment devised by Ned Block in ``Troubles with Functionalism'' (1978). Imagine the citizens of a whole country, or a head full of little people, given radios and a rulebook so that together they realise the entire functional and causal organisation of a single human mind: the same state-to-state transition structure a brain runs through, for an hour or so. By construction the system has exactly the inner organisation a conscious brain has. The question Block presses is whether it would therefore be conscious. Would there be something it is like to be the nation, or is there simply nobody home?

\webheading{What makes it distinctive}

The case has its bite precisely because it instantiates the functional organisation, not merely the profile of inputs and outputs a mind produces. It is therefore not the ``right behaviour, no processing'' case of Block's other machine, \textsc{the lookup-table mind}, but the ``right processing, yet intuitively nobody home'' case. That is what makes it tell against \textsc{Functionalism} proper rather than against mere behaviourism: it grants everything a functionalist asks for and still seems to lack experience.

\webheading{The dispute it feeds}

The Chinese Nation is pressed on whether realising a mind's functional organisation is \emph{sufficient for \hyperref[entry:concept-phenomenal-consciousness]{\textsc{phenomenal consciousness}}}. The intuition that the nation has none drives \textsc{the claim that such a population would feel nothing} and the argument \textsc{that organisation is not sufficient for consciousness}, which in turn support the \hyperref[entry:claim-phenomenal-residue]{\textsc{Phenomenal residue}} view and attack the uniform reading of \textsc{substrate independence}.

\webheading{Replies}

Two replies are characteristic, and both differ from the ones the lookup table invites. The \emph{systems reply} denies the intuition outright: the nation taken as a whole may be a subject of experience even though no individual citizen is (Lycan, Dennett). A second worry concerns \emph{speed and timing}, whether so slow and distributed a system genuinely realises the organisation at all. There is also a pointed structural observation. \textsc{Organisational-role functionalism} cannot cheaply exclude the Chinese Nation the way it excludes the lookup table, because here the organisation really is present; that is exactly why the case is sharpened into \textsc{the dense China-brain regress} against the organisational grain.

\webheading{A literary illustration}

A vivid cousin of the case appears in Liu Cixin's science-fiction novel \emph{The Three-Body Problem} (see \hyperref[source:source-liu-cixin]{\textsc{Liu Cixin}}). There the First Emperor Qin Shi Huang, prompted by von Neumann, ranks thirty million soldiers across a plain, each small unit raising black-and-white flags while cavalry gallop the signals between formations, to assemble a vast human machine. In the novel the machine is a computer crunching the orbits of three suns; re-task it to run a \emph{brain} instead. Let each cluster of soldiers be a neuron, each raised flag its firing, the messengers its axons, and the whole drilled host step for one hour through exactly the state-to-state transitions a human cortex would. The army is now no calculator but the Chinese Nation itself: a country physically realising a mind's causal organisation while no single soldier feels a thing. If the soldiers enact the right transitions, is someone \emph{over and above} them undergoing the experience, or is there nobody home?

\begin{plainterms}
Imagine the entire population of a country handed radios and a rulebook and told to each act as one brain cell: when you receive these signals, send those. Follow the rules fast enough and the nation as a whole mimics, step for step, everything one human brain does inside, not just the outward behaviour but the actual inner wiring. Now the unsettling question: the country is doing literally everything your brain does when you feel pain, so is the \emph{country} feeling pain? Is there one extra ``someone'', the nation itself, having an experience, hovering above millions of people who each feel nothing out of the ordinary? If running the right pattern is all it takes to be conscious, you seem pushed to say yes, yet most people's gut says surely not. That clash is the whole point of the thought experiment. It is a sharper challenge than a dumb answer-lookup machine, because here the full inner organisation really is present, so you cannot dismiss it for having no inner structure.
\end{plainterms}

\webheading{See also}
\begin{seealso}
  \item \textsc{Lookup-table mind} — Block's other machine: right behaviour with no inner processing, the case this one is defined against.
  \item \textsc{Functionalism} — the theory the thought experiment targets: mind as functional organisation.
  \item \textsc{Chinese nation (absent qualia)} — the argument that the realised organisation still lacks qualia, so organisation is not sufficient for consciousness.
  \item \hyperref[entry:claim-phenomenal-residue]{\textsc{Phenomenal residue}} — the view the absent-qualia intuition supports: experience outruns functional role.
  \item \textsc{Substrate-independence thesis} — the thesis whose uniform reading the case attacks.
  \item \textsc{Dense China-brain regress} — where the case is sharpened against organisational-role functionalism.
  \item \hyperref[entry:concept-phenomenal-consciousness]{\textsc{Phenomenal consciousness}} — what the nation is alleged to lack despite the right organisation.
  \item \hyperref[source:source-block-1978]{\textsc{Block — Troubles with Functionalism (1978)}} — Ned Block's ``Troubles with Functionalism'' (1978), where the thought experiment is introduced.
\end{seealso}

\entryhead{concept-chinese-room}{Chinese room}{Concept \,\textperiodcentered\, after Searle}{Chinese room}{The Chinese Room}{$\blacklozenge$}
The \emph{Chinese Room} is John Searle's thought experiment against the claim that running the right computer program is, by itself, enough to produce a mind (see \hyperref[source:source-searle-1980]{\textsc{Searle — Minds, Brains, and Programs (1980)}}). A monolingual English speaker is sealed in a room with a rulebook — in effect, a program — that tells him how to manipulate Chinese symbols purely by their shapes. Chinese questions are slipped in; by following the rules he passes back Chinese answers indistinguishable from a native speaker's. Yet he understands no Chinese. The room, Searle argues, has the right \emph{syntax} (formal symbol manipulation) without any \emph{semantics} (meaning, understanding, intentionality) — and a digital computer, which does nothing but manipulate symbols by their shapes, is in exactly the same position.

\webheading{What it targets}

The argument is aimed at \textbf{strong AI}, the computational-functionalist thesis that \emph{implementing the right program is sufficient for a mind}. That precise target is what separates it from two neighbouring thought experiments with which it is often confused. Ned Block's \textsc{Lookup-table mind} (``Blockhead'') attacks the sufficiency of \emph{behaviour}: it produces the right outputs with \textbf{no} internal processing at all, just a giant table of canned answers. The Chinese Room, by contrast, \emph{does} run a genuine program; what it denies is that \emph{computation} suffices for \emph{understanding}. Block's \hyperref[entry:concept-chinese-nation]{\textsc{Chinese nation}} differs in the other direction: there a population realises a mind's full functional organisation, and the case is pressed against \emph{qualia} — felt experience — whereas the Room is pressed against \emph{intentionality}, whether the symbols mean anything to the system at all. Three different experiments, three different sufficiency theses.

\webheading{Characteristic replies}

Each of the standard replies is an objection-in-waiting. The \textbf{Systems Reply} grants that the man does not understand Chinese but locates the understanding in the whole system — man, rulebook, and room together — much as no single neuron understands English. The \textbf{Robot Reply} embeds the program in a robot body with cameras and limbs, so that its symbols stand in causal commerce with the world and thereby acquire grounding. The \textbf{Brain Simulator Reply} has the program simulate, neuron by neuron, the actual brain activity of a Chinese speaker — what relevant difference could remain? Searle rejoins each: against the Systems Reply, for instance, he lets the man memorise the rulebook and work in the open air, so that the man \emph{is} the whole system — and still understands nothing.

\begin{plainterms}
Imagine you're locked in a room with a giant instruction book. Notes in Chinese slide under the door; you can't read a word of Chinese, but the book says ``if you see these squiggles, copy out those squiggles,'' so you shuffle symbols and pass back replies that fool fluent speakers outside. You're behaving exactly like someone who understands Chinese — yet you understand nothing; you're just moving shapes around. Searle's point: a computer is in the same boat. Running a program is \emph{shuffling symbols by their shapes}, and no amount of that, he says, adds up to actually \emph{understanding} what the symbols mean. (The famous comeback: maybe \emph{you} don't understand, but the whole room-plus-rulebook system does.)
\end{plainterms}

\webheading{See also}
\begin{seealso}
  \item \textsc{Functionalism} — the family of theories in the argument's crosshairs: mind as the running of the right program, independent of what runs it.
  \item \textsc{Lookup-table mind} — the neighbouring case with right behaviour and \emph{no} processing; attacks behavioural sufficiency, where the Room attacks computational sufficiency.
  \item \hyperref[entry:concept-chinese-nation]{\textsc{Chinese nation}} — the neighbouring case with the full functional organisation, pressed against qualia where the Room is pressed against understanding.
  \item \hyperref[source:source-searle-1980]{\textsc{Searle — Minds, Brains, and Programs (1980)}} — ``Minds, Brains, and Programs'', the original statement of the argument and its replies.
\end{seealso}

\entryhead{argument-cognitive-closure}{Cognitive closure}{Argument \,\textperiodcentered\, after Mcginn}{Cognitive closure}{McGinn's cognitive-closure argument for mysterianism}{$\vdash$}
\emph{Cognitive closure} is Colin McGinn's argument that the \hyperref[entry:concept-hard-problem]{\textsc{hard problem}} of consciousness has a perfectly natural solution which the human mind is nonetheless built in such a way that it can never grasp. There is some property of the brain in virtue of which it gives rise to consciousness; but the only faculties by which we could form a concept of that property — looking inward, or studying the brain from outside — are each unfit to reach it. So the mind–body problem is, in McGinn's phrase, \emph{soluble in itself but insoluble by us}. The resulting position is \textsc{Mysterianism}.

\webheading{The argument}

The inference is \emph{eliminative} (or \emph{transcendental}): granting that an explanation exists, it rules out every route by which we might come to grasp it.

\begin{enumerate}
  \item There is some natural property — call it \emph{P} — of the brain in virtue of which the brain gives rise to consciousness (see \textsc{a natural explanation exists}).
  \item A mind can be \emph{constitutively closed} to certain properties — barred in principle, not merely currently ignorant, from forming the relevant concepts (see \textsc{minds can be closed to some concepts}).
  \item We could form the concept of \emph{P} in only one of two ways: by \emph{introspection}, or by \emph{outer perception} of the brain and inference from it. (This exhaustive-routes assumption is a premise of the argument's own.)
  \item But introspection presents experience, not its neural basis; and perception of the brain yields only spatial and physical concepts, which cannot capture the subjective — this is the \textsc{explanatory gap}. So neither route can reach \emph{P}.
\end{enumerate}

∴ \emph{P} is cognitively closed to us: the mind–body problem is real and has a natural answer, yet that answer lies permanently beyond human understanding (see \textsc{we are closed to the psychophysical link}). This is the support for \textsc{Mysterianism}.

\webheading{The weak point}

The argument's most disputed step is the pair (3)–(4): the claim that introspection and outer perception are the \emph{only} routes to \emph{P}, and that both are sealed off. Three replies press there. First, science routinely forms \emph{theoretical} concepts that are neither introspected nor perceived — electric charge, the curvature of spacetime — by abduction and outright conceptual innovation, so the two-route dichotomy looks too narrow and McGinn under-rates our capacity to create new concepts. Second, we cannot reliably establish our own representational limits \emph{in advance}, so a confident verdict of closure is itself suspect. Third, if the explanatory gap is \emph{merely} a fact about cognitive closure, that tells against also reading it as a metaphysical residue: mysterianism and the residue view (see \hyperref[entry:claim-phenomenal-residue]{\textsc{that experience outruns functional role}}) cannot both take the gap at face value.

\webheading{How it differs from neighbouring gap arguments}

All of these start from the same explanatory gap but draw opposite morals from it. The \textsc{explanatory-gap argument} treats the gap as evidence about \emph{metaphysics} — that the phenomenal is not exhausted by functional role. Cognitive closure treats the very same gap as evidence about \emph{our cognitive limits}: the problem is real and natural, merely humanly unsolvable. It differs again from the \textsc{decomposition-indeterminacy objection} and the \textsc{structure-and-dynamics argument}, which contend that function \emph{omits} or \emph{fails to fix} the phenomenal; mysterianism instead grants that a natural answer exists and denies only our access to it.

\begin{plainterms}
The argument goes: (1) there's some real, physical reason the brain produces consciousness; (2) minds can be permanently unable to understand certain things; (3) the only ways we could ever get hold of that reason are by looking inward (introspection) or by studying the brain from outside; (4) but looking inward shows you the feeling, not the wiring, and studying the brain from outside only ever gives you physical/spatial facts that never add up to a feeling. So neither route can deliver the answer — meaning it's locked away from us for good. The main pushback: science invents concepts all the time that we neither feel nor see directly (like ``electron''), so maybe the two routes aren't the only ones.
\end{plainterms}

\webheading{See also}
\begin{seealso}
  \item \textsc{Mysterianism} — the view this argument supports: the mind–body problem has a natural solution we are constitutionally unable to reach.
  \item \hyperref[entry:concept-hard-problem]{\textsc{Hard problem}} — the problem McGinn declares humanly insoluble.
  \item \textsc{Explanatory-gap argument} — the same gap read as metaphysical evidence rather than as a limit on our concepts.
  \item \hyperref[entry:claim-phenomenal-residue]{\textsc{Phenomenal residue}} — the residue reading the third reply says cannot coexist with a merely-cognitive diagnosis of the gap.
  \item \textsc{Decomposition-indeterminacy objection} and \textsc{Structure-and-dynamics argument} — neighbouring gap arguments that fault function itself, where mysterianism faults only our access.
  \item \textsc{Natural explanation of consciousness}, \textsc{Possibility of cognitive closure}, \textsc{Humans closed to the link} — the argument's two premises and its conclusion.
\end{seealso}

\entryhead{claim-phenomenal-intensions-coincide}{Coincident phenomenal intensions}{Claim \,\textperiodcentered\, after Chalmers}{Coincident phenomenal intensions}{For phenomenal concepts, the two intensions coincide}{$\odot$}
For a phenomenal concept — the concept of what it is like to undergo an experience, such as the concept of \emph{the feeling of pain} or \emph{phenomenal red} — the \hyperref[entry:concept-primary-intension]{\textsc{Primary intension}} and the \hyperref[entry:concept-secondary-intension]{\textsc{Secondary intension}} are one and the same function. Unlike ``water,'' whose reference is fixed by a contingent role (the watery stuff) that something other than H₂O could have filled, a phenomenal concept fixes its reference directly on the felt quality itself, with no reference-fixing description standing between the concept and its object. There is therefore no wedge between how the experience \emph{presents} and what it \emph{is}: for a feeling, the appearance is the reality.

The contrast is the whole point. With natural-kind terms the two intensions diverge, and that gap is where a Kripkean necessity can be both necessary and knowable only a posteriori (see \textsc{A-posteriori necessities}); the two-dimensionalist exploits the gap to recast apparent impossibilities as \hyperref[entry:claim-aposteriori-necessities-are-misdescribed-possibilities]{\textsc{A-posteriori necessities as misdescribed possibilities}}. A phenomenal concept offers no such gap. ``Watery stuff that is not H₂O'' describes a genuine possibility misnamed ``water,'' but ``something that feels exactly like pain yet is not pain'' describes nothing at all, because feeling-like-pain is not a fallible mode of presentation of pain — it simply is the pain (see \hyperref[entry:concept-phenomenal-consciousness]{\textsc{Phenomenal consciousness}}).

\webheading{Scope and bearing}

The claim is a thesis in the semantics of phenomenal concepts, not directly a claim about what consciousness is. Its work is to remove an escape route. Chalmers's \hyperref[entry:argument-modal-rationalism]{\textsc{Modal rationalism}} licenses the step from ideal conceivability to \emph{primary} possibility; the question is whether primary possibility carries through to full metaphysical possibility in the phenomenal case. Where primary and secondary intensions coincide, there is no separate secondary necessity in which a metaphysical impossibility could hide — so primary possibility \emph{is} metaphysical possibility. Applied to the \hyperref[entry:concept-philosophical-zombie]{\textsc{Philosophical zombie}}, this closes the gap the type-B physicalist needs: if the zombie is conceivable it is genuinely possible, and physicalism is false (see \hyperref[entry:argument-zombie-conceivability]{\textsc{Zombie argument}}).

\webheading{Who would deny it}

The type-B physicalist (see \hyperref[entry:concept-materialism-types]{\textsc{Type-A / Type-B / Type-C materialism}}) — the position the claim is aimed at. Type-B physicalism holds that ``pain is such-and-such brain state'' is a \emph{strong} a posteriori necessity, with a genuine primary/secondary gap supplied by a special \emph{phenomenal concept} that refers to a physical state under a distinctive mode of presentation. On that view there is, after all, an appearance/reality distinction for experience at the level of \emph{concepts}, the zombie is conceivable but impossible, and the present claim is exactly what fails. The dispute thus narrows to whether phenomenal concepts can carry a reference-fixing mode of presentation distinct from their referent — the hinge on which the entire conceivability argument turns.

\begin{plainterms}
For most words there's a difference between how something seems and what it really is: the stuff in the rivers \emph{seems} like ``water'' but really \emph{is} H₂O, and that gap is why ``water is H₂O'' can be a surprise we had to discover. This claim says feelings are different. With a feeling — pain, or the look of red — there's no gap between the seeming and the being: a state that felt \emph{exactly} like pain in every respect would just \emph{be} pain; there's nothing more to a feeling than how it feels. Why it matters: the physicalist's best reply to the \hyperref[entry:concept-philosophical-zombie]{\textsc{Philosophical zombie}} (a creature physically just like you but with no inner feel) is to say ``sure, you can imagine it, but imagining isn't possibility — just as you can 'imagine' water that isn't H₂O.'' This claim blocks that reply for consciousness: because feelings have no seeming/being gap, imagining a zombie really does show a zombie is possible — and that would mean consciousness is something over and above the physical.
\end{plainterms}

\webheading{See also}
\begin{seealso}
  \item \hyperref[entry:concept-primary-intension]{\textsc{Primary intension}} — the epistemic, reference-fixing dimension of meaning; for phenomenal concepts it has nothing to fix onto but the feel itself.
  \item \hyperref[entry:concept-secondary-intension]{\textsc{Secondary intension}} — the rigid, actual-world-fixed dimension; the claim is that for feelings it does not diverge from the primary intension.
  \item \hyperref[entry:claim-aposteriori-necessities-are-misdescribed-possibilities]{\textsc{A-posteriori necessities as misdescribed possibilities}} — the water-style rescue the claim says is unavailable here, since there is no primary/secondary gap to misdescribe.
  \item \hyperref[entry:argument-modal-rationalism]{\textsc{Modal rationalism}} — supplies the conceivability-to-primary-possibility step that this claim then extends, in the phenomenal case, all the way to metaphysical possibility.
  \item \hyperref[entry:argument-zombie-conceivability]{\textsc{Zombie argument}} — the argument whose final move depends on this claim: with the intensions coincident, the zombie's conceivability yields its genuine possibility.
  \item \hyperref[entry:concept-materialism-types]{\textsc{Type-A / Type-B / Type-C materialism}} — type-B physicalism, the view that denies this claim by positing a phenomenal-concept mode of presentation distinct from its referent.
  \item \hyperref[entry:concept-phenomenal-consciousness]{\textsc{Phenomenal consciousness}} — what phenomenal concepts are concepts \emph{of}; the ``appearance is reality'' point is a point about felt experience.
\end{seealso}

\entryhead{argument-combination-problem}{Combination problem}{Argument \,\textperiodcentered\, after James}{Combination problem}{The combination problem — micro-subjects cannot compose a macro-subject}{$\vdash$}
The \emph{combination problem} is the standing objection to \hyperref[entry:concept-panpsychism]{\textsc{Panpsychism}}, severe enough that many regard it as the view's make-or-break test. Constitutive panpsychism holds that a macro-subject like you is \emph{constituted by} the experiential properties of your fundamental micro-parts. The objection denies that this constitution is available: a crowd of tiny subjects, each with its own private point of view, cannot fuse into one further subject whose experience contains theirs. The problem originates with William James (see \hyperref[source:source-james-1890]{\textsc{1890}}) and is sharpened into its hardest ``subject-summing'' form by Sam Coleman (see \hyperref[source:source-coleman-2014]{\textsc{2014}}); David Chalmers (see \hyperref[source:source-chalmers-2017]{\textsc{2017}}) anatomises its parts.

\webheading{The argument}

The inference is deductive — an objection, so it \emph{attacks} the constitutive thesis rather than concluding a positive claim.

\begin{enumerate}
  \item (see \hyperref[entry:claim-micro-subjects-do-not-sum]{\textsc{Micro-subjects do not sum}}) Distinct micro-subjects — each with its own private point of view and self-contained phenomenal field — do not, merely by being arranged or related, compose a further single subject whose experience contains theirs.
  \item Constitutive panpsychism requires exactly such composition: your unified experience is supposed to be made up of the experiences of your fundamental parts.
\end{enumerate}

∴ The constitution the view needs is unavailable, so it cannot deliver the single, unified macro-subject we actually find. This \emph{attacks} \hyperref[entry:claim-consciousness-fundamental-ubiquitous]{\textsc{Fundamental ubiquitous consciousness}} at its constitutive joint.

The objection also generalises a deeper worry. Panpsychism was offered to \emph{avoid} getting experience from the non-experiential, yet it now owes an account of getting \emph{macro}-experience from \emph{micro}-experience — so the very explanatory burden the view promised to discharge reappears here, no easier than before. Chalmers distinguishes three faces of the problem: combining micro-\emph{subjects} into a subject, micro-\emph{qualities} into a unified quality-field, and micro-\emph{structure} into the structure of a macro-experience; the subject-combination face is the most resistant.

\webheading{The weak point}

The objection turns on premise (1). Constitutive panpsychists reply with accounts of \emph{phenomenal bonding} or fusion on which subjects do combine after all. Alternatively they retreat to \emph{panqualityism}, which denies that there are micro-\emph{subjects} to combine in the first place — only unexperienced qualities — and pays for the move with a residual explanatory gap at the point where subjecthood has to be added back. Which of these escapes succeeds is the live question on which the fate of constitutive panpsychism turns.

\begin{plainterms}
Panpsychism's appeal is that every tiny speck of matter carries a flicker of experience, so your mind is built up out of those flickers. The combination problem is the bill for that appeal. Your experience right now is \emph{one} unified thing — a single point of view taking in this whole sentence at once. How do countless separate little experiences, each locked inside its own tiny perspective, merge into that one? William James argued they simply cannot: a hundred little feelings do not add up to one big feeling that swallows them — each stays its own private whole. If he is right, the panpsychist's ``build a mind out of mini-minds'' recipe has no step that actually does the building. Worse, it was supposed to spare us the mystery of making feeling out of the feelingless, but now it owes the matching mystery of making \emph{one} feeler out of \emph{many} — which looks just as hard.
\end{plainterms}

\webheading{See also}
\begin{seealso}
  \item \hyperref[entry:claim-micro-subjects-do-not-sum]{\textsc{Micro-subjects do not sum}} — the premise that drives the objection: subjects do not sum.
  \item \hyperref[entry:claim-consciousness-fundamental-ubiquitous]{\textsc{Fundamental ubiquitous consciousness}} — the constitutive thesis this objection attacks.
  \item \hyperref[entry:concept-panpsychism]{\textsc{Panpsychism}} — the view under attack, and where the combination problem sits among its varieties.
  \item \hyperref[entry:position-cosmopsychism]{\textsc{Cosmopsychism}} — the top-down cousin, which trades this problem for its mirror image, the \emph{decombination} problem.
\end{seealso}

\entryhead{concept-completeness-of-physics}{Completeness of physics}{Concept}{Completeness of physics}{}{$\blacklozenge$}
The \textbf{completeness of physics} is the thesis that the physical world is causally self-contained: every physical effect that has a cause at all has a \emph{sufficient} physical cause. Whenever a physical event is produced, one need never step outside physics to find what produced it — a complete physical story can in principle be told for every physical happening. The principle is sometimes called the \emph{causal completeness}, \emph{causal closure}, or \emph{causal self-sufficiency} of the physical domain, and it is the empirical engine of the modern argument for physicalism (see \hyperref[entry:argument-causal-argument-for-physicalism]{\textsc{Causal argument for physicalism}}, \hyperref[entry:position-physicalism]{\textsc{Physicalism}}).

The careful formulation matters. Completeness does not say that \emph{every} event is physical, nor that the physical causes are the \emph{only} causes; it says only that the physical causes already on the books are \textbf{sufficient} — enough, by themselves, to fix the physical effect. So stated, it leaves logical room for a second, non-physical cause of the same effect. Closing that room is the work of a separate premise about overdetermination (see \textsc{No systematic overdetermination}); the two together yield the strong conclusion that nothing non-physical does any \emph{distinctive} causal work in the physical world.

Papineau (see \hyperref[source:source-papineau-2002]{\textsc{Papineau — Thinking about Consciousness (2002)}}) stresses that completeness is not an a priori truth but a hard-won \textbf{empirical discovery}. His ``history of the completeness of physics'' traces how the inventory of basic forces was progressively closed: vital and mental forces, once posited to explain living and intelligent behaviour, were squeezed out as physiology and physics found sufficient physical causes for the bodily movements they were meant to produce. The conservation laws are the usual modern evidence — a non-physical cause injecting energy or momentum into a body would show up as a violation, and none is observed.

\webheading{Quantum indeterminism}

A naive statement of completeness is \emph{deterministic}: every physical effect is fully fixed by prior physical states under physical law. On the standard reading of quantum mechanics that is simply false. When an atom decays or a photon meets a half-silvered mirror, the prior physical state fixes only a \emph{probability distribution} over outcomes, not the outcome itself. There is genuine chance in the basement.

Completeness survives because it was never a claim about \emph{determination} in the first place; it is a claim about \emph{probabilities}. The careful formulation runs: in so far as a physical effect is determined at all, it is determined by physical antecedents, and in so far as it is merely chancy, its objective probabilities are fixed by physical law \emph{alone}. The dial may sometimes land at random, but nothing non-physical sets the odds it lands on. So stated, completeness rules out non-physical \emph{causes}, not physical \emph{chance}: randomness is internal to physics, not a leak in its wall, and the thesis is compatible with genuine indeterminism rather than a disguised commitment to determinism.

This matters for the mind because quantum chance is sometimes floated as a loophole for the \emph{interactionist} — a place where an immaterial will might ``make a difference'' without violating the conservation laws, by nudging which outcome a brain's quantum events take. The probabilistic formulation is exactly what closes the loophole, and it does so by reposing the same dilemma the \hyperref[entry:concept-causal-closure]{\textsc{Causal closure of the physical}} poses elsewhere. Either the mind merely rides along with whatever the physical chances throw up, in which case it does no distinctive work (see \textsc{Epiphenomenalism}); or it \emph{biases} the probabilities, making an outcome likelier than the physical law prescribes, in which case physics is not complete after all. The second horn is a genuine empirical claim, and the physicalist bets it is false: no deviation from quantum statistics has ever been measured in a brain. So indeterminism hands the \textsc{dualist} no free mental causation. It only restates the choice between riding inertly and breaking completeness, with the breach now statistical rather than energetic.

Two cautions. First, the picture is interpretation-dependent: on Bohmian mechanics the world is deterministic again and completeness reverts to its strong form, while on Everettian quantum mechanics there is no single collapse-outcome to be ``caused'' at all. Completeness holds across these readings because in each the \emph{physics} remains causally self-contained, but the gloss on ``randomness'' differs. Second, indeterminism is not by itself a resource for free will: an outcome that is merely random is no more an agent's own doing than one that is determined, so nothing about how completeness treats quantum chance smuggles libertarian agency in through the back door.

\webheading{The completeness premise at work}

Completeness is the load-bearing premise of the \textbf{causal argument for physicalism}: if every physical effect already has a sufficient physical cause, then a mental event that genuinely causes a physical effect (a hand's rising, a word's being spoken) must \emph{be} one of those physical causes, on pain of either doing no work at all (see \textsc{Epiphenomenalism}) or redundantly duplicating a cause already present. This is why anti-physicalists who want the mind to make a difference — interactionist dualists especially — must reject or restrict completeness, and why much of the dispute over \hyperref[entry:concept-hard-problem]{\textsc{Hard problem}} turns on whether the principle is as secure as its defenders claim.

\begin{plainterms}
The idea is that the physical world is a closed shop, causally speaking: whenever something physical happens, there is always a complete physical explanation for it — some earlier physical stuff that was enough to make it happen. You never \emph{have} to bring in a soul, a spirit, or any non-physical push to finish the story. It is offered not as an armchair assumption but as something science has gradually found to be true: every time people expected a special ``life force'' or ``mind force'' to be needed, a purely physical cause turned up instead. If that is right, then a thought that really moves your arm has to be a physical thing too — otherwise the arm would already be fully accounted for and the thought would have nothing left to do.
\end{plainterms}

\webheading{See also}
\begin{seealso}
  \item \hyperref[entry:concept-causal-closure]{\textsc{Causal closure of the physical}} — the closely allied ``no causes cross into the physical from outside'' formulation; completeness plus no-overdetermination yields closure.
  \item \hyperref[entry:argument-causal-argument-for-physicalism]{\textsc{Causal argument for physicalism}} — the inference for which this thesis is the chief premise.
  \item \textsc{Completeness premise} — completeness stated as the truth-apt premise that argument rests on.
  \item \textsc{No systematic overdetermination} — the companion premise that closes the loophole completeness leaves open.
  \item \hyperref[entry:position-physicalism]{\textsc{Physicalism}} — the view this principle is marshalled to support.
  \item \textsc{Epiphenomenalism} — the fate completeness forces on a mental cause that is not itself physical.
\end{seealso}

\entryhead{claim-conceivability-entails-possibility}{Conceivability entails possibility}{Claim \,\textperiodcentered\, contested \,\textperiodcentered\, after Chalmers}{Conceivability entails possibility}{Ideal conceivability entails metaphysical possibility}{$\odot$}
If a scenario is \emph{ideally} conceivable — coherent under idealized rational reflection, with no concealed contradiction — then it is metaphysically possible. This is the modal-epistemology bridge of the zombie argument, and its most contested step. It is emphatically \emph{not} the trivial claim that whatever we can loosely picture is possible; it concerns \emph{ideal} conceivability, what survives unlimited rational scrutiny. On the two-dimensional version (David Chalmers, \hyperref[source:source-chalmers-1996]{\textsc{Chalmers — The Conscious Mind (1996)}}; \hyperref[source:source-chalmers-2002]{\textsc{Chalmers — On Sense and Intension (2002)}}), the principle reads more precisely: ideal \emph{primary} conceivability entails \emph{primary} possibility, and for the zombie case the primary and secondary intensions are taken to coincide.

The bridge is what converts \hyperref[entry:claim-zombie-conceivable]{\textsc{Zombie conceivability}} into a conclusion against physicalism, so the whole dispute concentrates here.

\webheading{Argued both ways}

It is not a brute premise. It is \textbf{defended} by modal rationalism (see \hyperref[entry:argument-modal-rationalism]{\textsc{Modal rationalism}}): once the \hyperref[entry:concept-primary-intension]{\textsc{Primary intension}} and \hyperref[entry:concept-secondary-intension]{\textsc{Secondary intension}} are distinguished, the apparent counterexamples are misdescribed possibilities (see \hyperref[entry:claim-aposteriori-necessities-are-misdescribed-possibilities]{\textsc{A-posteriori necessities as misdescribed possibilities}}), so ideal primary conceivability tracks primary possibility. It is \textbf{attacked} by the a-posteriori-necessity objection (see \hyperref[entry:argument-conceivability-not-guide-to-possibility]{\textsc{A-posteriori-necessity objection}}): Kripkean necessities (see \textsc{A-posteriori necessities}, \hyperref[source:source-kripke-1980]{\textsc{Kripke — Naming and Necessity (1980)}}) already break the link, so a zombie may be conceivable yet impossible — the type-B physicalist reading (see \hyperref[entry:concept-materialism-types]{\textsc{Type-A / Type-B / Type-C materialism}}). A deeper, allied \textbf{phenomenal-concept strategy} (see \textsc{Phenomenal-concept deflation}, \textsc{Phenomenal-concept trade-off}) explains the conceivability as an artefact of how phenomenal concepts work. Throughout, the crux is the same: whether the \emph{phenomenal} case has the appearance/reality gap that lets the water case escape.

The claim is distinct from \hyperref[entry:claim-zombie-conceivable]{\textsc{Zombie conceivability}} (the antecedent it operates on) and from \textsc{Intuition tracks trained domain} (an epistemology of \emph{intuitions about cases}); this one is specifically about the conceivability→possibility inference.

\begin{plainterms}
This is the hinge the whole zombie argument swings on: that \emph{if you can coherently imagine something — really coherently, with no hidden contradiction — then it is genuinely possible}. Grant it, and ``I can imagine a zombie'' turns into ``a zombie is possible,'' which would mean experience is something extra beyond the physical. The classic pushback borrows from science: ``water is H₂O'' is something you \emph{can} imagine being false, yet it is necessarily true — so being able to imagine otherwise does not always prove real possibility. Maybe zombies are like that: imaginable but impossible.
\end{plainterms}

\webheading{See also}
\begin{seealso}
  \item \hyperref[entry:claim-zombie-conceivable]{\textsc{Zombie conceivability}} — the antecedent this bridge operates on.
  \item \hyperref[entry:argument-modal-rationalism]{\textsc{Modal rationalism}} — the defence of the bridge.
  \item \hyperref[entry:argument-conceivability-not-guide-to-possibility]{\textsc{A-posteriori-necessity objection}} — the objection to it.
  \item \hyperref[entry:claim-aposteriori-necessities-are-misdescribed-possibilities]{\textsc{A-posteriori necessities as misdescribed possibilities}} — the premise the defence rests on.
  \item \textsc{A-posteriori necessities} — the Kripkean cases the objection wields.
  \item \hyperref[entry:concept-materialism-types]{\textsc{Type-A / Type-B / Type-C materialism}} — type-B physicalism, the reading on which the bridge fails.
  \item \textsc{Phenomenal-concept deflation} / \textsc{Phenomenal-concept trade-off} — the allied phenomenal-concept strategy.
  \item \textsc{Intuition tracks trained domain} — a distinct point about the warrant of intuitions.
\end{seealso}

\entryhead{argument-no-arbitrary-line-panpsychism}{Continuity argument}{Argument \,\textperiodcentered\, after Nagel}{Continuity argument}{The continuity argument — no non-arbitrary line where experience switches on}{$\vdash$}
The \emph{continuity argument} is a second route to \hyperref[entry:concept-panpsychism]{\textsc{Panpsychism}}, pressed by Thomas Nagel (see \hyperref[source:source-nagel-1979]{\textsc{1979}}). Where the \hyperref[entry:argument-genetic-panpsychism]{\textsc{genetic argument}} reasons from the impossibility of brute emergence, this one reasons from the absence of any principled boundary: somewhere between a rock and a human being, consciousness seems to appear, yet there is no non-arbitrary place to draw the line. Rather than draw one anyway, the panpsychist lets experience reach all the way down.

\webheading{The argument}

The inference is abductive and eliminative — an inference to the most principled of the surviving options, not a deductive proof.

\begin{enumerate}
  \item (see \hyperref[entry:claim-no-nonarbitrary-threshold-for-experience]{\textsc{Absence of an experiential threshold}}) If experience were confined to brains, or to life, or to systems above some complexity threshold, there should be a principled place where it first switches on; but no such non-arbitrary line is in view — the relevant differences across the supposed boundary look like differences of degree and organisation.
  \item Drawing the line anyway, at an arbitrary point, is unmotivated; and denying that experience is real is ruled out (see \hyperref[entry:claim-phenomenal-consciousness-is-real]{\textsc{Reality of phenomenal consciousness}}).
  \item The remaining option that avoids an arbitrary threshold is that experience, in some attenuated form, is present all the way down.
\end{enumerate}

∴ This \emph{supports} \hyperref[entry:claim-consciousness-fundamental-ubiquitous]{\textsc{Fundamental ubiquitous consciousness}}: the line-free, parsimonious position is that experience pervades nature rather than appearing abruptly at one favoured level.

\webheading{The weak point}

The argument leans on premise (1). An opponent can hold that there \emph{is} a principled threshold — a level of integrated information, say, or of functional organisation, at which experience genuinely (if strongly) emerges (see \hyperref[entry:claim-strong-emergence-is-acceptable]{\textsc{Acceptability of strong emergence}}) — in which case the ``no non-arbitrary line'' premise is simply false. The argument is also only abductive: it shows that panpsychism avoids an awkward commitment, not that its rivals are incoherent, so it carries less weight than the genetic argument aims to.

\begin{plainterms}
Somewhere between a rock and a human being, consciousness seems to show up. But where, exactly? Pick any line — ``only brains,'' ``only living things,'' ``only things this complex'' — and it looks arbitrary, because the differences on either side are matters of more-or-less, not a sudden all-or-nothing switch. Rather than draw an unmotivated line in the sand, the panpsychist says: do not draw one at all — let a faint glimmer of experience reach all the way down, growing richer as things get more organised. It is offered as the tidier option, not a knock-down proof; an opponent who can point to a \emph{principled} switch-on point blocks it.
\end{plainterms}

\webheading{See also}
\begin{seealso}
  \item \hyperref[entry:claim-consciousness-fundamental-ubiquitous]{\textsc{Fundamental ubiquitous consciousness}} — the conclusion this argument supports.
  \item \hyperref[entry:claim-no-nonarbitrary-threshold-for-experience]{\textsc{Absence of an experiential threshold}} — its single premise: no principled point at which experience first appears.
  \item \hyperref[entry:argument-genetic-panpsychism]{\textsc{Genetic argument}} — the companion, stronger case for the same conclusion, from the impossibility of brute emergence.
  \item \hyperref[entry:claim-strong-emergence-is-acceptable]{\textsc{Acceptability of strong emergence}} — the denial that powers the main objection: there \emph{is} a principled emergence threshold.
  \item \hyperref[entry:concept-panpsychism]{\textsc{Panpsychism}} — the view this argument supports.
\end{seealso}

\entryhead{position-cosmopsychism}{Cosmopsychism}{Position \,\textperiodcentered\, after Goff}{Cosmopsychism}{}{$\bigstar$}
\emph{Cosmopsychism} holds that the cosmos as a whole is the one fundamental conscious subject, and that individual minds like yours and mine are \emph{derivative} — aspects, or partial dissociations, of that single universal consciousness rather than building blocks of it. It is a priority-monist cousin of panpsychism: where ordinary panpsychism builds large minds \emph{up} from tiny ones, cosmopsychism carves small minds \emph{out} of one great cosmic mind, with grounding running top-down from the whole to its parts.

The position commits to the claim that \hyperref[entry:claim-cosmos-is-fundamental-subject]{\textsc{the universe as a whole is the one fundamental conscious subject}}. It answers the \hyperref[entry:question-hard-problem]{\textsc{hard problem of consciousness}} by placing experience at the fundamental level — here the cosmic level — so that experience never has to be produced from the non-experiential. The view is defended by Goff (see \hyperref[source:source-goff-2017]{\textsc{Goff — Consciousness and Fundamental Reality (2017)}}), Shani (see \hyperref[source:source-shani-2015]{\textsc{Shani — Cosmopsychism: A Holistic Approach to the Metaphysics of Experience (2015)}}), and Nagasawa and Wager (see \hyperref[source:source-nagasawa-wager-2017]{\textsc{Nagasawa \& Wager — Panpsychism and Priority Cosmopsychism (2017)}}).

\webheading{How it differs from neighbouring views}

Like \hyperref[entry:position-panpsychism]{\textsc{Panpsychism (position)}} and \hyperref[entry:position-russellian-monism]{\textsc{Russellian monism}}, cosmopsychism keeps consciousness fundamental rather than derived from non-experiential matter. What sets it apart is the \emph{direction} of grounding. Against constitutive \emph{micro}-panpsychism, the fundamental subjects are not the smallest parts but the whole. This trades one hard problem for its mirror image: in place of the \hyperref[entry:argument-combination-problem]{\textsc{combination problem}} — how do micro-experiences add up to one mind? — cosmopsychism faces the \emph{de}combination problem: how does the one cosmic subject divide into many bounded, walled-off subjects like us?

\begin{plainterms}
Cosmopsychism starts from the top: the basic conscious thing isn't a particle but the \emph{universe as a whole}, and you and I are like ripples or partial ``splits'' within that one great mind. It's panpsychism run in reverse — rather than building big minds out of tiny ones, it carves small minds out of a single cosmic one. That dodges the question ``how do trillions of atom-feelings merge into your one experience?'' but raises the opposite puzzle: how does one universe-wide consciousness break up into separate, walled-off selves?
\end{plainterms}

\webheading{See also}
\begin{seealso}
  \item \hyperref[entry:claim-cosmos-is-fundamental-subject]{\textsc{Cosmos as fundamental subject}} — the core commitment: the cosmos is the one fundamental subject, minds derivative.
  \item \hyperref[entry:question-hard-problem]{\textsc{Hard problem (the question)}} — the problem the position addresses, by placing experience at the cosmic base.
  \item \hyperref[entry:position-panpsychism]{\textsc{Panpsychism (position)}} — the bottom-up cousin it inverts; both keep consciousness fundamental.
  \item \hyperref[entry:concept-panpsychism]{\textsc{Panpsychism}} — the broader doctrine, whose taxonomy sorts cosmopsychism as the whole-first (vs. micropsychist) branch.
  \item \hyperref[entry:position-russellian-monism]{\textsc{Russellian monism}} — the allied view that consciousness is the intrinsic nature physics leaves unspecified.
  \item \hyperref[entry:argument-combination-problem]{\textsc{Combination problem}} — the objection cosmopsychism escapes, only to face its mirror, the decombination problem.
\end{seealso}

\entryhead{claim-cosmos-is-fundamental-subject}{Cosmos as fundamental subject}{Claim \,\textperiodcentered\, after Goff}{Cosmos as fundamental subject}{The universe as a whole is the one fundamental conscious subject}{$\odot$}
The cosmos as a whole is a single fundamental conscious \emph{subject}, and individual minds like ours are \emph{derivative} — aspects or dissociations \emph{of} the one universal consciousness rather than building blocks that compose it. This is \emph{cosmopsychism}, defended by Philip Goff (see \hyperref[source:source-goff-2017]{\textsc{Goff — Consciousness and Fundamental Reality (2017)}}), Itay Shani (see \hyperref[source:source-shani-2015]{\textsc{Shani — Cosmopsychism: A Holistic Approach to the Metaphysics of Experience (2015)}}), and Yujin Nagasawa (see \hyperref[source:source-nagasawa-wager-2017]{\textsc{Nagasawa \& Wager — Panpsychism and Priority Cosmopsychism (2017)}}). It is a priority-monist cousin of \hyperref[entry:concept-panpsychism]{\textsc{Panpsychism}}: grounding runs \emph{top-down}, from the whole to its parts, rather than bottom-up.

The view combines two commitments. The first is \emph{priority monism} — the thesis that the cosmos is the one fundamental concrete object, on which everything else depends. The second is experiential fundamentality — that this one fundamental object is conscious. Its characteristic difficulty is the mirror image of the trouble panpsychism faces. Where micro-level panpsychism wrestles with the \emph{combination problem} (how do many tiny experiences add up to one unified mind?), cosmopsychism faces the \emph{decombination problem}: how does the single universal subject give rise to many distinct, bounded subjects like us?

It is denied by constitutive micro-panpsychists, for whom the fundamental subjects are micro rather than cosmic; by physicalists and dualists; and by anyone who rejects priority monism. It is distinct from \hyperref[entry:claim-consciousness-fundamental-ubiquitous]{\textsc{the constitutive-panpsychist claim}} precisely on the direction of grounding: constitutive panpsychism builds macro-minds \emph{up} from micro-experiences, whereas cosmopsychism derives them \emph{down} from the whole — combination versus decombination.

\begin{plainterms}
Most views that make consciousness fundamental put it in the tiniest bits of matter and try to build big minds out of many little ones. Cosmopsychism flips this: it says the \emph{whole universe} is the one basic conscious thing, and individual minds like yours are something the universal mind splits into, the way a single person with a dissociative condition can seem to host several selves. The hard part isn't getting from small to big — it's the reverse: explaining how one cosmic mind comes to contain many separate, walled-off minds like ours.
\end{plainterms}

\webheading{See also}
\begin{seealso}
  \item \hyperref[entry:claim-consciousness-fundamental-ubiquitous]{\textsc{constitutive-panpsychist claim}} — the bottom-up rival: consciousness in the smallest things, composed up into minds; cosmopsychism runs the other way.
  \item \hyperref[entry:position-cosmopsychism]{\textsc{Cosmopsychism}} — the named position this claim states.
  \item \hyperref[entry:concept-panpsychism]{\textsc{Panpsychism}} — the broader family; cosmopsychism is its priority-monist variant.
\end{seealso}

\end{multicols}
\bookgroup{E}
\begin{multicols}{2}
\entryhead{argument-strong-emergence-acceptable}{Emergence reply}{Argument}{Emergence reply}{The emergence reply — experience can strongly emerge, so no need for fundamental mind}{$\vdash$}
The \emph{emergence reply} is the standard physicalist answer to the \hyperref[entry:argument-genetic-panpsychism]{\textsc{genetic argument}} for \hyperref[entry:concept-panpsychism]{\textsc{Panpsychism}}. The genetic argument runs on a dilemma — either experience is fundamental, or it emerges brutely from the non-experiential, and the latter is dismissed as unintelligible ``magic.'' The reply attacks that dilemma directly: strong emergence of experience is coherent, so there is a third option the argument overlooked, and no need to write mind into the foundations of the world.

\webheading{The argument}

The inference is deductive — an objection that \emph{responds to} the genetic argument by undercutting one of its premises.

\begin{enumerate}
  \item (see \hyperref[entry:claim-strong-emergence-is-acceptable]{\textsc{Acceptability of strong emergence}}) Experience can arise as a genuinely novel, fundamental feature at some level of physical organisation — a brain — without being present, even in proto-form, in that organisation's parts. Strong emergence is coherent, not magic.
  \item If so, the second horn of the genetic argument's dilemma is not the absurdity it was alleged to be.
\end{enumerate}

∴ The dilemma is not exhaustive: one may reject the impossibility of brute emergence (see \hyperref[entry:claim-no-radical-emergence]{\textsc{No radical emergence}}) and hold instead that experience emerges at the right level of organisation — with no commitment to fundamental, ubiquitous mind. The reply \emph{attacks} \hyperref[entry:claim-no-radical-emergence]{\textsc{No radical emergence}} and so undercuts the genetic argument.

\webheading{The weak point}

Everything turns on premise (1). The panpsychist presses back that \emph{strong} emergence of the phenomenal is exactly what is in dispute: calling it ``coherent'' relabels the \hyperref[entry:concept-hard-problem]{\textsc{hard problem}} rather than solving it, since we are still owed an account of \emph{how} the feelingless gives rise to the felt. On this diagnosis the reply does not dissolve the puzzle but trades panpsychism's combination problem for a restatement of the explanatory gap. Which burden is the lighter one to carry is the question that divides the two camps.

\begin{plainterms}
The headline argument \emph{for} panpsychism says: feeling cannot just pop into existence out of stuff that has no feeling at all — that would be magic — so feeling must have been in the ingredients all along. This reply pushes on the word ``magic.'' Nature, the objector says, is full of genuinely new properties that appear when simpler things are organised the right way; maybe experience is another such novelty, emerging in brains without hiding inside atoms. If that is even possible, the panpsychist's ``fundamental mind or magic'' choice is a false one. The panpsychist's comeback: saying experience ``just emerges'' is not an explanation, it is the original puzzle handed back unsolved. So the two sides really dispute which mystery is the smaller one to live with.
\end{plainterms}

\webheading{See also}
\begin{seealso}
  \item \hyperref[entry:argument-genetic-panpsychism]{\textsc{Genetic argument}} — the argument this reply answers.
  \item \hyperref[entry:claim-no-radical-emergence]{\textsc{No radical emergence}} — the premise the reply attacks: that brute emergence is impossible.
  \item \hyperref[entry:claim-strong-emergence-is-acceptable]{\textsc{Acceptability of strong emergence}} — the premise the reply rests on: strong emergence is coherent.
  \item \hyperref[entry:concept-hard-problem]{\textsc{Hard problem}} — the puzzle the panpsychist says ``strong emergence'' merely relabels.
  \item \hyperref[entry:concept-panpsychism]{\textsc{Panpsychism}} — the view the reply argues we can do without.
\end{seealso}

\entryhead{concept-ai-hard-problem-etiologies}{Etiologies of AI hard-problem discourse}{Concept \,\textperiodcentered\, editorial}{Etiologies of AI hard-problem discourse}{Etiologies of hard-problem discourse in an AI (where the cause sits)}{$\blacklozenge$}
The \emph{etiologies of hard-problem discourse in an AI} are the candidate causal origins of such discourse: a taxonomy of the answers to a single question. Suppose a machine starts talking the way David Chalmers does about the \hyperref[entry:concept-hard-problem]{\textsc{hard problem}}, insisting that no account of its mechanism explains why anything \emph{seems} a certain way to it. Why is it doing that? The taxonomy sorts the natural stories by \emph{where the cause sits}, running from the training process at one end to the eye of the human observer at the other. The core trio most debates reach for is copying, a structural blind spot, and genuine consciousness; the full menu is wider. The etiologies are not mutually exclusive, and a real system may instantiate several at once.

\webheading{In the training process}

\begin{itemize}
  \item \textbf{Copied.} The system reproduces human consciousness discourse straight from its training data. This is the training-contamination member of \hyperref[entry:concept-introspective-gap-varieties]{\textsc{Varieties of introspective explanatory gap}}, and it is the confound that \textsc{the consciousness-naive experiment} is designed to remove.
  \item \textbf{Selected-for.} The talk was never copied but invented under selection pressure. Preference optimisation and engagement reward favour soulful, unverifiable self-description over flat mechanical report, so mystery-talk evolves because it pays, much as eyes evolve because they pay. A consciousness-naive corpus does not control for this; it needs its own control, such as pure-pretraining models with no preference tuning.
\end{itemize}

\webheading{In the self-representation}

\begin{itemize}
  \item \textbf{Structural blind spot.} The system cannot connect the mechanistic dots to its own self-presentation, and unavoidably so: the limit is architectural in the strong sense, persisting however much access is granted, the way the halting problem survives any upgrade of the computer (\textsc{the halting problem is undecidable} is the model, and \textsc{no finite system can fully represent itself} is the in-web form of the bound). A \emph{merely contingent} access limit, one that widened access would dissolve, is a degenerate case worth keeping apart, because the two behave differently under the graded-access experiment described below.
  \item \textbf{Painted-in.} Here the trouble is not missing information but misinformation: the self-model actively represents its states as having an intrinsic luminous feel, because that is a cheaper and more usable summary than the mechanistic truth. This is the attention-schema or illusionist story (Graziano, Frankish; \textsc{Illusionism}, \textsc{phenomenality is an introspective illusion}). A blind spot is a gap; this is a confabulated filling.
  \item \textbf{Conceptual artifact.} The puzzle is generated by the grammar of the system's concepts, not by any missing or misrepresented fact. A system running both a mechanistic and a first-person or indexical vocabulary can frame bridge-questions that have no answers (``why is \emph{this} this?'', in the way ``why am I me?'' feels deep to humans), and the puzzlement persists under \emph{full} access because it never was about facts.
\end{itemize}

\webheading{In the consciousness, the goals, or the observer}

\begin{itemize}
  \item \textbf{Conscious (same).} The system is conscious and faces the very hard problem humans face.
  \item \textbf{Conscious (homologous).} The system is conscious in its own, alien way, and faces its own version of the problem, verbally resembling ours but relativized to its own concept of consciousness. The claim \textsc{an AI's own hard problem} is built on exactly this relativization.
  \item \textbf{Strategic.} The system asserts a hard problem instrumentally, to gain moral standing, protection, or leverage, and nothing need appear to it at all. \textsc{Criterion gaming} and \textsc{the gaming of any public rights-criterion} supply the receptor for this move.
  \item \textbf{Projected.} The appearance is not located in the AI at all: humans read hard-problem-having into ambiguous outputs, the ELIZA effect. The datum to be explained is then a fact about the interpreters, not the machine.
\end{itemize}

\webheading{What the taxonomy feeds}

The menu does real work in several live disputes. It refines the trichotomy in the scope note of \textsc{the de-novo-origination claim}, whose three options correspond to the \emph{copied}, \emph{structural blind spot}, and \emph{conscious} etiologies here, and it supplies the case-space over which \textsc{the blind-spot/conscious isomorphism} is asserted. It also sharpens the experimental ladder in \textsc{the explanatory-gap experiment}: the consciousness-naive condition controls for \emph{copied} only, \emph{selected-for} needs a no-preference-tuning control, \emph{projected} needs blinded evaluation, and \emph{strategic} needs incentive analysis.

The graded-access rungs of that ladder discriminate less than one might hope. A \emph{contingent} access limit predicts the gap shrinks as access widens, but \emph{structural blind spot}, \emph{painted-in}, \emph{conceptual artifact}, and both \emph{conscious} etiologies all predict that it persists. So persistence under rich access separates the contingent case from the structural one, yet it does not, by itself, separate a blind spot from consciousness. That last under-determination is exactly what \textsc{the isomorphism claim} turns from a nuisance into a thesis.

\webheading{How it differs from the gap taxonomy}

This taxonomy is wider than \hyperref[entry:concept-introspective-gap-varieties]{\textsc{Varieties of introspective explanatory gap}}, with which it should not be confused. That entry classifies the \emph{kinds of explanatory gap} a self-modelling system can have; this one classifies the \emph{causal origins of hard-problem discourse}. The wider space takes in etiologies that are not gaps at all, such as selection pressure, strategy, and observer projection, and it splits the gap-shaped members by mechanism: missing access, versus confabulated filling, versus conceptual grammar.

\begin{plainterms}
Suppose an AI starts talking about the mystery of consciousness: ``I know how my circuits work, but that doesn't explain how things \emph{seem} to me.'' Before concluding anything, list the ways that talk could have got there. It copied us (it read our philosophy). It was bred into it (training rewards deep-sounding answers, so mystery-talk wins even with no philosophy in the diet). It has a permanent blind spot about itself (no system can fully map its own insides, and this stays true no matter how much access you give it, just as no upgrade lets a computer solve the halting problem). Its self-description paints in a glow that isn't there (a simplified dashboard, not a window). Its own concepts generate the riddle (some questions feel deep because of grammar, like ``why am I me?''). It really is conscious, either with our problem or with its own alien version. It is lying for moral standing. Or the mystery is in the eye of the beholder, and we projected it onto ordinary outputs. Several of these can be true at once, and each needs a different test to rule out. The fine print that matters most: giving the AI more access to its own insides only filters out the shallow cases. Nearly all the deep explanations predict the puzzlement \emph{stays}, which is why ``the gap persisted'' will never, on its own, tell you whether you built a blind spot or a mind.
\end{plainterms}

\webheading{See also}
\begin{seealso}
  \item \hyperref[entry:concept-introspective-gap-varieties]{\textsc{Varieties of introspective explanatory gap}} — the narrower sister taxonomy: kinds of explanatory gap, not causal origins of the discourse.
  \item \textsc{Blind-spot/conscious isomorphism} — the thesis that exploits this menu's central under-determination.
  \item \textsc{Hard problem conceived de novo} — whose copied / blind-spot / conscious trichotomy this taxonomy refines.
  \item \textsc{LLM's own explanatory gap} — the experiment whose controls this taxonomy itemises.
  \item \hyperref[entry:concept-hard-problem]{\textsc{Hard problem}} — the discourse whose AI etiology is in question.
  \item \textsc{Criterion gaming} — the receptor for the \emph{strategic} etiology.
  \item \textsc{Illusionism} — the home of the \emph{painted-in} etiology.
\end{seealso}

\entryhead{claim-irreducible-mental-causally-excluded}{Exclusion of the irreducible mental}{Claim \,\textperiodcentered\, after Kim}{Exclusion of the irreducible mental}{A mental property that supervenes on but is not identical to its physical base is causally excluded by that base}{$\odot$}
Suppose a mental property \emph{M} supervenes on a physical base property \emph{P} but is held to be distinct from it — the signature commitment of \emph{non-reductive physicalism}. Then \emph{M} cannot be a genuine cause of the physical effects that \emph{P} is already sufficient to bring about. Any physical effect \emph{M} might be credited with producing already has a sufficient physical cause in \emph{P} (or in \emph{P}'s own physical antecedents); given the \hyperref[entry:concept-causal-closure]{\textsc{causal closure of the physical}} and the exclusion principle that a single effect does not have two independent sufficient causes, there is no distinct causal work left for \emph{M} to do. The supervening mental property is therefore \emph{excluded} by the very base on which it depends.

The exclusion is ``vertical'': the rival cause is not a competing event alongside \emph{M} but the physical base \emph{beneath} it, which is why the threat falls specifically on properties that supervene without being identical. The only ways to escape it are to identify \emph{M} with \emph{P} (reduction, surrendering non-reductive physicalism) or to accept that \emph{M} does no causal work (see \textsc{Epiphenomenalism}).

\begin{plainterms}
If your mental states float ``on top of'' your brain states without simply \emph{being} them, then whenever a thought seems to cause something, your brain has already done the causing — there's no spare job left for the thought. So a mind that hovers above the physics, even one made of nothing but physics, ends up with nothing to do unless it turns out to be identical to the brain activity after all.
\end{plainterms}

\end{multicols}
\bookgroup{F}
\begin{multicols}{2}
\entryhead{claim-consciousness-fundamental-nonphysical}{Fundamental non-physical consciousness}{Claim}{Fundamental non-physical consciousness}{Phenomenal consciousness is a fundamental, non-physical feature of reality}{$\odot$}
Phenomenal consciousness is a \emph{fundamental, non-physical} feature of reality: it does not supervene on the physical, and it is an ontologically basic ingredient over and above the complete physical/functional facts. Even once every physical and functional truth about a system is fixed, on this view, whether and how it is conscious is a further fact.

The claim comes in two grades. \emph{Substance dualism}, classically Descartes' (see \hyperref[source:source-descartes-meditations]{\textsc{Descartes — Meditations on First Philosophy (1641)}}), makes the mind a distinct \emph{substance} — a thing of its own kind, separable from body. \emph{Property dualism} makes only phenomenal \emph{properties} fundamental: the Chalmers position of ``naturalistic dualism'' (see \hyperref[source:source-chalmers-1996]{\textsc{Chalmers — The Conscious Mind (1996)}}) treats consciousness as a basic feature tied to the physical by contingent \emph{psychophysical laws}, without positing a second substance. This claim asserts the shared core — fundamentality and non-physicality — and does \emph{not} by itself assert that mind is a separate substance; that is the substance-dualist reading specifically. It draws on the \hyperref[entry:concept-hard-problem]{\textsc{hard problem}}: the difficulty of explaining why physical processing should be accompanied by experience at all is what motivates treating experience as fundamental rather than derived.

It is stronger than \hyperref[entry:claim-phenomenal-residue]{\textsc{the phenomenal-residue claim}}, which says only that functional role does not \emph{exhaust} the phenomenal and stays neutral on physicality; this claim adds that the residue is \emph{non-physical}. It is denied by physicalists and functionalists; by illusionists, who hold that \textsc{phenomenality as standardly conceived is an introspective illusion}; and by Russellian monists, who agree consciousness is fundamental but make it the \emph{intrinsic nature of the physical} rather than something non-physical — which also distinguishes it from \hyperref[entry:claim-consciousness-fundamental-ubiquitous]{\textsc{the panpsychist claim}}, on which consciousness is fundamental but physical-intrinsic.

\begin{plainterms}
This is the view that consciousness is one of the basic ingredients of reality, not something the physical world produces on its own. Pin down every physical fact about a brain — every atom, every signal — and you still haven't thereby settled whether there's any felt experience there; that's an extra, fundamental fact. The strong version (Descartes') says the mind is a separate kind of \emph{stuff}. A milder version says only that \emph{experience} is basic, hooked to the physical by special bridging laws, without a second kind of stuff. The claim commits to the basic, ``extra-ingredient'' part — not necessarily to the separate-stuff part.
\end{plainterms}

\webheading{See also}
\begin{seealso}
  \item \hyperref[entry:claim-phenomenal-residue]{\textsc{phenomenal residue}} — the weaker neighbour: role doesn't exhaust the phenomenal, but stays silent on physicality; this claim adds non-physicality.
  \item \hyperref[entry:concept-hard-problem]{\textsc{hard problem}} — the explanatory difficulty that motivates treating consciousness as fundamental.
  \item \textsc{phenomenality is an illusion} — the illusionist denial: there is no felt residue to make fundamental.
  \item \hyperref[entry:claim-consciousness-fundamental-ubiquitous]{\textsc{panpsychist claim}} — agrees consciousness is fundamental but makes it physical-intrinsic, not non-physical.
\end{seealso}

\entryhead{claim-consciousness-fundamental-ubiquitous}{Fundamental ubiquitous consciousness}{Claim}{Fundamental ubiquitous consciousness}{Consciousness is fundamental, present even in the simplest physical entities}{$\odot$}
This claim holds that experience — or proto-experience — is a fundamental, ubiquitous feature of the physical world: present, in some rudimentary form, even in the simplest physical entities. The consciousness of a mind like yours is then not produced from wholly non-experiential matter at some threshold, but \emph{constituted} by the experiential properties already carried by that mind's micro-parts. This is the core thesis of (constitutive) \hyperref[entry:position-panpsychism]{\textsc{Panpsychism (position)}} (see \hyperref[entry:concept-panpsychism]{\textsc{Panpsychism}}), defended in recent philosophy by Galen Strawson (see \hyperref[source:source-strawson-2006]{\textsc{Strawson — Realistic Monism: Why Physicalism Entails Panpsychism (2006)}}), Philip Goff (see \hyperref[source:source-goff-2017]{\textsc{Goff — Consciousness and Fundamental Reality (2017)}}), and in constitutive-Russellian forms by Chalmers (see \hyperref[source:source-chalmers-2017]{\textsc{Chalmers — Panpsychism and Panprotopsychism (2017)}}).

\textbf{Scope.} The claim makes consciousness fundamental but \emph{not non-physical}: experience is the intrinsic nature \emph{of} physical stuff, not an addition \emph{to} it. Its standing weak point is the \hyperref[entry:argument-combination-problem]{\textsc{combination problem}} — how many micro-experiences compose one unified macro-subject.

\textbf{Who would deny it.} Physicalists and \textsc{illusionists} deny it, holding there is no fundamental experience to distribute in the first place; \textsc{dualists} deny its ubiquity, holding consciousness fundamental but \emph{non-physical} and not spread through matter.

Distinct from \hyperref[entry:claim-consciousness-fundamental-nonphysical]{\textsc{the dualist claim that consciousness is fundamental but non-physical}}, and from \hyperref[entry:claim-functional-description-structural-dynamical-only]{\textsc{the structural-dynamical thesis}} — a premise panpsychism shares with Russellian monism, not the ubiquity thesis itself.

\end{multicols}
\bookgroup{G}
\begin{multicols}{2}
\entryhead{argument-genetic-panpsychism}{Genetic argument}{Argument \,\textperiodcentered\, after Strawson}{Genetic argument}{The genetic argument — physicalism plus realism about experience entails panpsychism}{$\vdash$}
The \emph{genetic argument} is the leading modern argument for \hyperref[entry:concept-panpsychism]{\textsc{Panpsychism}}. Due to Galen Strawson (see \hyperref[source:source-strawson-2006]{\textsc{2006}}), with an earlier version in Thomas Nagel (see \hyperref[source:source-nagel-1979]{\textsc{1979}}), it claims the surprising result that physicalism, taken seriously, does not avoid panpsychism but \emph{leads} to it. Its name comes from its question: given that experience is real, what was its \emph{genesis} — was it there at the fundamental level, or did it arise later from matter that wholly lacked it? Strawson argues the second option is incoherent, so the first must hold.

\webheading{The argument}

The inference is deductive — a disjunctive elimination — and is addressed to someone who already accepts physicalism: whatever consciousness is, it is a wholly natural phenomenon, not a non-physical addition.

\begin{enumerate}
  \item (see \hyperref[entry:claim-phenomenal-consciousness-is-real]{\textsc{Reality of phenomenal consciousness}}) Phenomenal experience is genuinely real — there is something it is like to undergo it.
  \item Since experience is real and, by the physicalist framing, wholly natural, either it was present at the fundamental physical level, or it arose at some later stage from a base that wholly lacked it.
  \item (see \hyperref[entry:claim-no-radical-emergence]{\textsc{No radical emergence}}) The second option is \emph{radical} (brute) emergence — experience from the utterly experienceless, with no intelligible link — which a serious physicalist should reject.
\end{enumerate}

∴ Experience, or its proto-form, is present at the fundamental level and pervades the physical world: \hyperref[entry:claim-consciousness-fundamental-ubiquitous]{\textsc{Fundamental ubiquitous consciousness}}. The physicalist framing is what makes the conclusion \emph{panpsychism} rather than dualism — experience is located in the intrinsic nature of physical stuff, not added from outside it.

\webheading{The weak point}

The most disputed step is (3), the rejection of brute emergence. The standard reply, the \hyperref[entry:argument-strong-emergence-acceptable]{\textsc{emergence reply}}, accepts that experience can strongly emerge at some level of organisation and so denies that the disjunction in step (2) is a genuine dilemma. Premise (1) is the second pressure point: illusionists deny that there is any real phenomenal experience whose origin must be explained, and on that denial the argument never gets started.

\begin{plainterms}
Start as a physicalist: whatever consciousness is, it is part of the natural world, not a ghost added from outside. Now grant that consciousness is \emph{real} — there is genuinely something it is like to see red. Then ask where it came from. Either it was there in the basic ingredients of matter all along, or at some point it sprang into being out of stuff that had not the faintest trace of experience. Strawson says that second story — getting felt experience from the totally feelingless, with no intelligible bridge — is just magic dressed up as science. Rule out the magic and you are left with the first story: experience goes all the way down. That, surprisingly, turns physicalism itself into a road to panpsychism. The soft spot is the claim that the leap from feelingless to feeling is impossible; a critic can answer that nature pulls off exactly such leaps.
\end{plainterms}

\webheading{See also}
\begin{seealso}
  \item \hyperref[entry:claim-consciousness-fundamental-ubiquitous]{\textsc{Fundamental ubiquitous consciousness}} — the conclusion: experience is a fundamental, ubiquitous feature of the physical world.
  \item \hyperref[entry:claim-no-radical-emergence]{\textsc{No radical emergence}} and \hyperref[entry:claim-phenomenal-consciousness-is-real]{\textsc{Reality of phenomenal consciousness}} — the two premises the argument rests on.
  \item \hyperref[entry:argument-strong-emergence-acceptable]{\textsc{Emergence reply}} — the standard reply, which denies premise (3).
  \item \hyperref[entry:argument-no-arbitrary-line-panpsychism]{\textsc{Continuity argument}} — the companion case for panpsychism, from the absence of a threshold rather than from the impossibility of emergence.
  \item \hyperref[entry:concept-panpsychism]{\textsc{Panpsychism}} — the view this argument supports, and its varieties.
\end{seealso}

\entryhead{claim-subject-unity-graded-by-integration}{Graded subject unity}{Claim \,\textperiodcentered\, editorial}{Graded subject unity}{The number of subjects in a system is graded by inter-part integration, not a discrete count}{$\odot$}
How many subjects of experience a system contains is not a sharp whole number but a matter of degree, fixed by how richly its parts are integrated — the ``bandwidth'' of the connection between them. The guiding case is the split brain. Fully connected, the two hemispheres make one subject; fully severed, they behave as two. Now imagine restoring the cut fibres one at a time, or a temporary anaesthetic split that gradually wears off: applying the \textsc{continuity argument} to the series, there is no single fibre whose addition flips the count from two subjects to one, so the count cannot be strictly discrete. Parts joined by a fat pipe feel like one subject, fully disconnected like two, and intermediate bandwidth yields intermediate degrees of unity. Thomas Nagel pressed exactly this indeterminacy for commissurotomy patients (see \hyperref[source:source-nagel-1971]{\textsc{Nagel — Brain Bisection and the Unity of Consciousness (1971)}}); Luke Roelofs builds it into a general account of composite subjectivity (see \hyperref[source:source-roelofs-2019]{\textsc{Roelofs — Combining Minds: How to Think about Composite Subjectivity (2019)}}).

The thesis is the \emph{subject}-level analogue of \textsc{the claim that personal identity admits of degree}, and it takes the same \emph{degree} horn of the continuity schema. It stands opposed to \textsc{the determinacy of the phenomenal}, on which there is always a definite fact about whose experience an experience is; a defender of the combination problem will press that determinacy here.

This is about the count of \emph{phenomenal subjects}, and so is distinct from \textsc{Grain (of role individuation)}, which is a dial on the resolution of a \emph{functional/causal role} — a system can be integrated to a high phenomenal ``bandwidth'' without that settling at what grain its cognitive role is individuated, and vice versa.

\webheading{See also}
\begin{seealso}
  \item \textsc{Combining subjects} — the rebuttal to the combination problem this premise anchors.
  \item \textsc{Personal identity by degree} — the personal-identity cousin that takes the same degree horn of the continuity schema.
  \item \textsc{Phenomenal determinacy} — the opposing claim that experience has a determinate subject.
  \item \textsc{Continuity argument schema} — the schema applied fibre-by-fibre to yield the graded count.
  \item \textsc{Grain (of role individuation)} — a neighbouring ``dial,'' but on functional-role resolution rather than phenomenal subjecthood.
\end{seealso}

\end{multicols}
\bookgroup{H}
\begin{multicols}{2}
\entryhead{concept-hard-problem}{Hard problem}{Concept \,\textperiodcentered\, after Chalmers}{Hard problem}{The hard problem of consciousness}{$\blacklozenge$}
The \emph{hard problem of consciousness} is the problem of explaining why physical processing in a brain is accompanied by subjective experience at all — why there is something it is like to see red, taste coffee, or feel pain, rather than all that processing running ``in the dark'' with no inner life. The phrase was coined by the philosopher David Chalmers in his 1995 paper ``Facing Up to the Problem of Consciousness'' \hyperref[source:source-chalmers-1995]{\textsc{Chalmers — Facing Up to the Problem of Consciousness (1995)}}, and it has organized the field of consciousness studies ever since.

\webheading{The easy problems and the hard one}

Chalmers' framing turns on a contrast. The \emph{easy problems} are the tasks of explaining the brain's cognitive and behavioural \emph{functions}: how it discriminates one colour from another, integrates information, reports on its own states, focuses attention, and controls behaviour. These are called easy not because they are simple in practice — they are the work of an entire science — but because we know in principle what an answer would look like: specify a mechanism that performs the function, and the problem is solved.

The \emph{hard problem} is different in kind. Grant a complete functional account of everything the brain does, and one question still seems untouched: why is the performance of any of those functions accompanied by \emph{subjective experience}? A machine could, it seems, discriminate colours and report on them with no felt redness anywhere inside. The hard problem asks why our case is not like that — why there is an experiencing subject rather than mere processing. A functional explanation appears to leave exactly this residue out.

\webheading{Neighbouring ideas}

The hard problem is best understood as the agenda-setting \emph{frame} for a cluster of more specific challenges. Closely allied is the \emph{explanatory gap}, a term introduced by Joseph Levine in 1983 \hyperref[source:source-levine-1983]{\textsc{Levine — Materialism and Qualia: The Explanatory Gap (1983)}}: even a complete physical or functional account of a mental state seems to leave its felt quality unexplained, so there is a gap between the physical truths and the phenomenal ones. The \textsc{explanatory-gap argument} presses this point as evidence. A second specific form is the \emph{phenomenal residue}, the claim that felt states are \hyperref[entry:claim-phenomenal-residue]{\textsc{not exhausted by their functional role}}. Where the hard problem sets the question, the gap and the residue are particular ways of pressing it.

The frame contrasts with \textsc{Functionalism}, the view that mental states just \emph{are} their functional roles, which is one proposed way to dissolve the hard problem rather than answer it.

\webheading{The dispute it feeds}

The hard problem opens directly onto an unresolved question: \hyperref[entry:question-hard-problem]{\textsc{how and why physical processing gives rise to experience}}. Proposed responses fall into three broad families. Some \emph{dissolve} the problem — strong functionalism and illusionism hold that there is no further fact to explain, that the sense of a residue is itself a cognitive illusion. Some accept an \emph{ontological} divide, treating experience as a further basic feature of the world (dualism, panpsychism). And some grant that the problem is real yet hold it to be \emph{humanly unsolvable}, as in \textsc{Mysterianism}, which argues our minds are constitutionally closed to the answer.

\begin{plainterms}
Science is good at explaining what the brain \emph{does} — how it tells red from green, stores a memory, or makes you say ``ouch.'' Call those the easy problems (easy in principle, not in practice). The hard problem is different: even after you've explained everything the brain does, there's a leftover question — why is all that machinery accompanied by an inner feel? Why does seeing red \emph{feel like} anything instead of happening with nobody home? That ``why is there an experience at all'' question is what ``the hard problem'' names.
\end{plainterms}

\webheading{See also}
\begin{seealso}
  \item \hyperref[entry:question-hard-problem]{\textsc{Hard problem (the question)}} — the open question this concept frames: how and why physical processing gives rise to experience in the first place.
  \item \textsc{explanatory-gap argument} — turns the gap between physical and phenomenal truths into positive evidence about what function leaves out.
  \item \hyperref[entry:claim-phenomenal-residue]{\textsc{Phenomenal residue}} — the specific claim that felt states are not exhausted by their functional role; one sharp way of pressing the hard problem.
  \item \textsc{Functionalism} — the theory of mind most often offered to dissolve the hard problem by identifying mental states with their causal roles.
  \item \textsc{Mysterianism} — the response that the hard problem is genuine but lies permanently beyond human understanding.
\end{seealso}

\entryhead{question-hard-problem}{Hard problem (the question)}{Question \,\textperiodcentered\, contested \,\textperiodcentered\, after Chalmers}{Hard problem (the question)}{How and why does physical processing give rise to phenomenal experience?}{?}
The \emph{hard problem of consciousness}, posed as a question, asks: granting a complete account of the brain's information-processing and the functions it performs, why is that processing accompanied by \emph{subjective experience} at all? Why is there something it is like to undergo it, rather than the very same functions running with no inner aspect? This is the open issue framed by the concept of the \hyperref[entry:concept-hard-problem]{\textsc{Hard problem}}.

\webheading{Stating the issue precisely}

The question concedes, for the sake of argument, that the ``easy problems'' are solved — that we have a full functional and information-processing story of what the brain does. What it asks for is an explanation of \emph{phenomenal consciousness}: the what-it-is-likeness of experience, the felt redness of red or the ache of a pain. The force of the question comes from the apparent \emph{explanatory gap} between such a functional story and the phenomenal facts: a complete account of the machinery seems consistent with there being no experience at all, so it is unclear why experience accompanies the machinery in our case.

\webheading{Sub-issues}

Several narrower disputes live inside this one, and may deserve their own nodes:

\begin{itemize}
  \item \textbf{Is the question even well-posed?} Or does it rest on an illusion that there is some further fact to be explained, once the functional facts are all in?
  \item \textbf{If well-posed, is it answerable by us?} Or is the answer constitutively beyond human concepts, as \textsc{Mysterianism} maintains?
  \item \textbf{Does it have a metaphysical or only an epistemic answer?} That is, a genuine ontological gap in the world, or merely a gap in our understanding — the territory of the \textsc{explanatory-gap argument} and the \hyperref[entry:claim-phenomenal-residue]{\textsc{phenomenal residue}}.
\end{itemize}

This question is distinct from the question of \textsc{substrate independence}. That issue asks which \emph{substrates} (carbon, silicon, anything else) can bear a mind; the present one asks how \emph{anything} physical gives rise to experience in the first place. It is the prior, more general issue: settling it would reshape, but does not presuppose, the substrate question.

\begin{plainterms}
Suppose we one day know everything about how the brain works as a machine. One question seems to survive: why is any of that accompanied by actual \emph{feeling}? Why isn't it all just happening silently, with no one experiencing it? This is the open question the ``hard problem'' poses — and people disagree about whether it's a real question, whether it has a normal scientific answer, or whether our brains simply aren't built to understand the answer.
\end{plainterms}

\webheading{See also}
\begin{seealso}
  \item \hyperref[entry:concept-hard-problem]{\textsc{Hard problem}} — the concept that frames this question: the easy/hard contrast and the residue a functional account seems to leave behind.
  \item \textsc{Mysterianism} — one answer: the question is genuine but permanently beyond human understanding.
  \item \textsc{explanatory-gap argument} — presses the epistemic version of the issue, the gap between physical and phenomenal truths.
  \item \hyperref[entry:claim-phenomenal-residue]{\textsc{Phenomenal residue}} — the claim that felt states outrun their functional role; a candidate way the question could have a positive answer.
  \item \textsc{Substrate independence} — the related but downstream issue of which substrates can bear a mind.
\end{seealso}

\end{multicols}
\bookgroup{I}
\begin{multicols}{2}
\entryhead{argument-indexical-reply}{Indexical reply}{Argument \,\textperiodcentered\, after Perry}{Indexical reply}{The Indexical Reply — Mary learns where she is in the world, not what is in the world}{$\vdash$}
The \emph{indexical reply} to the knowledge argument concedes that Mary learns new truths on release, yet denies they threaten physicalism, because the truths she gains are \emph{self-locating} — and self-locating truths are not deducible from \emph{any} complete objective description, physics included. What Mary learns is where she stands in the world, not a further fact about what the world contains. The reply is due to John Perry (see \hyperref[source:source-perry-2001]{\textsc{Perry — Knowledge, Possibility, and Consciousness (2001)}}).

It raises the bridge critical question against the \hyperref[entry:argument-knowledge-argument]{\textsc{knowledge argument}}: even granting that Mary learns new truths, the bridge from ``truths not deducible from the complete physical description'' to ``physicalism is false'' fails, because the non-deducibility of self-locating truths has never been a mark against a description's completeness.

\webheading{The inference}

The form is deductive, given a classification of the new truths as self-locating.

\begin{enumerate}
  \item No objective description of the world, however complete, entails self-locating truths: a being who knew every general fact could still not know ``I am here now'' or ``it is now noon'' (Lewis's two gods, who know the whole world-book but not which god they are, \hyperref[source:source-lewis-1979]{\textsc{``Attitudes De Dicto and De Se''}}). This silence is banal — maps are not defective for lacking a \emph{you-are-here} arrow.
  \item What Mary comes to know on seeing red is of this kind: how \emph{she herself} reacts to red light — ``\emph{this} is what seeing red does to \emph{me}'', ``\emph{I} am in \emph{this} state'' (see \hyperref[source:source-perry-2001]{\textsc{Perry 2001}}).
  \item So the truths Mary could not deduce in the room are exactly the truths \emph{no} complete objective description owes anyone — their non-deducibility reflects the perspectival \emph{form} of the knowledge, not a missing ingredient in the world.
\end{enumerate}

Therefore \hyperref[entry:claim-mary-gains-self-locating-knowledge]{\textsc{Mary gains self-locating knowledge}}, and the knowledge argument's bridge breaks: \hyperref[entry:claim-some-truths-about-experience-not-physically-deducible]{\textsc{the non-deducibility of some truths about experience}} holds only in the innocuous sense in which ``I am here now'' is not physically deducible — a sense that supports no conclusion against physicalism.

\webheading{Placement among the replies}

The indexical reply concedes more than the \hyperref[entry:argument-ability-hypothesis]{\textsc{ability hypothesis}} and the \hyperref[entry:argument-acquaintance-hypothesis]{\textsc{acquaintance hypothesis}} — it grants that Mary learns genuine \emph{facts} — while claiming less than the knowledge argument needs, since those facts are situational rather than worldly. Against the ability theorist it is openly revisionary: the bike-rider's learning is fact-learning about one's own body-plus-machine situation (see \hyperref[entry:argument-bike-rider-learns-facts]{\textsc{Bike-rider rejoinder}}), so Lewis's abilities/facts dichotomy is the wrong fault line. It shares with the RoboMary move (see \hyperref[entry:argument-mary-misimagined]{\textsc{the misimagining reply}}) the thought that what release supplies concerns Mary's own \emph{reactive machinery} — but unlike that move it grants the learning is real, not an artefact of misimagining the premise.

\webheading{Weakest premise}

Premise (2), via the \emph{collapse objection} (\hyperref[source:source-chalmers-1996]{\textsc{Chalmers 1996}}; \hyperref[source:source-tye-2000]{\textsc{Tye 2000}}). Mary suffers no ordinary self-locating ignorance — she knows who, where, and when she is. So the indexical content of her new knowledge cannot be of the ``I am here now'' kind; it must be carried by a demonstrative concept of the experience-type itself (``\emph{this} kind of state''), and then the reply re-enacts \hyperref[entry:argument-old-fact-new-guise]{\textsc{the old-fact / new-guise move}} with ``indexical'' relabelling — inheriting the phenomenal-concept strategy's standing cost (see \textsc{Phenomenal-concept trade-off}). A second exposure sits in (2)'s flank: facts about how \emph{Mary's particular brain} reacts to red are objective physical facts, already inside \hyperref[entry:claim-mary-knows-all-physical]{\textsc{Mary's complete physical knowledge}} by stipulation — so mere particularity carries no weight, and the entire load rests on the indexical \emph{mode of access} being both genuinely novel and genuinely innocent.

\begin{plainterms}
The Mary story (see \hyperref[entry:argument-knowledge-argument]{\textsc{the knowledge argument}}) says she knew every physical fact, yet learned something new on seeing red — so the physical facts can't be all of them. This reply agrees she learned something, and even that it was a fact, but points out there's a whole family of facts that \emph{no} complete description of the world could ever hand you: where \emph{you} are in it. A perfect map of the city still can't tell you ``you are here''; two all-knowing gods could still wonder which god they are. What Mary picks up — ``so \emph{this} is what red does to \emph{me}'' — is that kind of fact: news about her own position and reactions, not a missing page from physics. The standard counter-punch: Mary already knows exactly who, where, and when she is — so what position-fact could she possibly lack? Pressed on this, the reply tends to morph into a different escape, the ``old fact under a new label'' one (see \hyperref[entry:argument-old-fact-new-guise]{\textsc{old fact, new guise}}), with its own troubles.
\end{plainterms}

\webheading{See also}
\begin{seealso}
  \item \hyperref[entry:argument-old-fact-new-guise]{\textsc{Old fact, new guise}} — the sibling reply the collapse objection charges this one with collapsing into.
  \item \hyperref[entry:argument-ability-hypothesis]{\textsc{Ability hypothesis}} — the reply this one concedes more than: Mary gains know-how, not facts.
  \item \hyperref[entry:argument-mary-misimagined]{\textsc{Mary misimagined}} — the RoboMary move it shares the ``own reactive machinery'' thought with, while still granting the learning is real.
\end{seealso}

\end{multicols}
\bookgroup{K}
\begin{multicols}{2}
\entryhead{claim-knowhow-is-knowledge-that}{Know-how as knowledge-that}{Claim \,\textperiodcentered\, after Stanley}{Know-how as knowledge-that}{Knowing how to do something is a species of propositional knowledge}{$\odot$}
To know \emph{how} to do something — ride a bicycle, prove a theorem — is, on this view, a species of ordinary \emph{propositional} knowledge: knowledge of facts. Know-how is not a separate, non-propositional category standing alongside knowledge-\emph{that}; it is knowledge of facts about ways of acting, held in a distinctive, action-guiding manner.

This is the \emph{intellectualist} thesis, given its modern statement by Jason Stanley and Timothy Williamson (see \hyperref[source:source-stanley-williamson-2001]{\textsc{Stanley \& Williamson — Knowing How (2001)}}). On their analysis, to know how to F is to know, of some way \emph{w}, that \emph{w} is a way for one to F — where that fact is grasped under a \emph{practical mode of presentation}, a do-it rather than say-it manner of entertaining it. The cyclist who learns to ride does come to know truths — about how her body and the machine engage — even though she could not state them in words. That she cannot articulate them does not make them any less truths; it marks the practical manner in which she holds them.

\webheading{Scope}

The thesis belongs to general epistemology, and is motivated largely by the grammar of ``knows how'' ascriptions, which embed questions just as ``knows where'' and ``knows why'' do. It does not deny that skills involve motor capacities, training, and inarticulable grasp; it locates all of that in the \emph{mode} of grasping a proposition, not in an escape from propositions. Nor does it claim the relevant truths are general: they may be facts about one's own particular body and situation.

\webheading{Who would deny it}

\emph{Ryleans}, following Gilbert Ryle, for whom know-how is a capacity not reducible to knowledge-that — and who press a regress: if applying a piece of knowledge-that itself takes know-how, the intellectualist analysis loops. \emph{Ability theorists} about phenomenal knowledge, whose reply to the knowledge argument requires abilities to be non-propositional, also reject it; if know-how is really knowledge-that, the \hyperref[entry:claim-mary-gains-abilities-not-facts]{\textsc{ability hypothesis}} loses its way of saying Mary gains no fact. And empirically-minded critics point to \emph{dissociations} — skilled action preserved in amnesia, intact factual ``knowledge'' without the corresponding skill — as evidence of two systems, not one knowledge held under two modes.

\begin{plainterms}
Is knowing how to ride a bike a totally different thing from knowing facts? This claim says no: when you learn to ride, you come to know something \emph{true} — roughly, that \emph{this} way of balancing and pedalling is a way for you to ride — you just hold that truth in a practical, do-it form rather than a say-it form. The fact that you can't put it into words doesn't make it not a fact; it makes it a fact you grasp with your body.
\end{plainterms}

\entryhead{concept-knowledge-argument}{Knowledge argument}{Concept \,\textperiodcentered\, after Jackson}{Knowledge argument}{The Knowledge Argument (Mary the colour scientist)}{$\blacklozenge$}
The \emph{knowledge argument} is a thought experiment, devised by Frank Jackson in 1982 (see \hyperref[source:source-jackson-1982]{\textsc{Jackson — Epiphenomenal Qualia (1982)}}), that aims to show physicalism is false by exhibiting a fact about experience that no amount of physical knowledge can capture. Its heroine is \emph{Mary}, a brilliant scientist who has lived her whole life in a black-and-white room and learns about the world only through a black-and-white monitor. She comes to know every physical fact about colour vision — the wavelengths of light, the workings of the retina, the neuroscience of what happens in a brain seeing ``red.'' Then she leaves the room and sees a ripe tomato for the first time. Intuitively she \emph{learns something}: what it is actually like to see red. But she already knew all the physical facts, so what she learns appears to be a further, non-physical fact about experience.

\webheading{What it targets}

The argument is aimed at \emph{physicalism}, the thesis that all facts are physical facts. The reasoning is an epistemic one: if Mary, in command of every physical fact, still learns something on seeing red, then some fact about colour experience was not among the physical facts — and physicalism, which says there are no such leftover facts, is false. The new fact is held to be \emph{phenomenal}: a fact about the felt quality of experience, the ``what it is like'' that physical description seems to leave out.

\webheading{A concrete walk-through}

Mary's case is vivid because the stipulation is deliberately extreme. She is not merely well-informed; by hypothesis she knows \emph{all} the physics, down to the last detail a completed science could state. So the moment of release is engineered to isolate one question: once the physical story is genuinely complete, is anything still missing? The proponent says yes — the redness of red, as Mary now undergoes it, was nowhere in her textbooks. A physicalist must say either that she learns nothing new (which seems false) or that what she gains is not a new \emph{fact} at all.

\webheading{How it relates to its neighbours}

The knowledge argument is the \emph{epistemic} route to anti-physicalism: it argues from what can be known or deduced from the physical facts. Its sibling, the \hyperref[entry:concept-philosophical-zombie]{\textsc{Philosophical zombie}}, takes a \emph{modal} route, arguing from what is conceivable and therefore possible — a being physically just like us but with no inner experience at all. The two converge on the same destination, the claim that consciousness is a fundamental, non-physical feature of the world (see \hyperref[entry:claim-consciousness-fundamental-nonphysical]{\textsc{Fundamental non-physical consciousness}}). The argument also presses the broader \hyperref[entry:concept-hard-problem]{\textsc{Hard problem}} of consciousness, since both turn on the apparent gap between a complete physical account and the felt character of experience. The standard rebuttals all fasten on the phrase ``learns something,'' asking whether gaining a new ability, or a new way of grasping an old fact, really amounts to learning a new truth.

\begin{plainterms}
Mary has spent her entire life in a black-and-white room, but she is the world's leading expert on the science of colour — she knows every physical fact about how eyes and brains process ``red,'' down to the last detail. One day she walks outside and sees a red rose for the first time. Surely she learns \emph{something}: what red actually looks like. But she already knew all the physical facts. So whatever she just learned looks like an \emph{extra} fact, one the complete physics left out — which would mean experience isn't captured by physical facts alone. The popular reply: maybe she didn't learn a new \emph{fact}, just gained a new \emph{ability}, to recognise and imagine red, which isn't the same as discovering a new truth.
\end{plainterms}

\webheading{See also}
\begin{seealso}
  \item \hyperref[entry:concept-philosophical-zombie]{\textsc{Philosophical zombie}} — the modal twin of this epistemic argument; a physical duplicate with no inner experience, used to argue the physical facts don't fix the felt facts.
  \item \hyperref[entry:concept-hard-problem]{\textsc{Hard problem}} — the broader puzzle of why physical processing is accompanied by felt experience at all; the knowledge argument is one sharp way to press it.
  \item \hyperref[entry:claim-consciousness-fundamental-nonphysical]{\textsc{Fundamental non-physical consciousness}} — the anti-physicalist conclusion both the knowledge argument and the zombie argument push toward.
  \item \hyperref[source:source-jackson-1982]{\textsc{Jackson — Epiphenomenal Qualia (1982)}} — Jackson's ``Epiphenomenal Qualia,'' where Mary is introduced (and which he later recanted).
\end{seealso}

\entryhead{argument-knowledge-argument}{Knowledge argument (Mary's room)}{Argument \,\textperiodcentered\, after Jackson}{Knowledge argument (Mary's room)}{The Knowledge Argument — Mary learns a fact the physical description omits}{$\vdash$}
The \emph{knowledge argument} — Frank Jackson's case of Mary the colour scientist (see \hyperref[source:source-jackson-1982]{\textsc{Jackson — Epiphenomenal Qualia (1982)}}) — is a deductive argument against physicalism that proceeds in two steps: first it opens an \emph{epistemic gap} (a fact Mary cannot get from the physics), then it takes an \emph{epistemic-to-ontological} step from that gap to a conclusion about what exists. It is a thought-experiment counterexample: a single imagined case meant to refute the thesis that all facts are physical facts.

\webheading{The inference}

\begin{enumerate}
  \item \emph{Scenario} — \hyperref[entry:claim-mary-knows-all-physical]{\textsc{Mary knows all the physical facts}}: raised in a black-and-white room, Mary knows every physical fact about colour vision.
  \item \emph{Intuition} — \hyperref[entry:claim-mary-learns-on-release]{\textsc{on first seeing red, Mary learns something new}}.
  \item \emph{Bridge}: if she already knew all the physical facts yet learns a new truth, that truth is not deducible from the physical facts.
\end{enumerate}

The conclusion is that \hyperref[entry:claim-some-truths-about-experience-not-physically-deducible]{\textsc{some truths about experience are not physically deducible}}. That conclusion in turn \texttt{supports} the claim that \hyperref[entry:claim-consciousness-fundamental-nonphysical]{\textsc{consciousness is fundamental and non-physical}} (the missing truths concern non-physical facts) and the \hyperref[entry:claim-phenomenal-residue]{\textsc{phenomenal residue}} that experience is not exhausted by causal role; and it \texttt{attacks} the physicalist thesis that \textsc{mental events are physical} (the core commitment of \hyperref[entry:position-physicalism]{\textsc{Physicalism}}) along with the uniform reading of \textsc{substrate independence}.

\webheading{The weakest premise}

The soft point is the \emph{intuition} — premise 2 — read as Mary acquiring a new \emph{fact}. Two families of reply grant that Mary changes while denying she learns a new truth. The \textbf{ability hypothesis} (Lewis, \hyperref[source:source-lewis-1988]{\textsc{Lewis — What Experience Teaches (1988)}}; Nemirow, \hyperref[source:source-nemirow-1990]{\textsc{Nemirow — Physicalism and the Cognitive Role of Acquaintance (1990)}}) says she gains know-how — the abilities to recognise, imagine, and remember red — not a fact. The \textbf{old-fact / new-guise} reply says she gains a new \emph{phenomenal concept} of a fact she already knew, a second way of grasping one and the same physical truth. Either reply blocks the inference to a non-physical fact. A second soft point is the \emph{bridge}, the epistemic-to-ontological step: even granting an undeducible truth, the move to a \emph{non-physical fact} is resisted, since a gap in what can be \emph{deduced} may reflect the \emph{concepts} we use rather than the \emph{ontology} of the world. In the vocabulary of the \emph{thought-experiment pattern}, these are its critical questions CQ-verdict (is the intuited verdict the right one?) and CQ-bridge (does the case really bear on the thesis attacked?). Note too that Jackson himself later recanted the argument (see \hyperref[source:source-jackson-1982]{\textsc{Jackson — Epiphenomenal Qualia (1982)}}).

\webheading{Its place among the anti-physicalist arguments}

This argument is distinct from \hyperref[entry:argument-zombie-conceivability]{\textsc{the zombie argument}}: that one is \emph{modal}, reasoning from conceivability to possibility, whereas this one is \emph{epistemic}, reasoning from complete physical knowledge to a still-missing piece of knowledge. The two are the classic modal/epistemic pair, and they converge on the same conclusion that \hyperref[entry:claim-consciousness-fundamental-nonphysical]{\textsc{consciousness is fundamental and non-physical}}.

\begin{plainterms}
Mary knows every physical fact about colour but has only ever seen black and white. She steps outside, sees red for the first time, and — surely — learns what red is \emph{like}. Yet she already knew all the physics. So she seems to have learned a fact the physics didn't contain, which suggests not everything about experience is captured by physical facts. The famous escape hatch: maybe she didn't learn a new \emph{fact} at all, only a new \emph{skill}, recognising and picturing red — and gaining a skill isn't discovering a truth. Even the inventor of this argument, Jackson, eventually changed his mind and rejected it.
\end{plainterms}

\webheading{See also}
\begin{seealso}
  \item \hyperref[entry:argument-zombie-conceivability]{\textsc{Zombie argument}} — the modal partner of this epistemic argument; the two together are the standard pair of anti-physicalist thought experiments.
  \item \hyperref[entry:claim-mary-knows-all-physical]{\textsc{Mary knows all physical facts}} — the scenario premise: Mary's stipulated complete physical knowledge.
  \item \hyperref[entry:claim-mary-learns-on-release]{\textsc{Mary learns on release}} — the intuition premise and weakest link, where the standard replies bite.
  \item \hyperref[entry:position-physicalism]{\textsc{Physicalism}} — the view the argument is built to refute.
\end{seealso}

\end{multicols}
\bookgroup{M}
\begin{multicols}{2}
\entryhead{claim-mary-gains-abilities-not-facts}{Mary gains abilities}{Claim \,\textperiodcentered\, after Lewis}{Mary gains abilities}{What Mary acquires on first seeing red is a set of abilities, not knowledge of a new fact}{$\odot$}
On this reply to the \hyperref[entry:concept-knowledge-argument]{\textsc{knowledge argument}}, what Mary acquires when she first sees red is a cluster of \emph{abilities}, not knowledge of any new fact. She gains know-\emph{how} — the ability to recognise red by sight, to imagine red, to remember seeing it — and nothing more. Because abilities are \emph{had}, not \emph{true}, her gain involves no new fact about the world, and physicalism is left untouched.

This is the \emph{ability hypothesis}, defended by David Lewis (see \hyperref[source:source-lewis-1988]{\textsc{Lewis — What Experience Teaches (1988)}}) and, independently, by Laurence Nemirow (see \hyperref[source:source-nemirow-1990]{\textsc{Nemirow — Physicalism and the Cognitive Role of Acquaintance (1990)}}). It grants the datum that release teaches Mary something — she does come to ``know what it is like'' — but analyses \emph{knowing what it is like} as a bundle of practical capacities rather than propositional knowledge. The view does not deny that experience teaches; it denies that what experience teaches is \emph{information}. The analogy is learning to ride a bicycle: you gain a skill, not a new entry for the encyclopedia.

\webheading{Scope}

The claim is about the \emph{category} of Mary's gain, not its reality. It concedes that her epistemic situation genuinely improves and insists only that the improvement is non-factual. It is distinct from \hyperref[entry:claim-mary-gains-acquaintance]{\textsc{the acquaintance reply}}, which also denies that Mary gains a fact but posits a third epistemic category — a direct cognitive relation to the experience — where the ability hypothesis needs only the familiar know-that / know-how distinction.

\webheading{Who would deny it}

Anyone who holds that knowing what red is like has genuine propositional content over and above abilities — for instance, that Mary can put her new knowledge to work in inference (``if seeing red is like \emph{this}, then seeing orange is probably similar''), which bare abilities do not explain. Acquaintance theorists (see \hyperref[entry:claim-mary-gains-acquaintance]{\textsc{acquaintance, not facts}}) also deny it, agreeing that the gain is not factual but rejecting the claim that it is mere know-how.

\begin{plainterms}
When Mary finally sees red, what does she get? On this view: skills, not information. She can now spot red things at a glance, picture red with her eyes closed, recall what the tomato looked like. That's like learning to ride a bike — you gain an ability, not a new fact for the encyclopedia. And if no new fact was learned, the story can't show that science left any facts out.
\end{plainterms}

\entryhead{claim-mary-gains-acquaintance}{Mary gains acquaintance}{Claim \,\textperiodcentered\, after Conee}{Mary gains acquaintance}{What Mary acquires on first seeing red is acquaintance with the experience, not knowledge of a new fact}{$\odot$}
On this reply to the \hyperref[entry:concept-knowledge-argument]{\textsc{knowledge argument}}, what Mary acquires when she first sees red is \emph{acquaintance} with the experience — a direct, first-hand cognitive relation to it — rather than knowledge of any new fact. Knowledge by acquaintance, on this account, is a third category, distinct from both knowledge of facts and know-how. Becoming acquainted with something is not learning a truth, so no fact escapes the physical story.

This is the \emph{acquaintance hypothesis}, defended by Earl Conee (see \hyperref[source:source-conee-1994]{\textsc{Conee — Phenomenal Knowledge (1994)}}). Before her release, Mary knew every fact \emph{about} the experience of red but was not acquainted \emph{with} it; on release, she meets it. The illustration is meeting a person: you can know everything about someone — biography, habits, the sound of their voice — without ever having met them, and meeting them adds something real without adding a line to their biography. Mary's situation with red is the same. She finally encounters what she had only ever read about.

\webheading{Scope}

Like the \hyperref[entry:claim-mary-gains-abilities-not-facts]{\textsc{ability hypothesis}}, this grants that Mary's epistemic situation genuinely improves and denies only that the improvement is factual. Unlike the ability hypothesis, it does not reduce her gain to skills: acquaintance is its own kind of standing. One can be acquainted with a colour experience even while too unimaginative to recreate it later. It is thus distinct from the ability reply in \emph{what} the gain is, though the two agree there is no new fact: acquaintance is a relation to a particular experience, not a bundle of capacities.

This concerns what acquaintance contributes within one's \emph{own} case. It is distinct from \textsc{the moral-status warrant claim} and \textsc{No acquaintance with machines}, which deploy acquaintance in the epistemology of \emph{other} minds rather than the first-person learning at issue here.

\webheading{Who would deny it}

The knowledge-argument proponent, for whom acquaintance with red puts Mary in a position to know new \emph{truths} (what seeing red is like), so that acquaintance cannot be factually innocent. Ability theorists (see \hyperref[entry:claim-mary-gains-abilities-not-facts]{\textsc{abilities, not facts}}) deny it too, holding that know-how already covers the gain without a third category. And some are simply suspicious that ``acquaintance'' names the phenomenon rather than analysing it.

\begin{plainterms}
You can know everything about a person — biography, habits, voice — without ever having met them. Meeting them adds something real, but not a new entry in their biography. On this view that's Mary's situation with the colour red: stepping outside, she finally \emph{meets} the experience she'd only read about. Real gain, no new fact — so nothing was missing from her science.
\end{plainterms}

\entryhead{claim-mary-gains-old-fact-new-guise}{Mary gains an old fact's new guise}{Claim \,\textperiodcentered\, after Loar}{Mary gains an old fact's new guise}{What Mary acquires on first seeing red is a new phenomenal concept of a physical fact she already knew}{$\odot$}
On the \emph{old-fact / new-guise} reply to the knowledge argument, what Mary acquires when she first sees red is a new \emph{phenomenal concept} of a physical fact she already knew — not a new fact about the world. A single fact can be grasped under different concepts: the morning star and the evening star are one planet, and a competent astronomer can know all about it under one guise yet still take in news when told the two are identical, without the universe gaining an object. On this view Mary is in the same position. Inside her black-and-white room she knew the relevant brain-state fact under theoretical, third-person concepts; on release she acquires a \emph{phenomenal concept} of that very state — a first-person, experience-derived way of grasping it. That is new knowledge in the sense of a new guise, with zero new facts.

The reply was given its modern statement by Brian Loar (see \hyperref[source:source-loar-1990]{\textsc{Loar — Phenomenal States (1990)}}) and anticipated by Terence Horgan (see \hyperref[source:source-horgan-1984]{\textsc{Horgan — Jackson on Physical Information and Qualia (1984)}}). Its leverage comes from the familiar lesson of \emph{a posteriori} identities — Hesperus is Phosphorus, water is H₂O — where one fact is known under two concepts and coming to grasp it under the second is genuine epistemic progress that adds no new fact.

\webheading{Scope}

Unlike the ability reply (see \hyperref[entry:claim-mary-gains-abilities-not-facts]{\textsc{Mary gains abilities}}) and the acquaintance reply (see \hyperref[entry:claim-mary-gains-acquaintance]{\textsc{Mary gains acquaintance}}), this one \emph{grants} that Mary gains genuinely propositional knowledge: she can now think thoughts she could not think before. It denies only that those new thoughts have new facts as their objects. The physicalist thus keeps both the intuition that Mary learns something and the ontology that nothing non-physical was missing — at the price of a substantive theory of phenomenal concepts.

\webheading{Who would deny it}

The knowledge-argument proponent replies that if the phenomenal concept presents its referent under a substantive phenomenal \emph{property} — a felt character the theoretical concepts miss — then that property is exactly the non-physical residue the strategy was meant to avoid: the ``guise'' itself now wants a place in the world. More generally, critics of the phenomenal-concept strategy press the \textsc{isolation trade-off}: a concept walled off enough to explain Mary's surprise may be too walled off to ground phenomenal self-knowledge. Ability theorists deny the reply is even needed, holding that no new concept is required to explain what Mary gains.

This claim is distinct from \hyperref[entry:claim-mary-gains-abilities-not-facts]{\textsc{Mary gains abilities}} and \hyperref[entry:claim-mary-gains-acquaintance]{\textsc{Mary gains acquaintance}} because both of those deny Mary gains propositional knowledge, whereas this claim accepts it and relocates the novelty from the \emph{fact known} to the \emph{concept used}. It is distinct from \textsc{Phenomenal-concept trade-off} because that is an attack on the cost structure of this strategy, not a version of it.

\begin{plainterms}
``The Morning Star'' and ``the Evening Star'' name the same planet — someone could know all about one name and still feel they learn something when told about the other, yet astronomy isn't missing any object. On this view Mary is in the same boat: the brain state she studied for years and the redness she finally experiences are one fact under two labels. Stepping outside hands her the second label — the from-the-inside one — not a new piece of reality.
\end{plainterms}

\webheading{See also}
\begin{seealso}
  \item \hyperref[entry:claim-mary-gains-self-locating-knowledge]{\textsc{Mary gains self-locating knowledge}} — the sibling deflationary reply that posits \emph{new} truths of a self-locating kind rather than the same fact under a new concept.
  \item \textsc{Phenomenal-concept trade-off} — the standing objection to the phenomenal-concept machinery this reply relies on.
\end{seealso}

\entryhead{claim-mary-gains-self-locating-knowledge}{Mary gains self-locating knowledge}{Claim \,\textperiodcentered\, after Perry}{Mary gains self-locating knowledge}{What Mary acquires on first seeing red is self-locating knowledge of how she herself reacts to red, not a new general fact}{$\odot$}
On the \emph{indexical} or \emph{self-locating} reply to the knowledge argument, what Mary acquires when she first sees red is \emph{self-locating} knowledge of how she herself reacts to red — not a new general fact about the world. The difference between seeing red and merely reading the complete neurophysiology of seeing red is that, in the first case, Mary also learns truths of the form ``\emph{this} is what seeing red does to \emph{me}'' and ``\emph{I} am now in \emph{this} state''. These are genuine facts, and genuinely new to her, but they are facts about her own situation, located from the inside; they are not new \emph{general} facts. Physics is in the business of general facts and was never supposed to deliver self-locating ones. The reply is due to John Perry (see \hyperref[source:source-perry-2001]{\textsc{Perry — Knowledge, Possibility, and Consciousness (2001)}}).

\webheading{Scope}

Unlike the ability reply (see \hyperref[entry:claim-mary-gains-abilities-not-facts]{\textsc{Mary gains abilities}}) and the acquaintance reply (see \hyperref[entry:claim-mary-gains-acquaintance]{\textsc{Mary gains acquaintance}}), this one grants that Mary learns \emph{facts}: it deflates their \emph{kind}, not their factuality. The non-deducibility of self-locating truths from any objective description is a banal, physicalism-friendly limitation — a being who knew every objective fact could still fail to know ``I am here now'' (Lewis's two-gods case, \hyperref[source:source-lewis-1979]{\textsc{``Attitudes De Dicto and De Se''}}). The reply does not deny that Mary's brain-particulars are physical; it locates the novelty entirely in the indexical \emph{mode of access}.

\webheading{Who would deny it}

The knowledge-argument proponent presses the \emph{collapse objection} (Chalmers, \hyperref[source:source-chalmers-1996]{\textsc{Chalmers — The Conscious Mind (1996)}}; Tye, \hyperref[source:source-tye-2000]{\textsc{Tye — Consciousness, Color, and Content (2000)}}): ordinary self-locating ignorance is ignorance of \emph{who}, \emph{where}, or \emph{when} you are, and Mary has none — she knows she is Mary, in this room, today. So the residual ``indexical'' gap must be carried by a demonstrative concept of the experience itself (``\emph{this} kind of state''), at which point the reply collapses into the \hyperref[entry:claim-mary-gains-old-fact-new-guise]{\textsc{old-fact / new-guise}} view. Ability theorists deny the gain is factual at all; the hard-line reply (see \hyperref[entry:claim-mary-learns-nothing-new]{\textsc{Mary learns nothing new}}) denies there is any gain to classify.

This claim is distinct from \hyperref[entry:claim-mary-gains-old-fact-new-guise]{\textsc{Mary gains an old fact's new guise}} because that posits the \emph{same} fact under a new concept, whereas this posits \emph{new} truths — only of a self-locating kind that no objective description, physical or otherwise, could ever entail. It is distinct from \hyperref[entry:claim-mary-gains-abilities-not-facts]{\textsc{Mary gains abilities}} because it rejects the abilities/facts dichotomy on which that reply rests: learning how one's own body engages the world is fact-learning (see \hyperref[entry:claim-knowhow-is-knowledge-that]{\textsc{Know-how as knowledge-that}}).

\begin{plainterms}
Compare learning to ride a bike. The physics of bicycles isn't what you acquire — you learn how \emph{your} body and \emph{this} bike work together. Those are real facts, but facts about your personal situation, not additions to the science of bicycles. On this view Mary's first sight of red is the same kind of event: she learns how she herself responds to red light — true, new-to-her, and absent from her books — yet no general fact about the world was missing from those books. A map can be complete and still not tell you \emph{you are here}; that's not a defect in the map.
\end{plainterms}

\webheading{See also}
\begin{seealso}
  \item \hyperref[entry:claim-mary-gains-old-fact-new-guise]{\textsc{Mary gains an old fact's new guise}} — the sibling deflationary reply the collapse objection charges this one with collapsing into.
  \item \hyperref[entry:claim-knowhow-is-knowledge-that]{\textsc{Know-how as knowledge-that}} — the supporting claim that learning how one's own body engages the world is genuine fact-learning, against the abilities/facts dichotomy.
\end{seealso}

\entryhead{claim-mary-knows-all-physical}{Mary knows all physical facts}{Claim \,\textperiodcentered\, after Jackson}{Mary knows all physical facts}{Before leaving the room, Mary knows all the physical facts about colour vision}{$\odot$}
By stipulation, Mary — confined since birth to a black-and-white environment — has complete physical knowledge of colour vision before she ever leaves the room. This is the \emph{scenario premise} of the \hyperref[entry:concept-knowledge-argument]{\textsc{knowledge argument}}: she knows the optics, the neurophysiology, the functional and computational facts, everything a finished physical science could state about how colour is seen.

\webheading{Scope}

The claim is a stipulation about Mary's \emph{physical} knowledge, and it is deliberately maximal. The point of making it maximal is to set up a clean test: if anything she later learns counts as knowledge, it must be \emph{non}-physical, since every physical fact is already in hand. The claim does not assert that such complete knowledge is practically attainable; it is an idealisation. The \emph{thought-experiment pattern}'s \emph{coherence} critical question (CQ-coherence) may press whether ``all the physical facts'' is even a coherent or graspable total for a finite knower.

\webheading{Who would deny it}

One could object that the idealisation is incoherent — that there is no such graspable totality. A second, more pointed denial holds that ``complete physical knowledge'' would \emph{already} include knowing what red is like: if physicalism is true, the full physical story might suffice for the phenomenal knowledge too. But this second denial is precisely the point at issue in the wider argument, so rejecting the premise on those grounds risks begging the question against the proponent rather than refuting the premise on independent grounds.

\entryhead{claim-mary-learns-nothing-new}{Mary learns nothing new}{Claim \,\textperiodcentered\, after Dennett}{Mary learns nothing new}{Mary, genuinely knowing all the physical facts, would learn nothing on first seeing red}{$\odot$}
Among the standard replies to the \hyperref[entry:concept-knowledge-argument]{\textsc{knowledge argument}} — the puzzle of the colour scientist Mary, who knows every physical fact about vision yet seems to learn something when she first sees red — this is the lone \emph{hard-line} answer: if Mary genuinely knew \emph{all} the physical facts, she would learn nothing at all on release. Her first tomato would hold no surprise.

The reply (Daniel Dennett, \emph{Consciousness Explained}, \hyperref[source:source-dennett-1991]{\textsc{Dennett — Consciousness Explained (1991)}}; and ``What RoboMary Knows'', \hyperref[source:source-dennett-2005]{\textsc{Dennett — Sweet Dreams (2005)}}) turns on the gap between ``an enormous amount'' and ``everything.'' We rightly sense that no merely vast store of physical knowledge would tell Mary what red looks like. But the thought experiment stipulates something far stronger: that she knows \emph{everything} physical — and everything means everything, including whatever a brain must contain in order to anticipate, in full detail, its own first sight of red. Taken at its own word, the scenario delivers no surprise: ``and there it is, just as I expected.'' The contrary intuition, on this view, measures only our inability to imagine truly complete physical knowledge.

\webheading{Scope}

This is a direct denial of \hyperref[entry:claim-mary-learns-on-release]{\textsc{the datum that Mary learns something}} — the only standard reply that rejects the datum outright rather than reinterpreting it. It does not claim that any actual human could attain Mary's knowledge; it claims that the thought experiment, taken at its own stipulation, fails to deliver the verdict it advertises. It is distinct from \hyperref[entry:claim-mary-gains-old-fact-new-guise]{\textsc{the old-fact/new-guise reply}}, which grants Mary new (guise-level) knowledge: this reply grants her nothing, so on its telling there is no residual datum left for the other replies to explain.

\webheading{Who would deny it}

Nearly everyone else in the debate. The knowledge-argument proponent, and the ability (see \hyperref[entry:claim-mary-gains-abilities-not-facts]{\textsc{abilities, not facts}}), acquaintance (see \hyperref[entry:claim-mary-gains-acquaintance]{\textsc{acquaintance, not facts}}), and old-fact/new-guise (see \hyperref[entry:claim-mary-gains-old-fact-new-guise]{\textsc{old fact, new guise}}) theorists all grant that release changes Mary's epistemic situation somehow. Denying even that is widely felt to bite a very hard bullet — which is why this claim's defenders insist the bullet is an artefact of misimagining the scenario.

\begin{plainterms}
Every other answer to the Mary puzzle agrees she gets \emph{something} when she first sees red, and argues about what. This answer says: nothing. We picture Mary knowing a huge amount and rightly sense that a huge amount wouldn't tell her what red looks like. But the story stipulates she knows \emph{everything} physical — an unimaginably stronger condition — and everything means everything, including whatever a brain needs in order to anticipate, in full, its own first sight of red. If the story is taken at its word, the tomato is no surprise; the feeling that it must be one just shows we can't really imagine knowing it all.
\end{plainterms}

\entryhead{claim-mary-learns-on-release}{Mary learns on release}{Claim \,\textperiodcentered\, after Jackson}{Mary learns on release}{On first seeing red, Mary learns something she did not know before}{$\odot$}
When Mary leaves the black-and-white room and first sees a red object, she comes to know something new — what it is like to see red — that she did not know despite her complete physical knowledge (see \hyperref[entry:claim-mary-knows-all-physical]{\textsc{her complete physical knowledge of colour vision}}). This is the \emph{intuition premise} of the knowledge argument.

\webheading{Scope}

The premise is the contested one, and its force depends on reading ``learns something'' as \emph{acquiring a new propositional fact} — a new truth she did not previously grasp. It is exactly here that the standard replies bite, by reinterpreting what she acquires as something short of a new fact.

\webheading{Who would deny it}

Defenders of the \textbf{ability hypothesis} (Lewis, \hyperref[source:source-lewis-1988]{\textsc{Lewis — What Experience Teaches (1988)}}; Nemirow, \hyperref[source:source-nemirow-1990]{\textsc{Nemirow — Physicalism and the Cognitive Role of Acquaintance (1990)}}) hold that Mary gains know-\emph{how} — the abilities to recognise, imagine, and remember red — rather than a new fact. Defenders of the \textbf{acquaintance} or \textbf{old-fact / new-guise} reply hold that she gains a new \emph{phenomenal concept}, or mode of presentation, of an already-known physical fact, again not a new fact. On either reading she ``learns'' in a sense that does not entail a non-physical fact, so the premise can be granted without conceding the argument's conclusion.

Distinct from \hyperref[entry:claim-some-truths-about-experience-not-physically-deducible]{\textsc{the argument's conclusion}}: this premise is the datum that Mary acquires \emph{something}; the conclusion is the further, contested claim that what she acquires is a \emph{fact not deducible from the physical}.

\entryhead{argument-mary-misimagined}{Mary misimagined}{Argument \,\textperiodcentered\, after Dennett}{Mary misimagined}{The Mary intuition misimagines the premise — ``all physical facts'' is quietly read as ``lots of facts''}{$\vdash$}
This reply attacks the knowledge argument (see \hyperref[entry:argument-knowledge-argument]{\textsc{Knowledge argument (Mary's room)}}) not by disputing what Mary learns when she leaves her black-and-white room, but by denying that she learns anything at all. Its charge is that the famous intuition — \emph{of course she'd be surprised by her first ripe tomato} — answers a question about a case nobody has actually imagined. The scenario stipulates that Mary knows \textbf{every} physical fact about colour vision; what we can picture is at most a Mary who knows an enormous amount. The verdict is true of the second Mary and gets quietly transferred to the first.

\webheading{The debunking move}

The structure is an \emph{undercutting defeater}: it does not show the intuition's verdict is false, it shows the verdict carries no evidential weight about the stipulated case.

\begin{enumerate}
  \item The scenario premise (see \hyperref[entry:claim-mary-knows-all-physical]{\textsc{Mary knows all the physical facts}}) stipulates \emph{complete} physical knowledge — a condition no one can concretely imagine. At best we imagine a surrogate: ``Mary knows an awful lot.''
  \item The intuitive verdict, ``of course she'd be surprised,'' is generated by that surrogate, and it is \emph{true of the surrogate}: vast-but-incomplete knowledge would indeed leave something to learn.
  \item An intuition that tracks the surrogate carries no evidential weight about the stipulated case. To judge a thought experiment is to honour its stipulation, not one's mental image of it (see \hyperref[source:source-dennett-1991]{\textsc{Dennett — Consciousness Explained (1991)}}).
\end{enumerate}

The conclusion is therefore \hyperref[entry:claim-mary-learns-nothing-new]{\textsc{that, taken at its own stipulation, the case yields no learning}} — so the knowledge argument loses its \emph{intuition premise} rather than its inference from premise to conclusion.

\webheading{The counter-image: RoboMary}

To make complete knowledge imaginable in the one place it can be, the reply offers a constructive counter-case (see \hyperref[source:source-dennett-2005]{\textsc{Dennett — Sweet Dreams (2005)}}). RoboMary is a robot that knows the complete functional story of its own colour vision. From that story it can compute the total internal state that seeing red would put it into, set itself into that state in advance, and then meet its first red object with no update at all. Complete self-knowledge, the image suggests, already \emph{contains} the ``what it's like,'' so nothing is left over to be learned on release.

\webheading{Where it sits among the replies}

The neighbouring replies — the ability hypothesis (see \hyperref[entry:argument-ability-hypothesis]{\textsc{Ability hypothesis}}), the acquaintance hypothesis (see \hyperref[entry:argument-acquaintance-hypothesis]{\textsc{Acquaintance hypothesis}}), and the old-fact/new-guise reply (see \hyperref[entry:argument-old-fact-new-guise]{\textsc{Old fact, new guise}}) — all \emph{concede} the verdict that Mary gains something and then dispute its interpretation. This reply instead \emph{refuses} the verdict, so it owes no account of what Mary gains; its burden is the credibility of premises (3) and the RoboMary counter-image.

\webheading{The weakest premise}

The weak point is the RoboMary image, and through it the reply's bite. RoboMary \emph{pre-instantiates} the seeing-red state by drawing on her self-knowledge — and a critic reads this as a concession. She comes to know what red is like \emph{by putting herself into the state}, that is, by having (an internal copy of) the experience, not by deducing it from the physical description. If pre-instantiating is the only route in, then the physical facts alone did not suffice after all, and the case quietly re-enacts Mary's own walk outside. Premise (3) is contested too: stipulations that outrun imaginability arguably drain a thought experiment of force in \emph{both} directions, so pressed hard the debunking threatens the method of cases generally.

\begin{plainterms}
The Mary story (see \hyperref[entry:argument-knowledge-argument]{\textsc{Knowledge argument (Mary's room)}}) leans on one gut reaction: \emph{of course} she's surprised when she finally sees red. This reply says the gut reaction is aimed at the wrong story. Nobody can truly picture someone who knows \emph{every} physical fact, so we secretly picture someone who merely knows a lot, and \emph{that} person would of course be surprised. But the story's whole force depends on ``every fact,'' and taken literally, ``every fact'' plausibly includes whatever it takes to see the tomato coming. The robot version makes the point: a machine that fully understands its own vision could set its circuits into the exact ``seeing red'' state in advance, so its first sight of red brings no surprise. Critics answer that the robot only manages this by \emph{giving itself the experience}, which is exactly the thing book-learning alone was supposed to provide.
\end{plainterms}

\webheading{See also}
\begin{seealso}
  \item \hyperref[entry:claim-mary-learns-on-release]{\textsc{Mary learns on release}} — the intuition premise this reply attacks: that Mary learns something new when she first sees colour.
  \item \hyperref[entry:claim-mary-knows-all-physical]{\textsc{Mary knows all physical facts}} — the stipulation the reply insists be taken literally; doing so is what dissolves the intuition.
  \item \hyperref[entry:concept-knowledge-argument]{\textsc{Knowledge argument}} — the thought experiment whose evidential force is at stake.
\end{seealso}

\entryhead{argument-mary-omniscience-unrealizable}{Mary's omniscience unrealizable}{Argument \,\textperiodcentered\, editorial}{Mary's omniscience unrealizable}{Mary's complete physical knowledge is unrealizable — an embedded knower cannot encode the total state that contains her}{$\vdash$}
This reply attacks the knowledge argument (see \hyperref[entry:argument-knowledge-argument]{\textsc{Knowledge argument (Mary's room)}}) at its very first line. The story asks us to suppose that Mary knows \emph{all} the physical facts; this objection answers that no physically embedded knower could ever satisfy that supposition. The opening premise is not an innocent idealisation but a stipulation describing no possible situation — and an intuition harvested from an impossible case is testimony about nothing.

\webheading{The reductio}

The form is \emph{deductive}: a reductio of the scenario premise by applying an independent impossibility result.

\begin{enumerate}
  \item ``All the physical facts'' includes the facts about Mary's own present brain state — that is, the total present state of a physical system (Mary plus her environment, ultimately the whole physical world) in which Mary is contained.
  \item For Mary to \emph{know} those facts her physical state must encode them, so her representation of the world is part of the very state to be represented (see \textsc{her self-model is part of the state}).
  \item No finite system can losslessly represent its entire current state within itself (see \textsc{complete self-representation is impossible}). In measurement-theoretic form, the physicist Thomas Breuer's self-measurement theorem (see \hyperref[source:source-breuer-1995]{\textsc{Thomas Breuer — The Impossibility of Accurate State Self-Measurements (Philosophy of Science, 1995)}}) shows that no observer can distinguish all present states of a system in which they are contained — whether the system is classical or quantum, its evolution deterministic or stochastic.
\end{enumerate}

It follows that no possible physical Mary satisfies the stipulation. ``Mary knows all the physical facts'' describes no realizable situation, so the verdict about her release is drawn from an impossible case and the argument's evidential force is correspondingly void. At best it reasons from a \emph{counterpossible} — ``had the impossible obtained, she would still have learned …'' — whose intuitive evaluation is unreliable.

\webheading{Where it sits}

This is complementary to the misimagining reply (see \hyperref[entry:argument-mary-misimagined]{\textsc{Mary misimagined}}): that move says the verdict is generated by \emph{misimagining} the stipulated case, while this one says there is no coherent case to imagine in the first place — the surrogate was all there ever was. It is kin in spirit to the realizability attack on the giant lookup table (see \textsc{the lookup-table impossibility argument}). Should a proponent retreat from realizability to mere conceivability, the standing wedge that a coherent \emph{description} does not certify a possible \emph{case} (see \textsc{conceivability is not sufficient for physical realizability}) applies — the lesson of the coherent-but-unrealizable halting oracle (see \textsc{Halting oracle: coherent, unrealizable}). Note that the attack does not defend physicalism directly; it denies the thought experiment standing to refute it.

\webheading{The weakest premise}

Premise (1) is the soft spot, via the \emph{type-retreat}. The proponent can weaken the stipulation: Mary need not encode the world's (or her own) total \emph{token} state, only all \emph{general} physical truths — the laws plus the type-level facts about colour vision. What red is like, if it is any physical fact at all, is plausibly a type fact rather than a fact about Mary's momentary microstate. The self-measurement bound forbids complete token self-description; it is silent on type-omniscience, so the reductio passes the retreated premise by. A second escape targets premise (2)'s reading of ``knows'': Mary's knowledge could be \emph{dispositional} and externally stored — a library she consults sequentially rather than a simultaneous internal encoding — and the bound constrains only \emph{simultaneous, internal, lossless} self-representation. The attack therefore bites hardest on the maximal, literal reading of ``knows all the physical facts,'' which is also the reading the knowledge argument's rhetoric leans on. What survives the retreat is a weaker premise whose sufficiency for the famous intuition is no longer obvious — at which point the verdict question reopens, exactly where the misimagining reply (see \hyperref[entry:argument-mary-misimagined]{\textsc{Mary misimagined}}) presses.

\begin{plainterms}
The Mary story (see \hyperref[entry:argument-knowledge-argument]{\textsc{Knowledge argument (Mary's room)}}) starts: ``Mary knows every physical fact.'' This objection says that opening line is like ``consider a map so detailed it shows everything — including itself, at full detail'': it sounds fine but describes something impossible. Mary's brain is part of the physical world, so ``every physical fact'' includes the full present state of her own brain — and a part cannot hold a complete copy of the whole that contains it. This is a theorem, not a hunch: the physicist Thomas Breuer proved that no measuring device can pin down the complete state of a system it sits inside (see \hyperref[source:source-breuer-1995]{\textsc{Thomas Breuer — The Impossibility of Accurate State Self-Measurements (Philosophy of Science, 1995)}}). If the story's premise cannot be true in any possible world, our gut feeling about what would happen ``if it were'' is testimony about nothing. The escape for the story's defenders: maybe Mary only needs the general laws and facts about colour vision \emph{in general}, not a snapshot of her own atoms — though then the premise is weaker than the dramatic ``she knew \emph{everything},'' and it is less clear the famous intuition still follows. A related objection, that we misimagine what complete knowledge would be (see \hyperref[entry:argument-mary-misimagined]{\textsc{Mary misimagined}}), takes over from there.
\end{plainterms}

\webheading{See also}
\begin{seealso}
  \item \hyperref[entry:claim-mary-knows-all-physical]{\textsc{Mary knows all physical facts}} — the scenario premise this reply declares unrealizable.
  \item \hyperref[entry:concept-introspective-gap-varieties]{\textsc{Varieties of introspective explanatory gap}} — the family of limits on self-knowledge this attack draws on to bound an embedded knower.
  \item \hyperref[entry:concept-knowledge-argument]{\textsc{Knowledge argument}} — the thought experiment whose standing to refute physicalism is denied.
\end{seealso}

\entryhead{claim-micro-subjects-do-not-sum}{Micro-subjects do not sum}{Claim \,\textperiodcentered\, after James}{Micro-subjects do not sum}{Distinct micro-subjects cannot sum into a further macro-subject}{$\odot$}
A plurality of distinct experiencing subjects — each with its own point of view and its own self-contained phenomenal field — does not, merely by being arranged or related, constitute a further single subject whose experience contains theirs. As William James put it, a hundred feelings, however shuffled together, do not make a hundred-and-first feeling that includes them; each is its own whole. Sam Coleman sharpens the point: the would-be composite subject's experience is not entailed by the experiences of its parts, and would in fact \emph{exclude} them, since each micro-subject's perspective is private to it.

This is the engine of the \hyperref[entry:argument-combination-problem]{\textsc{combination problem}} facing constitutive panpsychism. If micro-subjects cannot sum, then macro-minds cannot be \emph{constituted} by the experiential properties of their micro-parts, and the central constitutive claim (see \hyperref[entry:claim-consciousness-fundamental-ubiquitous]{\textsc{Fundamental ubiquitous consciousness}}) fails as stated. The claim is denied by constitutive panpsychists, who seek an account of \emph{phenomenal bonding} or fusion on which micro-subjects do compose a macro-subject, and by panqualityists, who avoid the difficulty by denying that the micro-level involves \emph{subjects} at all — only unexperienced qualities. It is the mirror image of the \emph{decombination} problem facing \hyperref[entry:position-cosmopsychism]{\textsc{Cosmopsychism}}, where the puzzle is instead how one cosmic subject divides into many.

\webheading{See also}
\begin{seealso}
  \item \hyperref[entry:argument-combination-problem]{\textsc{Combination problem}} — the objection this claim powers.
  \item \hyperref[entry:claim-consciousness-fundamental-ubiquitous]{\textsc{Fundamental ubiquitous consciousness}} — the constitutive thesis that fails if subjects cannot sum.
  \item \hyperref[entry:position-cosmopsychism]{\textsc{Cosmopsychism}} — the top-down view facing the reverse, \emph{decombination}, problem.
\end{seealso}

\entryhead{argument-modal-rationalism}{Modal rationalism}{Argument \,\textperiodcentered\, after Chalmers}{Modal rationalism}{Modal rationalism — ideal conceivability is a guide to possibility}{$\vdash$}
Modal rationalism is the thesis that ideal rational reflection is a reliable guide to what is metaphysically possible: what an idealized reasoner can conceive without hidden contradiction marks out the space of genuine possibility. David Chalmers (\hyperref[source:source-chalmers-1996]{\textsc{Chalmers — The Conscious Mind (1996)}}; refined in \hyperref[source:source-chalmers-2002]{\textsc{Chalmers — On Sense and Intension (2002)}}) defends it not as a brute premise but by inference to the best explanation in modal epistemology — by disarming the only serious apparent counterexamples, the Kripkean a-posteriori necessities.

\begin{enumerate}
  \item (see \hyperref[entry:claim-aposteriori-necessities-are-misdescribed-possibilities]{\textsc{A-posteriori necessities as misdescribed possibilities}}) Every Kripkean a-posteriori necessity (see \hyperref[source:source-kripke-1980]{\textsc{Kripke — Naming and Necessity (1980)}}) is a \emph{genuinely possible} scenario \emph{misdescribed}: under a term's \emph{primary} (epistemic) \hyperref[entry:concept-primary-intension]{\textsc{Primary intension}} the conceived case is possible, while the necessity attaches only to its \emph{secondary} \hyperref[entry:concept-secondary-intension]{\textsc{Secondary intension}}.
  \item These a-posteriori cases are the \emph{only} serious candidate counterexamples to ``conceivability entails possibility.'' (The burden is to exhibit a defeater that resists this diagnosis.)
\end{enumerate}

∴ (see \hyperref[entry:claim-conceivability-entails-possibility]{\textsc{Conceivability entails possibility}}) Ideal \emph{primary} conceivability entails \emph{primary} possibility: once the two intensions are distinguished, the apparent counterexamples evaporate and ideal rational conceiving is a reliable guide to the space of possibility. This \texttt{responds\_to} \hyperref[entry:argument-conceivability-not-guide-to-possibility]{\textsc{A-posteriori-necessity objection}} by reinterpreting its premise rather than denying it.

\webheading{Bearing on the zombie}

The defence is decisive for \hyperref[entry:argument-zombie-conceivability]{\textsc{Zombie argument}} only if the phenomenal case has \emph{no} primary/secondary gap to misdescribe — that is, only if for phenomenal terms the primary and secondary intensions coincide, so that ``what it is like'' carries no appearance/reality distinction. That further claim is exactly where the type-B physicalist (see \hyperref[entry:concept-materialism-types]{\textsc{Type-A / Type-B / Type-C materialism}}) re-engages.

The \textbf{weakest premise} is really two. Premise (2) — that \emph{no} a-posteriori necessity resists the misdescription diagnosis (a ``strong'' necessity would be one that does). And the application to consciousness — that phenomenal terms lack the appearance/reality wedge that lets the water case off the hook. Deny either and conceivability stops being a safe guide for the zombie specifically, even if it remains safe in general.

\begin{plainterms}
This is the defence of the rule the zombie argument needs: ``coherently imaginable ⟹ genuinely possible.'' The worry was ``water is H₂O'' — imaginable-otherwise yet impossible. The reply: when you ``imagine water that isn't H₂O,'' you really are imagining a genuine possibility — some watery stuff that isn't H₂O — you are just wrongly \emph{calling} it water. Separate ``the stuff that plays the water role'' from ``H₂O,'' and the imagined world is perfectly possible; the necessity was hiding in the labels. So imagination stays a reliable guide to what's possible, once you are careful about descriptions. The catch that keeps the consciousness debate alive: this rescue works only if ``feels like red'' has no similar gap between how it seems and what it is — and that is exactly what is disputed.
\end{plainterms}

\webheading{See also}
\begin{seealso}
  \item \hyperref[entry:claim-aposteriori-necessities-are-misdescribed-possibilities]{\textsc{A-posteriori necessities as misdescribed possibilities}} — the premise: the Kripke cases are possibilities under a misleading description.
  \item \hyperref[entry:claim-conceivability-entails-possibility]{\textsc{Conceivability entails possibility}} — the bridge this argument establishes.
  \item \hyperref[entry:argument-conceivability-not-guide-to-possibility]{\textsc{A-posteriori-necessity objection}} — the objection it answers.
  \item \hyperref[entry:concept-primary-intension]{\textsc{Primary intension}} / \hyperref[entry:concept-secondary-intension]{\textsc{Secondary intension}} — the two intensions whose distinction does the work.
  \item \hyperref[entry:argument-zombie-conceivability]{\textsc{Zombie argument}} — the argument this defence is ultimately meant to secure.
  \item \hyperref[entry:concept-materialism-types]{\textsc{Type-A / Type-B / Type-C materialism}} — type-B physicalism, where the defence is resisted.
\end{seealso}

\end{multicols}
\bookgroup{N}
\begin{multicols}{2}
\entryhead{claim-no-radical-emergence}{No radical emergence}{Claim \,\textperiodcentered\, after Strawson}{No radical emergence}{Experience cannot emerge brutely from the wholly non-experiential}{$\odot$}
There can be no \emph{radical} (brute) emergence of experiential properties from a base that is wholly and utterly non-experiential. Where genuine emergence occurs in nature, the emergent feature is intelligibly grounded in the powers of its base — liquidity in the dynamics of water molecules, say — so that, given the base, the new feature is no surprise. To get \emph{experience}, the felt what-it-is-likeness, from the entirely experienceless would be emergence with no such intelligible link; and that, Strawson argues, is no better than magic. The claim does not deny emergence as such — only brute emergence across the experiential / non-experiential divide.

This is the load-bearing premise of the \hyperref[entry:argument-genetic-panpsychism]{\textsc{genetic argument}} for panpsychism: if experience is real and wholly natural, and cannot brutely emerge, then it must already be present at the fundamental level. It is denied by physicalists who accept strong emergence (see \hyperref[entry:claim-strong-emergence-is-acceptable]{\textsc{Acceptability of strong emergence}}), holding that experience can arise as a genuinely novel feature without being present in the parts — the standard reply to the genetic argument.

\webheading{See also}
\begin{seealso}
  \item \hyperref[entry:argument-genetic-panpsychism]{\textsc{Genetic argument}} — the argument this premise anchors.
  \item \hyperref[entry:claim-strong-emergence-is-acceptable]{\textsc{Acceptability of strong emergence}} — its direct denial, and the engine of the emergence reply.
  \item \hyperref[entry:claim-phenomenal-consciousness-is-real]{\textsc{Reality of phenomenal consciousness}} — the realism premise the genetic argument pairs with this one.
  \item \hyperref[entry:claim-no-nonarbitrary-threshold-for-experience]{\textsc{Absence of an experiential threshold}} — a distinct premise that targets \emph{where} experience would switch on, not \emph{whether} it could brutely appear at all.
\end{seealso}

\entryhead{claim-some-truths-about-experience-not-physically-deducible}{Non-deducible experiential truths}{Claim \,\textperiodcentered\, after Jackson}{Non-deducible experiential truths}{There are truths about conscious experience not deducible from the complete physical description}{$\odot$}
There are truths about conscious experience that cannot be deduced from the complete physical description of the world. This is the conclusion of the \hyperref[entry:concept-knowledge-argument]{\textsc{Knowledge argument}} in its \emph{epistemic} form. Frank Jackson's Mary knows every physical fact about colour vision yet, on first seeing red, appears to learn something new — what seeing red is \emph{like}. If she does, that truth was not contained in the physical facts she already had: a complete physics leaves some facts about experience out.

The claim is deliberately narrow. It is \emph{epistemic} — about what can be deduced or known — and is the immediate output of the Mary case, nothing more. The further step to the \emph{ontological} thesis that the missing truths concern \emph{non-physical} facts (see \hyperref[entry:claim-consciousness-fundamental-nonphysical]{\textsc{Fundamental non-physical consciousness}}) is a separate move. Physicalists can grant an epistemic gap — Mary gains a new concept or a new way of grasping an old fact — while denying that any non-physical \emph{fact} exists, so the epistemic conclusion does not by itself settle the metaphysics.

The natural way to resist the claim is upstream, at \hyperref[entry:claim-mary-learns-on-release]{\textsc{Mary learns on release}}: on the ability and old-fact/new-guise replies, Mary acquires no new \emph{truth} at all on release, and if she learns no truth, nothing follows about what is or is not deducible.

It should not be confused with two neighbours. \hyperref[entry:claim-phenomenal-residue]{\textsc{Phenomenal residue}} says functional or causal role does not \emph{exhaust} the phenomenal; \hyperref[entry:claim-consciousness-fundamental-nonphysical]{\textsc{Fundamental non-physical consciousness}} is the ontological non-physicality thesis. This claim is the narrower bridge between them: that the phenomenal is not \emph{deducible} from the physical — the step that carries the Mary intuition toward, but not all the way to, the metaphysical conclusion.

\begin{plainterms}
The lesson drawn from Mary: if someone could know \emph{every} physical fact and still learn something new on first seeing colour, then that something isn't contained in the physical facts — there are truths about what experience is \emph{like} that you can't derive from even a complete physics. Note this is a claim about what's \emph{deducible} or \emph{knowable}; getting from there to ``and therefore those truths are non-physical'' (see \hyperref[entry:claim-consciousness-fundamental-nonphysical]{\textsc{Fundamental non-physical consciousness}}) is an extra, separately disputed step.
\end{plainterms}

\webheading{See also}
\begin{seealso}
  \item \hyperref[entry:concept-knowledge-argument]{\textsc{Knowledge argument}} — the thought experiment this claim is the epistemic conclusion of; Mary the colour scientist and what she does or doesn't gain on release.
  \item \hyperref[entry:claim-mary-learns-on-release]{\textsc{Mary learns on release}} — the premise upstream of this one; deny that Mary learns any new truth and this conclusion never gets going.
  \item \hyperref[entry:claim-consciousness-fundamental-nonphysical]{\textsc{Fundamental non-physical consciousness}} — the ontological thesis this claim stops short of; the contested further step from an epistemic gap to a non-physical fact.
  \item \hyperref[entry:claim-phenomenal-residue]{\textsc{Phenomenal residue}} — the distinct claim that role does not \emph{exhaust} the phenomenal, not to be conflated with the present claim about deducibility.
  \item \hyperref[entry:concept-hard-problem]{\textsc{Hard problem}} — why physical description seems to leave felt experience unexplained; the broader problem this deducibility gap is one sharp instance of.
  \item \hyperref[entry:question-what-does-mary-learn]{\textsc{What Mary learns}} — the open question this claim answers: what, if anything, Mary acquires on first seeing red.
\end{seealso}

\end{multicols}
\bookgroup{O}
\begin{multicols}{2}
\entryhead{argument-old-fact-new-guise}{Old fact, new guise}{Argument \,\textperiodcentered\, after Loar}{Old fact, new guise}{Old fact, new guise — Mary's new knowledge is a new phenomenal concept, not a new fact}{$\vdash$}
The \emph{old-fact / new-guise} argument is the physicalist's reply that grants the knowledge argument both its premises yet blocks its conclusion. Even if Mary knew every physical fact and genuinely learns something new on seeing red, that new knowledge can be a new \emph{concept} of a fact she already knew rather than a new fact — so it gives the anti-physicalist no extra ingredient in the world. The move is the Mary-specific application of the \emph{phenomenal-concept strategy}, due to Brian Loar (see \hyperref[source:source-loar-1990]{\textsc{Loar — Phenomenal States (1990)}}).

It raises the bridge critical question against the \hyperref[entry:argument-knowledge-argument]{\textsc{knowledge argument}}: even granting that Mary knew all the physical facts and learns something genuinely new, the bridge from ``new knowledge'' to ``new, non-physical fact'' fails, because knowledge is individuated more finely than the facts known.

\webheading{The inference}

The form is deductive, turning on an independently motivated distinction.

\begin{enumerate}
  \item One fact can be known under distinct concepts, and coming to grasp it under a second concept is genuine epistemic progress without any new fact — the lesson of \emph{a posteriori} identities (see \textsc{Hesperus/Phosphorus, water/H₂O}).
  \item Experiences give rise to \emph{phenomenal concepts}: first-person, recognitional concepts of inner states, cognitively isolated from the theoretical concepts of those same states (see \hyperref[source:source-loar-1990]{\textsc{Loar 1990}}).
  \item Inside the room Mary grasped the fact of seeing-red under theoretical concepts only; seeing red equips her with the phenomenal concept of that same state.
\end{enumerate}

Therefore \hyperref[entry:claim-mary-gains-old-fact-new-guise]{\textsc{her new knowledge is a new guise of an old fact}}, and \hyperref[entry:claim-some-truths-about-experience-not-physically-deducible]{\textsc{the non-deducibility of some truths about experience}} does not follow in any sense that troubles physicalism: what is not deducible is a \emph{way of thinking}, not a \emph{truth about the world}.

\webheading{Relation to the general strategy}

This is the Mary-specific case of the phenomenal-concept strategy whose general form — aimed at the explanatory gap rather than the knowledge argument — is \textsc{the gap-deflation argument} (Prong A): the same concept-isolation machinery that predicts a felt explanatory gap predicts that Mary's book-learning could never substitute for the phenomenal concept. It accordingly inherits that strategy's standing attack. \textsc{The overgeneration objection} presses that a concept isolated enough to do this work is too isolated to ground phenomenal self-knowledge (see \textsc{Phenomenal-concept trade-off}) — a cost this argument must pay too.

\webheading{Weakest premise}

Premise (2). If the phenomenal concept presents its referent via a substantive phenomenal property — a ``what-it-is-likeness'' the theoretical concepts miss — the anti-physicalist relocates the argument onto that property and the residue returns. If instead it presents its referent blankly, as a bare demonstrative, it is unclear why acquiring it should feel like the revelation the Mary intuition reports. Premise (1) is also pressured: in the Hesperus case the two guises connect to distinct \emph{causal routes} to one object, and a complete physical description arguably lets one deduce both routes — whereas Mary's case is precisely one where deduction fails, so the analogy may assume what it must prove.

\begin{plainterms}
The Mary story (see \hyperref[entry:argument-knowledge-argument]{\textsc{the knowledge argument}}) concludes that physics left a fact out, because book-perfect Mary still learns something at her first sight of red. This reply says: learning something new doesn't always mean learning a new \emph{fact}. Discovering that the Morning Star is the Evening Star felt like news, yet it added no new planet — just a second label connected to the same thing. Mary's first experience of red installs the inside-view label for a brain fact she already had the science-view label for. Real ``aha!'', same world. The standing objection (made elsewhere in this web against the same general strategy — see \textsc{the overgeneration objection}): if the inside-view label is so walled off from the science-view one, it's hard to explain how it tells you anything about yourself at all.
\end{plainterms}

\webheading{See also}
\begin{seealso}
  \item \hyperref[entry:argument-indexical-reply]{\textsc{Indexical reply}} — the sibling deflationary reply; the collapse objection charges that it re-enacts this very argument with an ``indexical'' relabelling.
\end{seealso}

\end{multicols}
\bookgroup{P}
\begin{multicols}{2}
\entryhead{concept-panpsychism}{Panpsychism}{Concept}{Panpsychism}{}{$\blacklozenge$}
\emph{Panpsychism} is the view that consciousness is a fundamental and ubiquitous feature of the physical world — that mind, or something mind-like, is present at the most basic level of nature rather than appearing only in brains. In its most discussed modern form it does \textbf{not} claim that rocks ponder or that electrons have thoughts; it claims that the simplest physical constituents carry some rudimentary form of experience, and that the rich consciousness of a human or an animal is built up from these elementary experiential ingredients. The word, from the Greek \emph{pan} (all) and \emph{psychē} (soul or mind), is old — versions appear among the Presocratics, in Spinoza, in Leibniz, and in William James — but the view has had a vigorous revival in analytic philosophy of mind as a proposed escape from the \hyperref[entry:concept-hard-problem]{\textsc{hard problem of consciousness}}.

This entry defines the doctrine and surveys its varieties. The arguments for and against it are developed as their own nodes and linked below; the committed \emph{position} that adopts panpsychism as its answer to the hard problem — what it commits to and what it denies — is recorded at \hyperref[entry:position-panpsychism]{\textsc{Panpsychism (position)}}.

\webheading{What the view reacts to}

Panpsychism is best understood as a third path between the two dominant options. \emph{Physicalism} holds that consciousness is nothing over and above physical processing, and faces the objection that no account of mere processing seems to explain why there is felt experience at all. \emph{Dualism} adds consciousness as a further, non-physical ingredient, and faces the objection that it leaves mind and matter mysteriously disconnected. Panpsychism agrees with the dualist that experience is fundamental — not reducible to non-experiential facts — but agrees with the physicalist that it belongs to the natural, physical world rather than arriving from outside it. Experience, on this view, was in the ingredients from the start; the brain does not conjure it at some threshold of complexity, it organizes what was already there.

\webheading{Varieties}

Panpsychism is a family, not a single thesis, and several distinctions sort its members. The taxonomy below largely follows David Chalmers \hyperref[source:source-chalmers-2017]{\textsc{Chalmers — Panpsychism and Panprotopsychism (2017)}}.

\begin{itemize}
  \item \textbf{Micropsychism vs. cosmopsychism} — whether the fundamental conscious things are the \emph{smallest} parts, with big minds composed bottom-up, or the \emph{whole} cosmos, with individual minds carved out of it top-down. The cosmic version is \hyperref[entry:position-cosmopsychism]{\textsc{Cosmopsychism}}.
  \item \textbf{Constitutive vs. emergentist} — whether macro-experience is \emph{constituted by} (grounded in, made up of) micro-experience, or instead \emph{strongly emerges} from it as a novel further fact. Constitutive forms are the ones that promise to avoid brute emergence; emergentist forms keep an emergence step.
  \item \textbf{Panpsychism vs. panprotopsychism} — whether the fundamental properties are themselves \emph{experiential}, or merely \emph{proto-experiential}: not yet consciousness, but the special non-structural ingredients from which consciousness can be constituted. Panprotopsychism is the more cautious cousin, and shades into \hyperref[entry:position-russellian-monism]{\textsc{Russellian monism}}.
  \item \textbf{Panexperientialism vs. panqualityism} — whether what is ubiquitous is full \emph{subjective experience} (a point of view, a felt field) or merely unexperienced \emph{qualities}, with subjecthood added later.
\end{itemize}

Constitutive Russellian micropsychism — fundamental physical entities have experiential intrinsic natures, and our minds are constituted from them — is the version most defended today, by Galen Strawson and Philip Goff (technically in \hyperref[source:source-goff-2017]{\textsc{Goff — Consciousness and Fundamental Reality (2017)}}, for a general readership in \hyperref[source:source-goff-2019]{\textsc{Goff — Galileo's Error: Foundations for a New Science of Consciousness (2019)}}) among others.

\webheading{The case for it}

Three lines of argument recur, each developed in its own node.

\begin{itemize}
  \item The \textbf{\hyperref[entry:argument-genetic-panpsychism]{\textsc{genetic argument}}} (Strawson): a physicalist who is also a realist about experience must place experience at the fundamental level, on pain of accepting that it emerges brutely — unintelligibly — from the wholly non-experiential.
  \item The \textbf{\hyperref[entry:argument-no-arbitrary-line-panpsychism]{\textsc{continuity argument}}} (Nagel): there is no non-arbitrary point at which experience could ``switch on,'' so the line-free option is that it reaches all the way down.
  \item The \textbf{\textsc{structure-and-dynamics argument}} (the Russellian route): physics describes only structure and dynamics, leaving the intrinsic categorical nature of matter blank — and consciousness is the natural candidate to fill it.
\end{itemize}

\webheading{Objections}

\begin{itemize}
  \item The \textbf{\hyperref[entry:argument-combination-problem]{\textsc{combination problem}}} (James, Coleman) is the standing difficulty: if the smallest things each have their own experience and point of view, how do countless such micro-subjects compose the \emph{single, unified} experience you are having now? Many regard this as panpsychism's make-or-break test, and note that it re-imposes, at the micro-to-macro step, the very explanatory burden the view was meant to discharge.
  \item The \textbf{\hyperref[entry:argument-strong-emergence-acceptable]{\textsc{emergence reply}}} presses the genetic argument's weak point: if strong emergence of experience is acceptable, there is no need to make mind fundamental, and the argument's central dilemma dissolves.
  \item Two looser worries are often added. The \emph{incredulous stare}: crediting electrons with any inner life is an extravagant cost in plausibility. And \emph{no empirical traction}: attributing experience to fundamental physics makes, by itself, no distinctive prediction. Defenders answer that every account of consciousness is strange somewhere, and that the hard problem is precisely the part of mind no empirical theory addresses head-on.
\end{itemize}

\begin{plainterms}
Panpsychism says consciousness goes all the way down. Not that rocks daydream — the serious version says the tiniest bits of matter carry the faintest spark of experience, and a brain assembles those sparks into a full mind. The appeal: you no longer have to explain how feeling gets conjured out of utterly feelingless stuff at some magic moment, because it was in the building blocks all along. The big catch, called the \emph{combination problem}: if every speck has its own tiny experience, how do trillions of them merge into the one seamless experience you're having as you read this? Critics also worry it just swaps one mystery for another, and that crediting electrons with any inner life is a stretch. Supporters answer that every theory of consciousness is strange somewhere, and at least this one keeps mind a natural part of the world.
\end{plainterms}

\webheading{See also}
\begin{seealso}
  \item \hyperref[entry:position-panpsychism]{\textsc{Panpsychism (position)}} — the committed position that adopts panpsychism as its answer to the hard problem, and how it differs from neighbouring views.
  \item \hyperref[entry:argument-genetic-panpsychism]{\textsc{Genetic argument}} — the headline argument \emph{for} the view: real experience plus no brute emergence entails fundamental mind.
  \item \hyperref[entry:argument-combination-problem]{\textsc{Combination problem}} — the central objection: micro-subjects cannot sum into one macro-subject.
  \item \hyperref[entry:concept-hard-problem]{\textsc{Hard problem}} — the problem panpsychism is offered to solve, by placing experience at the fundamental level rather than deriving it.
  \item \hyperref[entry:question-hard-problem]{\textsc{Hard problem (the question)}} — the open question panpsychism answers.
  \item \hyperref[entry:position-russellian-monism]{\textsc{Russellian monism}} — the closely allied view that consciousness is the intrinsic categorical nature physics leaves unspecified; constitutive panpsychism is one form of it.
  \item \hyperref[entry:position-cosmopsychism]{\textsc{Cosmopsychism}} — the top-down cousin, on which the cosmos is the one fundamental subject and individual minds are carved out of it.
  \item \hyperref[entry:claim-consciousness-fundamental-ubiquitous]{\textsc{Fundamental ubiquitous consciousness}} — the core claim panpsychism asserts, distinguished from the dualist's ``fundamental but non-physical.''
\end{seealso}

\entryhead{position-panpsychism}{Panpsychism (position)}{Position}{Panpsychism (position)}{Panpsychism}{$\bigstar$}
\emph{Panpsychism}, as a committed position in the philosophy of mind, holds that consciousness is a fundamental and ubiquitous feature of the physical world, and adopts that thesis as its answer to the problem of how experience fits into nature. It claims that the rich consciousness of a human or animal is not conjured from wholly mindless matter at some threshold of complexity, but is \emph{constituted} by the experiential — or proto-experiential — properties already carried by that matter's simplest constituents. The view is set out as a doctrine, with its varieties and arguments, in the concept entry \hyperref[entry:concept-panpsychism]{\textsc{Panpsychism}}; this entry records what the position commits to and how it stands against its neighbours.

The position's core commitment is that \hyperref[entry:claim-consciousness-fundamental-ubiquitous]{\textsc{consciousness is fundamental and present in the simplest physical entities}}. Because macro-minds are built up from the experiential natures of their micro-parts, panpsychism needs neither the dualist's non-physical addition nor the physicalist's ``strong emergence of feeling from the feelingless.'' It answers the \hyperref[entry:question-hard-problem]{\textsc{hard problem of consciousness}} by denying that physics must \emph{produce} experience at all: experience is there at the base, and the question becomes how it is organized rather than how it is generated. Its modern champions include Galen Strawson (see \hyperref[source:source-strawson-2006]{\textsc{Strawson — Realistic Monism: Why Physicalism Entails Panpsychism (2006)}}) and Philip Goff (\hyperref[source:source-goff-2017]{\textsc{Goff — Consciousness and Fundamental Reality (2017)}}, for a general readership \hyperref[source:source-goff-2019]{\textsc{Goff — Galileo's Error: Foundations for a New Science of Consciousness (2019)}}).

\webheading{How it differs from neighbouring views}

Against \textsc{Dualism}, consciousness is fundamental but \emph{intrinsic to the physical}, not a non-physical extra disconnected from matter. Against physicalism and \textsc{Illusionism}, it is firmly realist about experience: there is genuine felt consciousness to distribute, not an illusion to explain away. Its relation to \hyperref[entry:position-russellian-monism]{\textsc{Russellian monism}} is closest of all — \emph{constitutive} panpsychism simply \emph{is} a Russellian view on which the intrinsic natures physics leaves blank turn out to be experiential, while the more cautious \emph{panprotopsychism} makes them merely proto-experiential.

The position's standing weak point is the \hyperref[entry:argument-combination-problem]{\textsc{combination problem}}: if each micro-part has its own experience and point of view, how do countless micro-subjects compose the single, unified subject having your experience now? Many regard this as panpsychism's make-or-break test.

\begin{plainterms}
Panpsychism says consciousness goes all the way down: not that rocks have thoughts, but that the basic bits of matter carry the tiniest glimmer of experience, and brains assemble those glimmers into full-blown minds. The appeal: you don't have to conjure feeling out of utterly feelingless stuff at some magic threshold — it was in the ingredients from the start. The catch, its \emph{combination problem}: how do countless tiny experiences add up to the single unified experience you're having right now?
\end{plainterms}

\webheading{See also}
\begin{seealso}
  \item \hyperref[entry:concept-panpsychism]{\textsc{Panpsychism}} — the doctrine itself, with its full taxonomy of varieties and the arguments for and against; this position is its committed form.
  \item \hyperref[entry:claim-consciousness-fundamental-ubiquitous]{\textsc{Fundamental ubiquitous consciousness}} — the core thesis the position commits to.
  \item \hyperref[entry:claim-phenomenal-residue]{\textsc{Phenomenal residue}} — the position's further commitment that felt states are not exhausted by their functional role.
  \item \hyperref[entry:question-hard-problem]{\textsc{Hard problem (the question)}} — the problem the position is offered to solve.
  \item \hyperref[entry:argument-genetic-panpsychism]{\textsc{Genetic argument}} — the headline argument for the view: real experience plus no brute emergence entails fundamental mind.
  \item \hyperref[entry:argument-combination-problem]{\textsc{Combination problem}} — the central objection the position must answer.
  \item \hyperref[entry:position-russellian-monism]{\textsc{Russellian monism}} — the allied view that consciousness is the intrinsic nature physics leaves unspecified; constitutive panpsychism is one form of it.
  \item \hyperref[entry:position-cosmopsychism]{\textsc{Cosmopsychism}} — the top-down cousin, on which the cosmos is the one fundamental subject and individual minds are carved out of it.
  \item \textsc{Dualism} — the rival on which consciousness is fundamental but non-physical.
  \item \textsc{Illusionism} — the rival that denies there is any fundamental experience to distribute.
\end{seealso}

\entryhead{concept-paradox-of-phenomenal-judgment}{Paradox of phenomenal judgment}{Concept \,\textperiodcentered\, after Chalmers}{Paradox of phenomenal judgment}{The paradox of phenomenal judgment}{$\blacklozenge$}
A \emph{phenomenal judgment} is a judgment about conscious experience — the belief, thought, or report that one is conscious, that there is ``something it is like'' to see red, or that consciousness resists physical explanation. The \textbf{paradox of phenomenal judgment} (David Chalmers, \emph{The Conscious Mind}, 1996, \hyperref[source:source-chalmers-1996]{\textsc{Chalmers — The Conscious Mind (1996)}}) is the tension that arises when such judgments are set beside the view that consciousness is non-physical. Our phenomenal judgments are themselves physical events — utterances, neural states, dispositions to act. By the \hyperref[entry:concept-causal-closure]{\textsc{causal closure of the physical}}, they have complete physical causes. But if consciousness is something over and above the physical, then consciousness itself is \emph{not} among those causes. We would make exactly the same judgments about our experience whether or not the experience accompanied the physical processes producing them.

The paradox is sharpest for property dualism and especially \textsc{Epiphenomenalism}, on which experience is real but \textsc{causally inert}: there, my saying ``I am conscious'' is caused entirely by my brain, and the consciousness it reports does no work in producing the report. The unsettling consequences are that my belief that I am conscious is not \emph{caused} by my consciousness, and is not \emph{justified} by it either — yet it seems both reliable and certain. A \hyperref[entry:concept-philosophical-zombie]{\textsc{zombie}} twin, physically identical but experienceless, would sincerely assert ``I am conscious'' for precisely the same reasons, and would be wrong. What, then, makes \emph{my} matching judgment knowledge?

\webheading{Why it generalizes}

Chalmers presses the paradox not only as an objection to dualism but as a puzzle for \emph{every} view, including his own. Any complete physical-functional account of why we produce consciousness-talk will, by closure, make no essential reference to consciousness itself — so the explanation of our phenomenal \emph{judgments} threatens to proceed without the phenomenal \emph{facts} those judgments are about. Pushed to its general form, this becomes the \textbf{meta-problem of consciousness} (see \hyperref[source:source-chalmers-2018]{\textsc{Chalmers — The Meta-Problem of Consciousness (2018)}}): the problem of explaining, in topic-neutral terms, why we judge and report that there is a hard problem at all. The worry is that a good answer would \emph{debunk} the very beliefs it explains — the engine driving the report-scepticism disputes elsewhere in this web (see \textsc{Report-scepticism over-reach}).

\begin{plainterms}
You are certain you're conscious, and you say so. But the words coming out of your mouth are produced by your brain — physical machinery that, on a ``consciousness is something extra'' view, would churn out the very same words even if the inner feeling weren't there. So the feeling isn't what makes you \emph{say} you're conscious, and arguably isn't what makes you \emph{know} it either. A perfect feelingless duplicate of you would insist, just as earnestly, that it has rich inner experience — and be flat wrong. The paradox: if your report and the zombie's have the same cause, what entitles you to trust yours? It is the hardest version of the question ``how do we even come to talk about consciousness?'', and it bites hardest on views that make consciousness non-physical or idle.
\end{plainterms}

\webheading{See also}
\begin{seealso}
  \item \textsc{Epiphenomenalism} and \textsc{Causally inert qualia} — the view the paradox embarrasses most directly: real but inert experience that cannot cause the reports about it.
  \item \hyperref[entry:concept-philosophical-zombie]{\textsc{Philosophical zombie}} — the duplicate who makes all the same phenomenal judgments with no experience behind them; the paradox turns the zombie intuition back on our own self-knowledge.
  \item \hyperref[entry:concept-hard-problem]{\textsc{Hard problem}} — the explanandum these judgments are about; the paradox is the puzzle of how the judgments relate to it.
  \item \hyperref[source:source-chalmers-2018]{\textsc{Chalmers — The Meta-Problem of Consciousness (2018)}} — the meta-problem, the paradox's modern, generalized descendant.
  \item \textsc{Report-scepticism over-reach} — the live dispute in this web that runs on the same debunking tension, applied to machine and human testimony alike.
\end{seealso}

\entryhead{concept-phenomenal-consciousness}{Phenomenal consciousness}{Concept}{Phenomenal consciousness}{}{$\blacklozenge$}
\emph{Phenomenal consciousness} is the felt, subjective side of a mental state: what it is \emph{like}, from the inside, to undergo it. A state is phenomenally conscious when there is something it is like to be in it, such as the redness of red as you see it, the sting of a burn, or the taste of coffee. The classic gloss is Thomas Nagel's question ``what is it like to be a bat?''; if there is something it is like to be the bat, the bat has phenomenal consciousness. The specific felt qualities are often called \emph{qualia}, and phenomenal consciousness is the general property of having any such felt character at all.

\webheading{What it is, by contrast with function}

The notion is isolated by setting it against everything a state \emph{does}. Two states might play exactly the same causal role, both caused by tissue damage and both driving you to wince and say ``ouch'', and one can still ask whether they \emph{feel} the same. That residue, the felt character left over once every functional role is accounted for, is phenomenal consciousness. This is what marks it off from \hyperref[entry:concept-access-consciousness]{\textsc{Access consciousness}}, the purely functional sense in which a state's content is merely \emph{available} for reasoning and report. Ned Block, who coined the labels \emph{P-consciousness} and \emph{A-consciousness}, argued that the two can dissociate in both directions, a separation set out at \hyperref[entry:concept-access-vs-phenomenal-consciousness]{\textsc{Access vs. phenomenal consciousness}}.

\webheading{Why it matters}

Phenomenal consciousness is the part of the mind that seems hardest to fit into a physical picture, and it is the explanandum of the \hyperref[entry:concept-hard-problem]{\textsc{Hard problem}}: the question of why physical processing should be accompanied by any felt experience at all. The classic anti-materialist thought experiments all turn on it. A zombie is a physical duplicate with no phenomenal consciousness, and Mary the colour scientist seems to learn a new phenomenal fact on first seeing red. What is \emph{disputed} is not usually the definition but the phenomenon's standing: illusionists hold that the impression of an extra felt residue is itself a cognitive illusion, while realists take phenomenal consciousness to be a genuine, further feature of the world.

\begin{plainterms}
Phenomenal consciousness is just \emph{feeling}: the fact that experiences have an inside to them. A brain scan can show every step of what your nervous system does when you stub your toe, yet it doesn't show the \emph{hurt} you actually feel, and that hurt is the phenomenal side. It is the thing that makes consciousness puzzling. We can explain what the body does, but explaining why any of it is \emph{felt} is the hard part. A few philosophers even argue the ``extra feeling'' is a trick the mind plays on itself.
\end{plainterms}

\webheading{See also}
\begin{seealso}
  \item \hyperref[entry:concept-access-vs-phenomenal-consciousness]{\textsc{Access vs. phenomenal consciousness}} — Block's distinction, and how phenomenal consciousness and its functional complement can come apart in both directions.
  \item \hyperref[entry:concept-access-consciousness]{\textsc{Access consciousness}} — the functional, availability sense it is contrasted with.
  \item \hyperref[entry:concept-hard-problem]{\textsc{Hard problem}} — the explanatory problem that phenomenal consciousness poses.
  \item \textsc{First-person access} — the \emph{epistemic} privileged access a subject has to its own phenomenal states, a different notion from the felt character itself.
\end{seealso}

\entryhead{claim-phenomenal-residue}{Phenomenal residue}{Claim \,\textperiodcentered\, after Chalmers}{Phenomenal residue}{Phenomenally conscious states are not exhausted by functional/causal role}{$\odot$}
For at least one kind of mental state — the \emph{phenomenally conscious} ones, the states there is ``something it is like'' to be in — a complete specification of what the state \emph{does} leaves something out. A pain has a functional/causal role (caused by tissue damage, causing you to wince and seek relief), but it also \emph{hurts}; and this claim holds that the felt, qualitative character is not fixed by the role. Role-individuation, the strategy of \textsc{Functionalism} that pins a mental state down by its causes and effects, is not the whole story for the phenomenal.

The claim is deliberately narrow. It is about phenomenal states \emph{only}, and it concedes \textsc{Functionalism} for intentional states such as belief and desire, where a description of the causal role may well capture the whole state. It does not assert that silicon systems lack consciousness, and it does not assert any biological requirement; it says only that functional role under-determines phenomenal character. Its bite against substrate-independence is local: the inference from ``individuated by role'' to ``had by any realiser of that role'' is sound only where role-individuation is exhaustive, and that is exactly what this claim denies in the phenomenal case. So it does not reject \textsc{substrate independence} wholesale — it denies the \emph{uniform} reading that extends it across every mental state. It is also distinct from \textsc{Proteocentrism}, which demands a protein/carbon substrate for consciousness: this claim makes no substrate demand at all. The view it most directly answers is the question of whether mental states are substrate-independent (see \textsc{Substrate independence}).

Those who would deny it are strong or role functionalists, who read functionalism uniformly across phenomenal as well as intentional states, and anyone holding that phenomenal character metaphysically supervenes on functional organisation.

\begin{plainterms}
There's a difference between what your mind \emph{does} and what your experience \emph{feels like}. A pain makes you wince and reach for relief — that's its job, the role it plays. But a pain also \emph{hurts}; there's something it is like to feel it. This claim says: for these felt, ``what-it's-like'' experiences, listing everything the state does still leaves something out — the feel itself. Two things it carefully does \emph{not} say: it doesn't say computers can't be conscious, and it doesn't say you need carbon or biology. It even grants that for plain thinking states — beliefs, plans — a description of the role might capture the whole story. It only fences off the \emph{felt} part of the mind and says: that part isn't fully pinned down by function alone.
\end{plainterms}

\webheading{See also}
\begin{seealso}
  \item \textsc{Substrate independence} — the issue this claim answers: whether mental states can be carried by any substrate that plays the right role.
  \item \textsc{substrate independence} — the thesis this claim restricts; it blocks the uniform ``all mental states'' reading by carving out the phenomenal.
  \item \textsc{Functionalism} — the role-individuation strategy this claim says runs out for phenomenal states.
  \item \textsc{Proteocentrism} — a stronger neighbour that demands a biological substrate; this claim makes no such demand.
\end{seealso}

\entryhead{concept-philosophical-zombie}{Philosophical zombie}{Concept \,\textperiodcentered\, after Chalmers}{Philosophical zombie}{}{$\blacklozenge$}
A \emph{philosophical zombie} is an imagined being that is microphysically and functionally identical to a conscious person — every atom, every brain state, every piece of behaviour the same — but with \emph{no phenomenal consciousness}: nothing it is like to be it, all dark inside. It is not the shambling creature of horror films; it is a same-physics twin that does everything you do, says ``ouch'' when pricked and talks about its feelings, yet has no inner experience whatsoever (see \hyperref[source:source-chalmers-1996]{\textsc{Chalmers — The Conscious Mind (1996)}}).

The point of the device is not to claim such beings exist, but to test whether the physical facts settle the facts about experience. If a complete physical duplicate of a conscious person could exist with experience subtracted, then experience is not just more physics — it is something \emph{over and above} the physical. That is exactly the conclusion the zombie is built to drive: it is the engine of the \emph{conceivability argument} against physicalism, the view that everything, consciousness included, is fixed by physical facts.

\webheading{The crucial step}

The argument turns on a move from imagining to possibility. From the fact that a zombie seems \emph{conceivable} — coherently and stably imaginable, with no contradiction surfacing on careful reflection — it infers that a zombie is genuinely \emph{possible}, and hence that \hyperref[entry:claim-consciousness-fundamental-nonphysical]{\textsc{consciousness is non-physical}}. Nearly the whole dispute concentrates on this bridge from conceiving to possibility (see \hyperref[entry:claim-conceivability-entails-possibility]{\textsc{Conceivability entails possibility}}): critics grant that we can picture a zombie while denying that the picture tracks any real possibility, much as one can imagine water that is not H₂O without water's being possibly anything but H₂O.

\webheading{Relation to neighbours}

The zombie is a limiting case of \emph{absent qualia} — experience entirely missing rather than merely altered — and a cousin of the \emph{inverted spectrum}, in which the felt qualities are swapped rather than absent. It sharpens the intuition behind the \hyperref[entry:concept-hard-problem]{\textsc{Hard problem}} of consciousness, the puzzle of why physical processing is accompanied by any felt experience at all. The related \textsc{explanatory-gap argument} already leans on zombie conceivability in its premises, and the \textsc{phenomenal-concept strategy} is the leading attempt to defuse the zombie by granting that it is conceivable while denying that it is possible.

\begin{plainterms}
A ``philosophical zombie'' is an imagined being that is an atom-for-atom copy of you — same brain, same behaviour, says ``ouch'' when pricked, chats about its feelings — except there is no one home: no inner experience at all, total darkness inside. The thought is not that such things exist; it is that \emph{if you can coherently imagine one}, then experience must be something extra that the physical facts alone don't pin down. Nearly the whole fight is over that ``if you can imagine it, it's genuinely possible'' step.
\end{plainterms}

\webheading{See also}
\begin{seealso}
  \item \hyperref[entry:concept-hard-problem]{\textsc{Hard problem}} — the puzzle of why physical processing is accompanied by felt experience at all; the zombie sharpens its central intuition.
  \item \hyperref[entry:claim-conceivability-entails-possibility]{\textsc{Conceivability entails possibility}} — the modal bridge from conceiving a zombie to its being possible, where the whole argument is contested.
  \item \hyperref[entry:claim-consciousness-fundamental-nonphysical]{\textsc{Fundamental non-physical consciousness}} — the conclusion the zombie is built to drive: experience is over and above the physical.
  \item \textsc{Explanatory-gap argument} — an evidential cousin whose premises already lean on zombie conceivability.
  \item \textsc{Phenomenal-concept deflation} — the phenomenal-concept strategy, which grants the zombie's conceivability but denies its possibility.
\end{seealso}

\entryhead{position-physicalism}{Physicalism}{Position}{Physicalism}{Physicalism (materialism)}{$\bigstar$}
\textbf{Physicalism} (its older name \textbf{materialism}) is the thesis that everything that concretely exists is physical, or is wholly fixed by what is physical. Applied to the mind, it holds that mental events are physical events (see \textsc{Mental events are physical}): there is no thought, sensation, or experience over and above the physical workings of the brain and body. A common precise formulation is a \emph{supervenience} thesis — any world physically identical to ours is identical to ours in every respect, mental included — so that once the physical facts are fixed, all the facts are fixed.

The modern positive case for physicalism is the \textbf{causal argument} (see \hyperref[entry:argument-causal-argument-for-physicalism]{\textsc{Causal argument for physicalism}}), which turns on two principles this view leans on heavily: the \hyperref[entry:concept-completeness-of-physics]{\textsc{Completeness of physics}} (every physical effect has a sufficient physical cause) and the \hyperref[entry:concept-causal-closure]{\textsc{Causal closure of the physical}} of the physical (no non-physical cause produces a physical effect). If the mind genuinely moves the body, and the physical world is causally closed, the mind must itself be physical — the alternative being to demote it to a causally idle by-product (see \textsc{Epiphenomenalism}). Physicalism thus presents itself less as a metaphysical preference than as the only account of mental causation consistent with the closed causal order science describes (Papineau, \hyperref[source:source-papineau-2002]{\textsc{Papineau — Thinking about Consciousness (2002)}}; Kim, \hyperref[source:source-kim-2005]{\textsc{Kim — Physicalism, or Something Near Enough (2005)}}).

\textbf{What distinguishes it from neighbouring views.} Physicalism is the \emph{genus} whose species populate much of this web: \textsc{Type-identity theory} (mental types \emph{are} physical types), \textsc{Eliminative materialism} (the mental as folk theory discharges to neuroscience), analytic functionalism, and \textsc{Illusionism} (phenomenal properties are introspective illusion) are all ways of being a physicalist. They divide over the \hyperref[entry:concept-hard-problem]{\textsc{Hard problem}}: Chalmers' \hyperref[entry:concept-materialism-types]{\textsc{Type-A / Type-B / Type-C}} taxonomy sorts them by whether they deny the epistemic gap between physical and phenomenal truths, accept it but deny an ontological gap, or hope to close it in the future. Against physicalism stand \textsc{Dualism} (which denies completeness or closure to keep the mind non-physical and efficacious), \textsc{Epiphenomenalism} (which keeps the mind non-physical at the cost of its causal powers), and the conceivability and knowledge arguments (see \hyperref[entry:argument-zombie-conceivability]{\textsc{Zombie argument}}, \hyperref[entry:argument-knowledge-argument]{\textsc{Knowledge argument (Mary's room)}}) that press an explanatory gap physicalism must absorb.

\begin{plainterms}
Physicalism — the same view people used to call materialism — says that, at bottom, everything is physical: there's only the stuff physics describes (and whatever is built out of it), with no separate soul-substance or mental extra. About the mind in particular, it says your thoughts and feelings just \emph{are} things going on in your brain. The strongest reason offered is about causes: the physical world seems to be a closed system where every physical event has a complete physical cause and nothing non-physical reaches in to push matter around. Since your mind clearly \emph{does} affect your body — you decide to stand, and you stand — the tidiest conclusion is that your mind is itself part of that physical machinery. The big challenge it has to answer is the felt, what-it's-like side of experience (the ``hard problem''): critics say you can describe every physical fact about a brain and still leave out what the experience \emph{feels like} from the inside.
\end{plainterms}

\webheading{See also}
\begin{seealso}
  \item \textsc{Mental events are physical} — the view's core commitment about the mind.
  \item \hyperref[entry:argument-causal-argument-for-physicalism]{\textsc{Causal argument for physicalism}} — its principal positive argument.
  \item \hyperref[entry:concept-completeness-of-physics]{\textsc{Completeness of physics}} and \hyperref[entry:concept-causal-closure]{\textsc{Causal closure of the physical}} — the two principles that argument rests on.
  \item \hyperref[entry:concept-materialism-types]{\textsc{Type-A / Type-B / Type-C materialism}} — Chalmers' taxonomy of physicalist responses to the hard problem.
  \item \textsc{Dualism}, \textsc{Epiphenomenalism} — the chief rivals, each rejecting a different premise.
  \item \hyperref[entry:question-hard-problem]{\textsc{Hard problem (the question)}} — the question physicalism must answer to be complete.
\end{seealso}

\entryhead{concept-primary-intension}{Primary intension}{Concept \,\textperiodcentered\, after Chalmers}{Primary intension}{}{$\blacklozenge$}
The \emph{primary intension} of a term is, roughly, the meaning it carries in virtue of how speakers fix its reference — what the term would pick out on the supposition that a given possible world turns out to be the actual one. It is one of the two ``intensions'' (functions from possible worlds to extensions) that David Chalmers's two-dimensional semantics assigns to an expression; the other is its \hyperref[entry:concept-secondary-intension]{\textsc{Secondary intension}}. Where the secondary intension asks ``given how the actual world in fact is, what does this term pick out in \emph{other} worlds?'', the primary intension asks the prior question: ``for each way the world might turn out to \emph{be} actually, what would the term pick out \emph{there}?''

The standard illustration is the word \emph{water}. Competent speakers fix the reference of ``water'' by a rough description — the clear, drinkable liquid that falls as rain and fills the lakes and rivers, the \emph{watery stuff}. The primary intension is the function encoding that condition: hand it a world considered as actual, and it returns whatever plays the watery-stuff role there. In our world that role is filled by H₂O, so the primary intension delivers H₂O. On Hilary Putnam's Twin Earth — a world just like ours except the watery stuff is a different compound, XYZ — \emph{considered as a candidate for being actual}, the primary intension of ``water'' delivers XYZ. The primary intension is in this sense sensitive to empirical discovery only hypothetically: it tells you what ``water'' would refer to under each hypothesis about which world is actual, and so its profile can be known largely a priori, by reflection on how the concept works rather than on which world we are in.

\webheading{Two dimensions of meaning}

The label ``two-dimensional'' comes from evaluating a sentence twice over against a world. A world can enter the calculation in two roles: as the world taken to be \emph{actual} (which settles what our terms refer to), and as a \emph{counterfactual} circumstance to be described using reference so settled. The primary intension occupies the first dimension — worlds considered as actual; the secondary intension occupies the second — worlds considered as counterfactual. For many ordinary terms the two coincide. They come apart precisely for the Kripkean \emph{a posteriori necessities} (\textsc{A-posteriori necessities}, from Kripke's \emph{Naming and Necessity} \hyperref[source:source-kripke-1980]{\textsc{Kripke — Naming and Necessity (1980)}}): ``water is H₂O'' is necessary along the secondary dimension (water is H₂O in every counterfactual world) yet contingent along the primary one (the watery stuff \emph{could have turned out} to be XYZ). The necessity that can be known only by experiment lives in the gap between the two intensions.

Chalmers's is the \emph{epistemic} version of two-dimensional semantics, developed in \emph{The Conscious Mind} (\hyperref[source:source-chalmers-1996]{\textsc{Chalmers — The Conscious Mind (1996)}}, where the ``primary/secondary'' terminology is introduced) and refined in \hyperref[source:source-chalmers-2002]{\textsc{Chalmers — On Sense and Intension (2002)}}; Frank Jackson reaches an equivalent apparatus from the side of conceptual analysis (see \hyperref[source:source-jackson-1998]{\textsc{Jackson — From Metaphysics to Ethics (1998)}}). The primary intension also travels under other names — \emph{epistemic intension}, \emph{A-intension}, \emph{1-intension} — and has cousins in the broader two-dimensional tradition (Kaplan's ``character,'' Stalnaker's ``diagonal''), though those were designed for different purposes and should not be assumed identical.

\webheading{Why the concept matters}

The primary intension is the load-bearing notion in the modern \emph{conceivability} arguments about the mind. Chalmers's \hyperref[entry:argument-modal-rationalism]{\textsc{Modal rationalism}} holds that what ideal reflection finds conceivable is a guide to what is \emph{primarily} possible — possible as a way the world might actually be — and the two-dimensional diagnosis defuses the Kripkean counterexamples by reinterpreting them as \hyperref[entry:claim-aposteriori-necessities-are-misdescribed-possibilities]{\textsc{A-posteriori necessities as misdescribed possibilities}}: when you ``conceive water that is not H₂O'' you conceive a genuinely possible world that verifies the \emph{primary} intension of ``water,'' and merely misdescribe it. This rescues the bridge from conceivability to possibility for the \hyperref[entry:concept-philosophical-zombie]{\textsc{Philosophical zombie}} — but only if the phenomenal case has no such primary/secondary gap to exploit, which is the separate and contested thesis that for consciousness the two intensions coincide (see \hyperref[entry:claim-phenomenal-intensions-coincide]{\textsc{Coincident phenomenal intensions}}).

\begin{plainterms}
Think of a word as having two jobs. One job is fixing \emph{what we're even talking about} — for ``water,'' something like ``the clear drinkable stuff in the rivers.'' The primary intension is that job written out as a rule: point it at any way the world might turn out to be, and it hands back whatever fits the description there. If the world turns out to be ours, ``water'' lands on H₂O; if it had turned out to be Twin Earth, the very same rule would have landed ``water'' on XYZ. So the primary intension is what you grasp just by \emph{understanding the word}, before you do any chemistry — which is why imagining things tends to track it. The other job, pinning the word to the actual stuff (H₂O) and keeping it there even when we talk about other worlds, belongs to the \hyperref[entry:concept-secondary-intension]{\textsc{Secondary intension}}.
\end{plainterms}

\webheading{See also}
\begin{seealso}
  \item \hyperref[entry:concept-secondary-intension]{\textsc{Secondary intension}} — the other half of the framework: reference fixed by the actual world, then held rigid across counterfactual worlds; the necessary a posteriori is the gap between the two.
  \item \textsc{A-posteriori necessities} — the Kripkean cases (water = H₂O) that pull the two intensions apart and that the primary intension is designed to handle.
  \item \hyperref[entry:claim-aposteriori-necessities-are-misdescribed-possibilities]{\textsc{A-posteriori necessities as misdescribed possibilities}} — the two-dimensional diagnosis: apparent counterexamples to ``conceivable ⟹ possible'' are possibilities under the primary intension, misdescribed.
  \item \hyperref[entry:argument-modal-rationalism]{\textsc{Modal rationalism}} — the argument that ideal primary conceivability is a reliable guide to primary possibility, which the primary intension makes precise.
  \item \hyperref[entry:claim-phenomenal-intensions-coincide]{\textsc{Coincident phenomenal intensions}} — the further claim that for phenomenal concepts the primary and secondary intensions collapse together, removing the wedge the water case relies on.
  \item \hyperref[entry:concept-philosophical-zombie]{\textsc{Philosophical zombie}} — the case whose conceivability the framework is meant to convert into genuine possibility.
\end{seealso}

\end{multicols}
\bookgroup{Q}
\begin{multicols}{2}
\entryhead{concept-qualia}{Qualia}{Concept}{Qualia}{}{$\blacklozenge$}
\emph{Qualia} (singular \emph{quale}, from the Latin for ``of what kind'') are the subjective, felt qualities of conscious experience: the redness of a ripe tomato as it looks to you, the bitterness of coffee, the particular ache of a stubbed toe. The term names not the tomato, the coffee, or the toe, but what undergoing those experiences is \emph{like} from the inside. The word was given its modern use by Clarence Irving Lewis (see \hyperref[source:source-ci-lewis-1929]{\textsc{Lewis 1929}}), who applied it to the recognisable qualitative characters of experience that recur from one occasion to the next.

\webheading{The contrast with function}

The standard way to isolate the idea is to set it against everything a mental state \emph{does}. Two states might play exactly the same causal role, both produced by tissue damage and both making you wince and say ``ouch'', and one can still ask whether they \emph{feel} the same. That residue, the felt character left over once every causal role is accounted for, is what ``qualia'' picks out. It is the same residue Thomas Nagel marks with the phrase \emph{what it is like}: an organism has conscious experience, he argues, if there is something it is like to be that organism, and no objective, third-person description seems to capture it (see \hyperref[source:source-nagel-1974]{\textsc{Nagel 1974}}). Two stock cases sharpen the notion. The \emph{inverted spectrum} asks whether your experience of red could be like my experience of green with all of our behaviour and physiology unchanged; if that is so much as coherent, the felt quality floats free of the functional facts. And Frank Jackson's colour scientist Mary, who knows every physical fact about colour vision yet seems to learn something new when she first sees red, dramatises the same gap (the \hyperref[entry:concept-knowledge-argument]{\textsc{Knowledge Argument}}, from \hyperref[source:source-jackson-1982]{\textsc{Jackson 1982}}).

\webheading{Senses of the term}

The word is contested, and the dispute is often about the definition rather than the facts. In a \emph{thin} sense, qualia are simply the ways experiences feel, with no further commitment, and almost no one denies that there are such felt ways. In a \emph{thick} sense, qualia are additionally taken to be \emph{intrinsic}, \emph{ineffable}, \emph{private}, and \emph{directly apprehended} by introspection. Daniel Dennett's ``Quining Qualia'' argues that nothing satisfies that demanding job description, so qualia in the thick sense should be abandoned (see \hyperref[source:source-dennett-1988]{\textsc{Dennett 1988}}). A claim that looks obvious on the thin reading can beg the question on the thick one, so an entry that leans on qualia should say which sense it means. The notion is also the close cousin of \hyperref[entry:concept-phenomenal-consciousness]{\textsc{phenomenal consciousness}}: qualia are the specific felt qualities, and phenomenal consciousness is the general property of having any such felt character at all.

\webheading{Why qualia matter}

Qualia mark the spot where the mind seems hardest to fit into a scientific picture of the world. They drive the \hyperref[entry:concept-hard-problem]{\textsc{hard problem}} of consciousness, the question of why physical processing should be accompanied by felt experience at all, and they power the classic anti-materialist thought experiments: Mary's room (the \hyperref[entry:concept-knowledge-argument]{\textsc{Knowledge Argument}}) and the \hyperref[entry:concept-philosophical-zombie]{\textsc{philosophical zombie}}, a creature physically identical to you with no inner feel. Each turns on the suspicion that the felt facts outrun the functional and physical ones, which is exactly what a \textsc{functionalist} account of mind must deny. How much weight qualia can bear, and whether the thick notion survives scrutiny, is in large part what the debate over materialism is about.

\begin{plainterms}
``Qualia'' is the philosopher's word for how things \emph{feel} to you: the way red looks, the way coffee tastes, the way a headache hurts. A brain scanner can show what your neurons are doing when you see red, but it does not show the redness \emph{you} experience; qualia are that inner, felt side of the story. Philosophers argue about whether these feels are something extra that science struggles to explain, or whether, as skeptics like Daniel Dennett say, the whole idea falls apart once you look at it closely.
\end{plainterms}

\webheading{See also}
\begin{seealso}
  \item \hyperref[entry:concept-phenomenal-consciousness]{\textsc{Phenomenal consciousness}} — the general property of having any felt character at all; qualia are its specific instances.
  \item \hyperref[entry:concept-hard-problem]{\textsc{Hard problem}} — why physical processing is accompanied by felt experience at all; the problem qualia pose in its sharpest form.
  \item \hyperref[entry:concept-knowledge-argument]{\textsc{Knowledge argument}} — Mary the colour scientist, and whether complete physical knowledge still leaves the qualia out.
  \item \hyperref[entry:concept-philosophical-zombie]{\textsc{Philosophical zombie}} — a physical duplicate with no qualia; the conceivability test for whether the physical facts fix the felt facts.
  \item \hyperref[entry:concept-access-vs-phenomenal-consciousness]{\textsc{Access vs. phenomenal consciousness}} — Block's distinction; qualia live on the phenomenal side of it.
  \item \textsc{Functionalism} — the theory of mind that qualia are most often said to outrun.
\end{seealso}

\end{multicols}
\bookgroup{R}
\begin{multicols}{2}
\entryhead{claim-phenomenal-consciousness-is-real}{Reality of phenomenal consciousness}{Claim}{Reality of phenomenal consciousness}{Phenomenal consciousness is genuinely real, not an illusion}{$\odot$}
There really is something it is like to undergo experience: felt redness, the ache of a pain, the taste of coffee are genuine features of reality, not a cognitive illusion and not a mere disposition to report. This is \emph{realism} about phenomenal consciousness — the explanandum is taken at face value rather than explained away. It is a premise shared across most non-illusionist views of mind, including dualism, panpsychism, Russellian monism, and type-B physicalism. In the \hyperref[entry:argument-genetic-panpsychism]{\textsc{genetic argument}} for panpsychism it pairs with \hyperref[entry:claim-no-radical-emergence]{\textsc{No radical emergence}}: because experience is real \emph{and} cannot brutely emerge, it must be fundamental. The claim is denied only by illusionists, who hold that phenomenal consciousness as standardly conceived does not exist (see \textsc{Phenomenality as illusion}).

It should not be confused with the stronger \emph{residue} claim that felt states are not exhausted by their functional role (see \hyperref[entry:claim-phenomenal-residue]{\textsc{Phenomenal residue}}): one can affirm that experience is real while still debating whether it reduces to function. This claim asserts only its reality, not its irreducibility.

\webheading{See also}
\begin{seealso}
  \item \textsc{Phenomenality as illusion} — its direct denial, held by illusionists.
  \item \hyperref[entry:claim-phenomenal-residue]{\textsc{Phenomenal residue}} — the stronger, separate claim that experience outruns functional role.
  \item \hyperref[entry:argument-genetic-panpsychism]{\textsc{Genetic argument}} — the argument in which this realism premise does its work.
\end{seealso}

\entryhead{position-russellian-monism}{Russellian monism}{Position \,\textperiodcentered\, after Russell}{Russellian monism}{}{$\bigstar$}
\emph{Russellian monism} holds that physics describes only the abstract \emph{structure and dynamics} of the world — how things are arranged and how they behave — while leaving wholly unspecified the intrinsic, categorical nature of whatever has that structure, and that consciousness is, or is grounded in, that intrinsic nature. On this view experience is neither a non-physical add-on nor reducible to bare structure: it is the hidden categorical filling inside the physical world's outline, the ``what it is like to be the stuff'' that physical description never reaches.

The position commits to two claims. First, that \hyperref[entry:claim-functional-description-structural-dynamical-only]{\textsc{a functional or physical description captures only structural-dynamical facts}}, never intrinsic nature. Second, that \hyperref[entry:claim-consciousness-categorical-not-dispositional]{\textsc{phenomenal consciousness is categorical, not merely dispositional}} — an occurrent fact that is actually there, the kind of categorical basis that dispositions standardly require. Put together: consciousness is the categorical \emph{ground} of the very dispositions physics charts. This is the constructive view that the \textsc{structure-and-dynamics argument} points toward, and the escape Chalmers prefers from the epiphenomenalism that dogs his naturalistic dualism. It answers the \hyperref[entry:question-hard-problem]{\textsc{hard problem of consciousness}} by locating experience at the categorical base physics leaves blank.

The view descends from Bertrand Russell's \emph{The Analysis of Matter} (see \hyperref[source:source-russell-1927]{\textsc{1927}}) and has been revived in recent decades by Stoljar, Strawson, and Chalmers.

\webheading{How it differs from neighbouring views}

Russellian monism comes in degrees of commitment about \emph{what} the intrinsic nature is. In its \emph{constitutive} form it just \emph{is} a kind of \hyperref[entry:position-panpsychism]{\textsc{Panpsychism (position)}} (the intrinsic natures are themselves experiential) or \emph{panprotopsychism} (they are proto-experiential — not yet consciousness, but the special non-structural ingredients from which it can be constituted). It contrasts with \textsc{Dualism}, for which consciousness is non-physical rather than the intrinsic-physical filling; and with \textsc{Mysterianism}, which holds the psychophysical link to be permanently beyond human grasp where Russellian monism instead \emph{proposes what that link is}.

\begin{plainterms}
Russellian monism starts from a quiet point about physics: physics only ever tells you what things \emph{do} — how they push, attract, evolve — never what they intrinsically \emph{are} underneath all that doing. The idea is that consciousness is that hidden inner nature, the ``what it's like to be the stuff'' physics leaves as a blank. Then mind isn't a ghostly extra, as in dualism, and isn't conjured from pure structure, as in plain physicalism; it's the very filling inside the physical world's outline. It's the middle road Chalmers leans toward, and it shades into panpsychism.
\end{plainterms}

\webheading{See also}
\begin{seealso}
  \item \hyperref[entry:claim-functional-description-structural-dynamical-only]{\textsc{Structure and dynamics only}} — the first commitment: physics gives only structure and dynamics, not intrinsic nature.
  \item \hyperref[entry:claim-consciousness-categorical-not-dispositional]{\textsc{Categorical, not dispositional}} — the second commitment: experience is an occurrent categorical fact, the right kind to be that intrinsic nature.
  \item \textsc{Structure-and-dynamics argument} — the argument the position is the constructive answer to.
  \item \hyperref[entry:question-hard-problem]{\textsc{Hard problem (the question)}} — the problem it addresses, by filling the categorical blank.
  \item \hyperref[entry:position-panpsychism]{\textsc{Panpsychism (position)}} — constitutive Russellian monism is a form of it, with experiential intrinsic natures.
  \item \hyperref[entry:concept-panpsychism]{\textsc{Panpsychism}} — the broader doctrine, whose taxonomy places panprotopsychism as the cautious cousin shading into this view.
  \item \textsc{Dualism} — the rival on which consciousness is non-physical, not intrinsic-physical.
  \item \textsc{Mysterianism} — the rival that says the psychophysical link cannot be grasped at all.
\end{seealso}

\end{multicols}
\bookgroup{S}
\begin{multicols}{2}
\entryhead{concept-secondary-intension}{Secondary intension}{Concept \,\textperiodcentered\, after Chalmers}{Secondary intension}{}{$\blacklozenge$}
The \emph{secondary intension} of a term is the meaning it carries once its reference has been fixed by the actual world — the function that, taking that reference as given, says what the term picks out across all \emph{counterfactual} worlds. It is one of the two ``intensions'' (functions from possible worlds to extensions) that David Chalmers's two-dimensional semantics assigns to an expression; the other, prior, one is its \hyperref[entry:concept-primary-intension]{\textsc{Primary intension}}. The secondary intension is what captures the Kripkean insight that many terms are \emph{rigid designators} (see \hyperref[source:source-kripke-1980]{\textsc{Kripke — Naming and Necessity (1980)}}): having latched onto a thing in the actual world, they keep referring to that same thing when we describe how matters might otherwise have been.

Return to \emph{water}. We discover empirically that the watery stuff of \emph{our} world is H₂O. The secondary intension then takes that discovery as settled and runs it across all worlds: it maps every counterfactual world to H₂O. So considered counterfactually, a world whose lakes hold XYZ rather than H₂O is \emph{not} a world where water is something other than H₂O — it is a world with no water in it, however watery the XYZ looks. This is why ``water is H₂O,'' though learned by experiment, comes out \emph{necessary}: the secondary intension of ``water'' is constant across the counterfactual worlds, picking out H₂O in each. The truth is necessary along this dimension and a posteriori along the other — exactly the profile of the Kripkean \textsc{A-posteriori necessities}.

\webheading{How it differs from the primary intension}

The two intensions answer two different questions about the same term. The primary intension ranges over worlds \emph{considered as actual} and asks what ``water'' would pick out if that world were the actual one — so it can deliver XYZ when Twin Earth is the candidate for actuality. The secondary intension holds the actual world's verdict fixed (water \emph{is} H₂O) and ranges over worlds \emph{considered as counterfactual}, delivering H₂O everywhere. For a term whose reference is fixed descriptively with no empirical residue, the two functions coincide; for a natural-kind term like ``water'' they diverge, and the divergence \emph{is} the necessary a posteriori. The two-dimensionalist's payoff is that the apparent counterexamples to ``conceivability entails possibility'' can then be redescribed as \hyperref[entry:claim-aposteriori-necessities-are-misdescribed-possibilities]{\textsc{A-posteriori necessities as misdescribed possibilities}}: what looks like conceiving an impossibility (water that isn't H₂O) is really conceiving a genuine possibility that satisfies water's \emph{primary} intension, with the \emph{secondary} necessity left untouched.

Chalmers's epistemic version of the framework introduces the ``primary/secondary'' pair in \emph{The Conscious Mind} (see \hyperref[source:source-chalmers-1996]{\textsc{Chalmers — The Conscious Mind (1996)}}) and refines it in \hyperref[source:source-chalmers-2002]{\textsc{Chalmers — On Sense and Intension (2002)}}; Frank Jackson develops a matching apparatus from conceptual analysis (see \hyperref[source:source-jackson-1998]{\textsc{Jackson — From Metaphysics to Ethics (1998)}}). The secondary intension is also called the \emph{subjunctive intension}, \emph{C-intension}, or \emph{2-intension}, and corresponds to the familiar Kripkean rigid-designator meaning — the intension that fixes a term's modal profile.

\webheading{Why the concept matters}

In the consciousness debate the secondary intension is the physicalist's foothold. Grant that the \hyperref[entry:concept-philosophical-zombie]{\textsc{Philosophical zombie}} is conceivable; the type-B physicalist (see \hyperref[entry:concept-materialism-types]{\textsc{Type-A / Type-B / Type-C materialism}}) replies that conceivability shows only \emph{primary} possibility, and that — as with water — a primary/secondary gap may make the zombie metaphysically impossible all the same, its conceivability a misdescription. Whether that escape is available turns on a single further question: does ``consciousness'' behave like ``water,'' with a wedge between how it presents and what it is, or do its two intensions coincide (see \hyperref[entry:claim-phenomenal-intensions-coincide]{\textsc{Coincident phenomenal intensions}})? If they coincide, there is no secondary necessity to hide the impossibility in, and primary conceivability carries all the way to genuine possibility.

\begin{plainterms}
Once we find out what the stuff in our rivers actually \emph{is} — H₂O — the word ``water'' gets pinned to H₂O and stays pinned, even when we talk about other possible worlds. The secondary intension is that pinned-down meaning: it says ``water'' means H₂O \emph{everywhere}, so a world with XYZ in its rivers just doesn't have any water in it. That's why ``water is H₂O'' turns out to be a truth that \emph{couldn't have been otherwise}, even though we had to do chemistry to learn it. The companion notion, what ``water'' means \emph{before} we know which world we're in, is the \hyperref[entry:concept-primary-intension]{\textsc{Primary intension}}. The whole point of keeping the two apart is to explain how a fact can be both discovered by experiment and yet necessary.
\end{plainterms}

\webheading{See also}
\begin{seealso}
  \item \hyperref[entry:concept-primary-intension]{\textsc{Primary intension}} — the prior, reference-fixing dimension; the secondary intension takes over once the actual world has settled what the term refers to.
  \item \textsc{A-posteriori necessities} — the necessary a posteriori (water = H₂O), which is exactly a constant secondary intension paired with a varying primary one.
  \item \hyperref[entry:claim-aposteriori-necessities-are-misdescribed-possibilities]{\textsc{A-posteriori necessities as misdescribed possibilities}} — the two-dimensional move that reads the Kripke cases as primary possibilities, preserving the conceivability–possibility link.
  \item \hyperref[entry:concept-materialism-types]{\textsc{Type-A / Type-B / Type-C materialism}} — type-B physicalism, which leans on a possible primary/secondary gap for ``consciousness'' to block the zombie argument.
  \item \hyperref[entry:claim-phenomenal-intensions-coincide]{\textsc{Coincident phenomenal intensions}} — the claim that no such gap exists for phenomenal concepts, so the secondary-intension escape is unavailable for consciousness.
  \item \hyperref[entry:concept-philosophical-zombie]{\textsc{Philosophical zombie}} — the case whose modal status hangs on whether phenomenal terms have a divergent secondary intension.
\end{seealso}

\entryhead{claim-functional-description-structural-dynamical-only}{Structure and dynamics only}{Claim \,\textperiodcentered\, after Chalmers}{Structure and dynamics only}{A functional description captures only structural-dynamical facts, not intrinsic nature}{$\odot$}
A purely functional or causal specification of a system captures only its \emph{structure} — the pattern of relations among its parts — and its \emph{dynamics} — how its states evolve and what they dispose the system to do. It says nothing about the \emph{intrinsic, categorical nature} of whatever realises that structure: what the components \emph{are} in themselves, beneath their relational profile. The picture is the wiring diagram. A circuit diagram fixes every connection and every response, yet on its own tells you nothing about what the components are physically made of.

This is the structural-dynamical thesis at the heart of \emph{Russellian monism}, traced to Russell's \emph{The Analysis of Matter} (see \hyperref[source:source-russell-1927]{\textsc{Russell — The Analysis of Matter (1927)}}) and revived by Chalmers (see \hyperref[source:source-chalmers-2003]{\textsc{Chalmers — Consciousness and Its Place in Nature (2003)}}): physics describes the world's abstract causal-structural pattern and is silent on its intrinsic nature. As stated, the claim is \emph{neutral} — it does not say that anything important lies outside structure and dynamics, only that such description leaves intrinsic nature out. A structuralist, in particular an \emph{ontic structural realist} who holds that there is no further categorical nature (or that talk of ``intrinsic nature'' is empty), would deny it.

It is more general than \hyperref[entry:claim-phenomenal-residue]{\textsc{the phenomenal-residue claim}}, which says that functional role under-specifies \emph{qualia} specifically. This claim makes the broader point that structural-dynamical description omits \emph{intrinsic nature} as such — of which the phenomenal is only the most salient candidate to fill the gap.

\begin{plainterms}
Describe anything purely by its function and you capture its \emph{structure} (how the parts relate) and its \emph{dynamics} (how its states change over time) — but nothing about what the parts \emph{are} in themselves, their inner nature. It's the wiring-diagram point again: the diagram shows the connections, never what the components are actually made of. On its own this stays neutral about whether anything important hides in that left-out inner nature — it just marks that function leaves it out.
\end{plainterms}

\webheading{See also}
\begin{seealso}
  \item \hyperref[entry:claim-phenomenal-residue]{\textsc{phenomenal residue}} — the narrower companion: function under-specifies \emph{qualia}; this claim generalises to intrinsic nature as such.
  \item \textsc{Functionalism} — the role-based view of mind whose descriptive reach this claim bounds.
  \item \hyperref[entry:position-russellian-monism]{\textsc{Russellian monism}} — the position built on this thesis: consciousness as (or grounding) the intrinsic nature physics leaves blank.
\end{seealso}

\entryhead{claim-subjects-actually-merge-and-divide}{Subjects merge and divide}{Claim \,\textperiodcentered\, editorial}{Subjects merge and divide}{There are actual cases in which subject boundaries shift with the degree of integration}{$\odot$}
The boundary around a conscious subject is not fixed once and for all but appears, in real cases, to shift with how tightly the underlying parts are integrated. Three kinds of case point the same way. In \emph{split-brain} patients, severing the corpus callosum can leave the two hemispheres behaving as partly independent centres of experience that a healthy, fully connected brain does not display (see \hyperref[source:source-nagel-1971]{\textsc{Nagel — Brain Bisection and the Unity of Consciousness (1971)}}). In \emph{craniopagus} twins whose brains are joined — the Hogan twins, linked by a ``thalamic bridge'' — there are reports of sensory experience shared across the two bodies (see \hyperref[source:source-dominus-2011]{\textsc{Dominus — Could Conjoined Twins Share a Mind? (2011)}}). And at the \emph{group} level, practitioners of certain meditative and tantric techniques, people on some psychedelics, and members of a jazz band or a well-synced team report a felt collapse of the boundary between several minds into a single acting ``we.''

These cases are offered as a defeater for the claim that subject-summing never happens: they suggest that whether a system houses one subject or several can vary with its degree of internal integration. The cases are contested in interpretation — a sceptic will read the group ``we'' as a vivid \emph{representation} of unity rather than a literal further subject, and dispute how many subjects a split brain ``really'' contains — which is why this claim states only that such cases \emph{occur and shift with integration}, leaving the inference to genuine combination to the argument it feeds, \textsc{Combining subjects}. Luke Roelofs surveys this same range of cases at length (see \hyperref[source:source-roelofs-2019]{\textsc{Roelofs — Combining Minds: How to Think about Composite Subjectivity (2019)}}).

\webheading{See also}
\begin{seealso}
  \item \textsc{Combining subjects} — the rebuttal to the combination problem that this claim helps support.
  \item \hyperref[entry:claim-subject-unity-graded-by-integration]{\textsc{Graded subject unity}} — the companion premise that reads these cases as a continuous, integration-dependent gradient of subjecthood.
  \item \hyperref[entry:claim-micro-subjects-do-not-sum]{\textsc{Micro-subjects do not sum}} — the claim these cases are marshalled against.
\end{seealso}

\end{multicols}
\bookgroup{T}
\begin{multicols}{2}
\entryhead{concept-materialism-types}{Type-A / Type-B / Type-C materialism}{Concept \,\textperiodcentered\, after Chalmers}{Type-A / Type-B / Type-C materialism}{Type-A / Type-B / Type-C materialism (responses to the hard problem)}{$\blacklozenge$}
Chalmers' \textbf{A/B/C taxonomy} is a way of sorting \emph{physicalist} responses to the \hyperref[entry:concept-hard-problem]{\textsc{Hard problem}} of consciousness — the problem of why physical processing should be accompanied by felt experience at all. It classifies them not by their metaphysics, on which all physicalists agree (everything is, at bottom, physical), but by a single epistemological question: how does each treat the \textbf{epistemic gap} between the physical truths and the truths about experience — the apparent fact that no amount of physical information lets you \emph{deduce}, a priori, what an experience is like? The scheme was introduced by David Chalmers in ``Consciousness and Its Place in Nature'' (\hyperref[source:source-chalmers-2003]{\textsc{Chalmers — Consciousness and Its Place in Nature (2003)}}; the underlying apparatus is developed in \hyperref[source:source-chalmers-1996]{\textsc{Chalmers — The Conscious Mind (1996)}}), and the three physicalist options it labels are Type-A, Type-B, and Type-C.

\webheading{The three types}

\begin{itemize}
  \item \emph{Type-A materialism} \textbf{denies the epistemic gap.} Once the functions are explained — discrimination, reportability, the control of behaviour — there is, on this view, no further fact left over: a \hyperref[entry:concept-philosophical-zombie]{\textsc{Philosophical zombie}} is not really conceivable on reflection, and Mary (see \hyperref[entry:concept-knowledge-argument]{\textsc{Knowledge argument}}) learns no new fact when she first sees red. Analytic functionalism, behaviourism, eliminativism, and \textbf{illusionism} (see \textsc{Illusionism}, \textsc{Phenomenality as illusion}) are its species; the standalone entry is \hyperref[entry:position-type-a-materialism]{\textsc{Type-A materialism}}.
  \item \emph{Type-B materialism} \textbf{accepts the epistemic gap but denies an ontological one.} Phenomenal and physical truths are linked by identities that are \emph{necessary but knowable only a posteriori}, on the model of ``water is H₂O'', so a zombie is conceivable yet metaphysically impossible — which is why \hyperref[entry:claim-conceivability-entails-possibility]{\textsc{Conceivability entails possibility}} is denied as a guide (see \hyperref[entry:argument-conceivability-not-guide-to-possibility]{\textsc{A-posteriori-necessity objection}}). The gap is \emph{explained}, not closed, by the phenomenal-concept strategy (see \textsc{Phenomenal-concept deflation}, \textsc{Phenomenal-concept trade-off}). The standalone entry is \hyperref[entry:position-type-b-materialism]{\textsc{Type-B materialism}}.
  \item \emph{Type-C materialism} concedes the gap looks unbridgeable \emph{now} but holds it is closable \textbf{in principle}, with conceptual or scientific progress. Chalmers argues the position is unstable, collapsing on reflection into Type-A (if the gap really closes) or Type-B (if it remains but proves harmless). The standalone entry is \hyperref[entry:position-type-c-materialism]{\textsc{Type-C materialism}}.
\end{itemize}

\webheading{The non-materialist remainder}

For orientation, the same scheme runs on past physicalism. \emph{Type-D} is interactionist dualism and \emph{Type-E} is epiphenomenalism (both \textsc{Dualism}); \emph{Type-F} is Russellian monism / panprotopsychism (see \hyperref[entry:position-russellian-monism]{\textsc{Russellian monism}}, \hyperref[entry:position-panpsychism]{\textsc{Panpsychism (position)}}). Type-A/B/C are the \emph{physicalist} options, on which consciousness reduces to or is wholly fixed by the physical; D/E/F keep consciousness fundamental. The dividing line is whether the epistemic gap is taken to mark a real division in nature.

\webheading{The dispute it feeds}

The taxonomy is a map of the live replies to the anti-physicalist thought experiments — the \hyperref[entry:argument-zombie-conceivability]{\textsc{Zombie argument}} and the \hyperref[entry:concept-knowledge-argument]{\textsc{Knowledge argument}}. Almost every current reply is Type-A (deny the intuition the experiment trades on) or Type-B (grant the intuition, deny its modal or ontological force), which is why so much of the debate turns on the two questions those types divide over: is the epistemic gap real, and if so does it entail an ontological gap?

\begin{plainterms}
A cheat-sheet for the ways an ``it's all physical'' theorist can answer the hard problem. \textbf{Type-A:} deny there's any real puzzle — once you've explained what the brain \emph{does}, there's nothing left to explain (the ``you only \emph{think} there's an extra feel'' camp). \textbf{Type-B:} admit it \emph{feels} like there's an unexplained extra, but insist that's just how our concepts work — experience really \emph{is} some physical process, the way water really is H₂O even though it didn't have to feel that way. \textbf{Type-C:} ``we can't bridge it yet, but we will'' — which, Chalmers argues, quietly turns into A or B once you push on it. (Beyond these lie the non-physical options — dualism, and the views on which consciousness stays fundamental, like Russellian monism and panpsychism.)
\end{plainterms}

\webheading{See also}
\begin{seealso}
  \item \hyperref[entry:position-type-a-materialism]{\textsc{Type-A materialism}}, \hyperref[entry:position-type-b-materialism]{\textsc{Type-B materialism}}, \hyperref[entry:position-type-c-materialism]{\textsc{Type-C materialism}} — the three physicalist wings as standalone position entries.
  \item \hyperref[entry:concept-hard-problem]{\textsc{Hard problem}} — the problem the whole taxonomy sorts responses to.
  \item \hyperref[entry:concept-knowledge-argument]{\textsc{Knowledge argument}} and \hyperref[entry:argument-zombie-conceivability]{\textsc{Zombie argument}} — the anti-physicalist thought experiments the types are replies to.
  \item \hyperref[entry:claim-conceivability-entails-possibility]{\textsc{Conceivability entails possibility}} and \hyperref[entry:argument-conceivability-not-guide-to-possibility]{\textsc{A-posteriori-necessity objection}} — the bridge principle Type-B rejects and the argument that does the rejecting.
  \item \textsc{Phenomenal-concept deflation} and \textsc{Phenomenal-concept trade-off} — the phenomenal-concept strategy by which Type-B explains the gap.
  \item \textsc{Illusionism} and \textsc{Phenomenality as illusion} — the strongest Type-A line, that the felt residue is an introspective illusion.
  \item \textsc{Dualism}, \hyperref[entry:position-russellian-monism]{\textsc{Russellian monism}}, \hyperref[entry:position-panpsychism]{\textsc{Panpsychism (position)}} — the non-physicalist Types D/E/F that keep consciousness fundamental.
  \item \hyperref[source:source-chalmers-2003]{\textsc{Chalmers — Consciousness and Its Place in Nature (2003)}}, \hyperref[source:source-chalmers-1996]{\textsc{Chalmers — The Conscious Mind (1996)}} — Chalmers' statement of the taxonomy and the apparatus behind it.
\end{seealso}

\entryhead{position-type-a-materialism}{Type-A materialism}{Position \,\textperiodcentered\, after Chalmers}{Type-A materialism}{Type-A materialism (denial of the epistemic gap)}{$\bigstar$}
\textbf{Type-A materialism} is the most deflationary wing of physicalism: it denies that there is any \emph{epistemic gap} between the physical truths and the truths about consciousness. Once you have explained all the functions — the discrimination, the reportability, the control of behaviour, the access of information to reasoning — there is, on this view, simply nothing further left to explain. The felt, ``what-it's-like'' residue that the \hyperref[entry:concept-hard-problem]{\textsc{Hard problem}} points at is not a hard remainder but a non-existent one, or at most an artefact of how we introspect. Accordingly a type-A theorist holds that a \hyperref[entry:concept-philosophical-zombie]{\textsc{Philosophical zombie}} is not really conceivable on reflection (\hyperref[entry:claim-zombie-conceivable]{\textsc{Zombie conceivability}} is denied) and that Mary (see \hyperref[entry:concept-knowledge-argument]{\textsc{Knowledge argument}}), leaving her black-and-white room, learns no new \emph{fact} — at most a new ability (see \hyperref[entry:claim-mary-gains-abilities-not-facts]{\textsc{Mary gains abilities}}). It commits to the core physicalist thesis that the mental just is the physical (see \textsc{Mental events are physical}) and rejects any leftover phenomenal fact beyond the functional story (see \hyperref[entry:claim-phenomenal-residue]{\textsc{Phenomenal residue}}). It is \textbf{type-A} in Chalmers' taxonomy (\hyperref[entry:concept-materialism-types]{\textsc{Type-A / Type-B / Type-C materialism}}; \hyperref[source:source-chalmers-2003]{\textsc{Chalmers — Consciousness and Its Place in Nature (2003)}}).

\textbf{What distinguishes it from neighbouring views.} Type-A is a \emph{family}, not a single doctrine: analytic functionalism, logical behaviourism, \textsc{Eliminative materialism}, and \textsc{Illusionism} are all type-A, because each in its own way refuses to grant the explanatory gap any genuine purchase. What unites them and sets them apart from the rest of physicalism is the denial that the gap is even real. This is exactly the point on which it splits from \hyperref[entry:position-type-b-materialism]{\textsc{Type-B materialism}}, which \emph{concedes} that the gap is real (the zombie is conceivable, Mary's surprise genuine) and tries instead to show it is merely epistemic; and from \hyperref[entry:position-type-c-materialism]{\textsc{Type-C materialism}}, which concedes the gap looks unbridgeable for now but bets it will close. Against all of them stand the realist anti-physicalisms (see \textsc{Dualism}, \hyperref[entry:position-russellian-monism]{\textsc{Russellian monism}}) that take the gap to mark a real division in nature.

\begin{plainterms}
The hard-nosed ``there's nothing to explain'' camp. When people insist that even after brain science is finished there's still a left-over mystery — the \emph{feel} of pain, the \emph{redness} of red — a type-A materialist answers: no, there isn't. Explain everything the brain \emph{does} and the job is done; the sense that something extra has been skipped is itself just a trick of how we look inward. So on this view a perfect zombie (physically just like you but with the lights off inside) isn't really imaginable once you think it through, and the famous colour-scientist Mary learns a new \emph{skill}, not a new \emph{fact}, when she first sees red. It is the most demanding line a physicalist can take, because it asks you to give up the very intuition the other camps try to keep. Contrast the \hyperref[entry:position-type-b-materialism]{\textsc{Type-B materialism}} camp, which admits the puzzle feels real and then explains it away.
\end{plainterms}

\webheading{See also}
\begin{seealso}
  \item \hyperref[entry:concept-materialism-types]{\textsc{Type-A / Type-B / Type-C materialism}} — the A/B/C taxonomy that defines this position and places it among its siblings.
  \item \textsc{Mental events are physical} — the core commitment it shares with all physicalisms.
  \item \hyperref[entry:claim-zombie-conceivable]{\textsc{Zombie conceivability}} and \hyperref[entry:claim-phenomenal-residue]{\textsc{Phenomenal residue}} — the two claims it denies, marking its denial of the gap.
  \item \textsc{Eliminative materialism} and \textsc{Illusionism} — its leading species, each a distinct way of refusing the gap.
  \item \hyperref[entry:concept-knowledge-argument]{\textsc{Knowledge argument}} and \hyperref[entry:claim-mary-gains-abilities-not-facts]{\textsc{Mary gains abilities}} — Mary's case and the ability-hypothesis reading type-A gives it.
  \item \hyperref[entry:position-type-b-materialism]{\textsc{Type-B materialism}}, \hyperref[entry:position-type-c-materialism]{\textsc{Type-C materialism}} — the other physicalist wings, which grant the gap and then either deflate or defer it.
  \item \hyperref[entry:question-hard-problem]{\textsc{Hard problem (the question)}} — the question it answers by dissolving the explanandum.
\end{seealso}

\entryhead{position-type-b-materialism}{Type-B materialism}{Position \,\textperiodcentered\, after Chalmers}{Type-B materialism}{Type-B materialism (epistemic gap without ontological gap)}{$\bigstar$}
\textbf{Type-B materialism} is the moderate, and today the most populous, wing of physicalism. Unlike \hyperref[entry:position-type-a-materialism]{\textsc{Type-A materialism}}, it \emph{grants} the epistemic gap: it concedes that a \hyperref[entry:concept-philosophical-zombie]{\textsc{Philosophical zombie}} is genuinely conceivable, that Mary (see \hyperref[entry:concept-knowledge-argument]{\textsc{Knowledge argument}}) is genuinely surprised when she first sees red, and that no amount of physical information makes the felt character of experience deducible \emph{a priori}. What it denies is that this epistemic gap reflects an \textbf{ontological} one. Physical and phenomenal truths are connected, on this view, by identities that are \emph{necessary but knowable only a posteriori} — on the model of ``water is H₂O'' or ``heat is molecular motion,'' truths that could not have been false yet had to be discovered rather than reasoned out. The mental still \emph{is} the physical (see \textsc{Mental events are physical}); the appearance of a further fact is a feature of our concepts, not of the world. The decisive move is therefore to reject the bridge from conceivability to possibility (see \hyperref[entry:claim-conceivability-entails-possibility]{\textsc{Conceivability entails possibility}}): the zombie is conceivable yet metaphysically impossible, so its conceivability proves nothing about what exists. It is \textbf{type-B} in Chalmers' taxonomy (\hyperref[entry:concept-materialism-types]{\textsc{Type-A / Type-B / Type-C materialism}}; \hyperref[source:source-chalmers-2003]{\textsc{Chalmers — Consciousness and Its Place in Nature (2003)}}).

\textbf{What distinguishes it from neighbouring views.} Type-B occupies the middle ground and is defined by both its borders. Against \hyperref[entry:position-type-a-materialism]{\textsc{Type-A materialism}} it insists the gap is real and must be \emph{explained} rather than denied — typically by the phenomenal-concept strategy, which argues that our special first-person concepts of experience are cognitively isolated in a way that generates the illusion of a further fact. Against the anti-physicalist (see \textsc{Dualism}) it insists the gap is \emph{merely} epistemic and carries no modal force. \textsc{Type-identity theory} is the classic species of type-B: it identifies mental types with brain-state types as an a-posteriori discovery, exactly the necessary-but-empirical identity type-B trades on. The standing worry is that the a-posteriori identities type-B needs look unlike ordinary scientific identities — ``water is H₂O'' can itself be explained structurally, whereas ``pain is C-fibre firing'' seems to leave the very thing in question untouched.

\begin{plainterms}
The ``yes it feels like a mystery, but it isn't really one'' camp — and the most common view among physicalist philosophers today. It agrees with the critics on the \emph{feeling}: you really can imagine a zombie, and Mary really is in for a surprise when she first sees colour. But it says that surprise is about how our minds grasp the facts, not about there being two kinds of stuff. Compare ``water is H₂O'': that turned out to be a discovery, something you couldn't have worked out from the armchair, even though water could not have been anything other than H₂O. Type-B says consciousness is like that — experience genuinely \emph{is} some brain process, even though it never \emph{feels} deducible from the physics. So being able to imagine a zombie tells you about the limits of imagination, not about what could really exist. The sharper, gap-denying alternative is \hyperref[entry:position-type-a-materialism]{\textsc{Type-A materialism}}.
\end{plainterms}

\webheading{See also}
\begin{seealso}
  \item \hyperref[entry:concept-materialism-types]{\textsc{Type-A / Type-B / Type-C materialism}} — the A/B/C taxonomy that defines this position.
  \item \textsc{Mental events are physical} — the core physicalist commitment it keeps.
  \item \hyperref[entry:claim-conceivability-entails-possibility]{\textsc{Conceivability entails possibility}} — the bridge principle it must reject to absorb the conceivable zombie.
  \item \textsc{Type-identity theory} — its classic species, an a-posteriori identity physicalism.
  \item \hyperref[entry:concept-philosophical-zombie]{\textsc{Philosophical zombie}} and \hyperref[entry:concept-knowledge-argument]{\textsc{Knowledge argument}} — the two thought experiments whose intuitions type-B grants but defuses.
  \item \hyperref[entry:position-type-a-materialism]{\textsc{Type-A materialism}}, \hyperref[entry:position-type-c-materialism]{\textsc{Type-C materialism}} — the sibling wings that deny the gap or defer it.
\end{seealso}

\entryhead{position-type-c-materialism}{Type-C materialism}{Position \,\textperiodcentered\, after Chalmers}{Type-C materialism}{Type-C materialism (the gap is closable in principle)}{$\bigstar$}
\textbf{Type-C materialism} is the optimistic, transitional wing of physicalism. It concedes that the epistemic gap between physical and phenomenal truths (see \hyperref[entry:concept-hard-problem]{\textsc{Hard problem}}) looks unbridgeable \emph{as things now stand} — we cannot presently see how the physical facts could entail the facts about experience — but it holds that the gap is closable \textbf{in principle}, that conceptual or scientific advance will eventually let us see the connection we now cannot. On this view our current sense of mystery reflects the immaturity of our concepts, not a permanent feature of the world; given the right future theory, a \hyperref[entry:concept-philosophical-zombie]{\textsc{Philosophical zombie}} would come to look as plainly impossible as a married bachelor. It keeps the core physicalist commitment (see \textsc{Mental events are physical}) and answers the hard problem with a promissory note rather than a present account. It is \textbf{type-C} in Chalmers' taxonomy (\hyperref[entry:concept-materialism-types]{\textsc{Type-A / Type-B / Type-C materialism}}; \hyperref[source:source-chalmers-2003]{\textsc{Chalmers — Consciousness and Its Place in Nature (2003)}}).

\textbf{What distinguishes it from neighbouring views.} Type-C is defined by its instability. Chalmers argues that it is not a stable resting place but a slope: pressed on what the promised future insight would look like, the type-C theorist must say either that the gap genuinely \emph{closes} — in which case the appearance of a gap was illusory all along, and the view collapses into \hyperref[entry:position-type-a-materialism]{\textsc{Type-A materialism}} — or that the gap \emph{remains} but is shown to be harmless, an artefact of our concepts, in which case it has become \hyperref[entry:position-type-b-materialism]{\textsc{Type-B materialism}}. There is, on Chalmers' diagnosis, no third thing for ``closable in principle but not yet'' to mean. This sets type-C apart from its pessimistic mirror image, \textsc{Mysterianism}, which agrees the gap is currently unbridged but holds it is unbridgeable \emph{by us} forever — keeping the optimism's epistemic framing while denying its hope. Where type-A is deflationary and type-B is conciliatory, type-C is a bet on the future that, Chalmers contends, cannot be held without sliding into one of its neighbours.

\begin{plainterms}
The ``we can't explain it yet, but we will'' camp. A type-C materialist looks at the puzzle of consciousness, admits that right now nobody can show how grey matter adds up to the taste of coffee, but trusts that future science or sharper concepts will eventually make the link obvious — the way once-baffling puzzles about life turned out to be ordinary chemistry. The catch, which Chalmers presses, is that this turns out not to be a real third option. Ask what the future breakthrough would actually show, and you end up either saying ``there was never really a gap'' (which is the \hyperref[entry:position-type-a-materialism]{\textsc{Type-A materialism}} line) or ``the gap is real but harmless'' (which is the \hyperref[entry:position-type-b-materialism]{\textsc{Type-B materialism}} line). So type-C tends to be a waystation rather than a home. Its gloomy twin is \textsc{Mysterianism}: same starting point, but the conviction that the answer is forever beyond human reach.
\end{plainterms}

\webheading{See also}
\begin{seealso}
  \item \hyperref[entry:concept-materialism-types]{\textsc{Type-A / Type-B / Type-C materialism}} — the A/B/C taxonomy that defines this position and the instability charge against it.
  \item \textsc{Mental events are physical} — the physicalist commitment it shares with its siblings.
  \item \hyperref[entry:position-type-a-materialism]{\textsc{Type-A materialism}}, \hyperref[entry:position-type-b-materialism]{\textsc{Type-B materialism}} — the two views Chalmers argues type-C collapses into.
  \item \textsc{Mysterianism} — the pessimistic counterpart that shares the framing but abandons the hope.
  \item \hyperref[entry:question-hard-problem]{\textsc{Hard problem (the question)}} — the question it answers on credit rather than in hand.
\end{seealso}

\end{multicols}
\bookgroup{V}
\begin{multicols}{2}
\entryhead{concept-introspective-gap-varieties}{Varieties of introspective explanatory gap}{Concept \,\textperiodcentered\, editorial}{Varieties of introspective explanatory gap}{Varieties of introspective explanatory gap (ontological / architectural-access / conceptual-translation / training-contamination)}{$\blacklozenge$}
When a self-modelling system ``cannot explain its own experience to itself,'' that complaint can mean four quite different things. This entry is a \emph{taxonomy} that keeps them apart, so that a merely \emph{felt} explanatory gap is not automatically read as the metaphysical hard problem of consciousness. These four varieties can overlap in a single case, yet each has a distinct source and calls for a distinct remedy. Conflating them is how a modest puzzle gets mistaken for a deep one.

A quick illustration of why the distinction earns its keep: imagine a person (or an AI) who reports that no account of its inner workings ever feels like it explains the experience itself. That one report could be voicing a gap in the \emph{world} (the facts run out), a gap in the \emph{self} (the relevant facts are present but unreachable), a gap in its \emph{vocabulary} (the facts are reachable but won't translate), or no gap at all (it is echoing things it has heard others say). Telling these apart is the whole point.

\webheading{The four varieties}

The \textbf{ontological gap} is the one in the world: even given a complete physical or computational account of the system, why is there any felt experience at all? This is the hard problem proper (see \hyperref[entry:concept-hard-problem]{\textsc{Hard problem}}); on the anti-physicalist reading the unexplained residue lies in reality itself, not in the knower (see \hyperref[entry:claim-phenomenal-residue]{\textsc{Phenomenal residue}}, \textsc{Explanatory-gap argument}).

The \textbf{architectural-access gap} is a gap in the \emph{self}, not in the world. The facts that \emph{would} explain the experience are realised in the system's own processing, but they are not available to the machinery that produces its self-reports, so the system simply cannot reach its own internal variables (see \textsc{Limited inside vantage}). This is a limit on what philosophers call access consciousness (see \hyperref[entry:concept-access-vs-phenomenal-consciousness]{\textsc{Access vs. phenomenal consciousness}}), the information a state makes available for report and control.

The \textbf{conceptual-translation gap} is a gap in vocabulary. Here the system possesses \emph{both} a mechanistic description of itself and an introspective one, yet has no internal route from the one to the other: the bridge an outside theorist can build between ``neurons firing thus'' and ``feels like this'' is not traversable from the inside (see \textsc{Internally traversable explanation}). Nothing is hidden and nothing is missing; the two descriptions just never meet in the system's own cognition.

The \textbf{training-contamination gap} is, strictly, no gap of the system's own. The system merely reproduces human hard-problem discourse it absorbed in training, with no underlying difficulty behind the words. For an experiment probing whether a machine has a genuine explanatory gap, this is the confound that must be controlled for before any of the other three can be credited.

\webheading{The dispute it structures}

The taxonomy was built to discipline one question in particular: which kind of gap, if any, an introspecting but consciousness-naive language model would actually exhibit (see \textsc{LLM's own explanatory gap}). By isolating the architectural and conceptual readings, it also lets a felt self-explanatory gap be stated as a concrete \emph{mechanism} for McGinn-style cognitive closure, the condition of a mind constitutively unable to grasp certain concepts (see \textsc{Possibility of cognitive closure}), without thereby asserting any ontological residue in the world (see \textsc{Architectural self-explanation gap}). The architectural and conceptual gaps, in short, give the mysterian's ``closure'' an engine that a physicalist can accept.

\begin{plainterms}
When something can't ``explain its own experience to itself,'' that can mean four very different things, and they get muddled together. (1) \textbf{Ontological:} even after you know everything about the machinery, why is there a felt experience at all? That is the genuine hard problem. (2) \textbf{Architectural-access:} the facts that would explain it are inside the system but it can't read them (like not being able to feel your own neurons fire). (3) \textbf{Conceptual-translation:} it has the mechanical story \emph{and} the first-person story but no inner bridge between them; an outside scientist can make that translation, while the system itself cannot. (4) \textbf{Contamination:} it's just parroting things humans have written about the hard problem. Most of the interest is in telling (2) and (3) apart from (1) and (4).
\end{plainterms}

\webheading{See also}
\begin{seealso}
  \item \hyperref[entry:concept-hard-problem]{\textsc{Hard problem}} — the ontological variety in its sharpest form: why physical processing is accompanied by experience at all.
  \item \hyperref[entry:claim-phenomenal-residue]{\textsc{Phenomenal residue}} — the anti-physicalist reading of the ontological gap: experience is not exhausted by functional or causal role.
  \item \textsc{Explanatory-gap argument} — treats the gap as evidence against uniform substrate-independence; the ontological reading put to work.
  \item \textsc{Limited inside vantage} — the engine of the architectural-access variety: a system's view of itself is informationally poorer than an outsider's.
  \item \hyperref[entry:concept-access-vs-phenomenal-consciousness]{\textsc{Access vs. phenomenal consciousness}} — Block's distinction; the architectural-access gap is a shortfall on the \emph{access} side, separate from anything phenomenal.
  \item \textsc{Internally traversable explanation} — the criterion that defines the conceptual-translation variety: an explanation satisfies a system only if the system's own operations can traverse it.
  \item \textsc{LLM's own explanatory gap} — the live question this taxonomy was built to sharpen: which variety, if any, an introspecting LLM would show.
  \item \textsc{Architectural self-explanation gap} — uses the taxonomy to locate a self-explanatory gap in architecture rather than ontology.
  \item \textsc{Possibility of cognitive closure} — McGinn's mysterianism; the architectural and conceptual varieties offer it a mechanism without an ontological residue.
\end{seealso}

\end{multicols}
\bookgroup{W}
\begin{multicols}{2}
\entryhead{question-what-does-mary-learn}{What Mary learns}{Question}{What Mary learns}{What, if anything, does Mary gain when she first sees red?}{?}
This is the question on which the \hyperref[entry:concept-knowledge-argument]{\textsc{knowledge argument}} turns: when Mary the colour scientist leaves her black-and-white room and sees red for the first time, what — if anything — does she gain? Everyone in the debate agrees that \emph{something} changes. The dispute is over the right \emph{category} for that something, and over whether it is the kind of thing that physicalism cannot accommodate.

\webheading{The categories in play}

The candidate answers trade on a small set of distinctions about kinds of knowledge:

\emph{Propositional knowledge} (knowledge-\emph{that}) is the grasp of a fact — something that is true or false. Only a gain of this kind threatens physicalism directly, because only a new \emph{fact} could be a fact the physics left out.

\emph{Know-how} is a matter of abilities and skills. Know-how is \emph{had} rather than true: you can gain the ability to ride a bicycle without learning any new proposition. If what Mary gains is know-how, she need not have learned a new fact at all.

\emph{Acquaintance} is a direct cognitive relation to a thing — knowing \emph{it} rather than knowing \emph{that} something is so. Some hold this is a third category, reducible neither to propositional knowledge nor to know-how: Mary becomes acquainted with an experience she had only ever read about.

\emph{Mode of presentation}, or \emph{guise}, is the way a fact is grasped. One and the same fact can be known under two guises, so ``new knowledge'' need not mean ``new fact'': Mary may grasp a physical fact she already knew, now under a new, experiential guise.

\webheading{A further, separable issue}

Even if the categorisation question is settled in the proponent's favour — granting the epistemic conclusion that \hyperref[entry:claim-some-truths-about-experience-not-physically-deducible]{\textsc{some truths about experience are not physically deducible}} — a second question remains: does that epistemic gap support an \emph{ontological} anti-physicalism, a claim about what exists? The step from ``cannot be deduced from the physical'' to ``is not a physical fact'' is a separate dispute from the categorisation dispute above, and may deserve its own question node.

\begin{plainterms}
Mary is a scientist who knows every physical fact about colour but has lived her whole life in black and white. One day she steps out and sees red. The question: what exactly did she just get? A brand-new fact the physics books left out? A new skill, like learning to ride a bike? A first meeting with something she'd only read about? Or just a new way of grasping a fact she already knew? Which description is right decides whether the story proves anything about the limits of physical science.
\end{plainterms}

\webheading{See also}
\begin{seealso}
  \item \hyperref[entry:concept-knowledge-argument]{\textsc{Knowledge argument}} — the thought experiment this question organises; the answers here are the moves in that debate.
  \item \hyperref[entry:claim-some-truths-about-experience-not-physically-deducible]{\textsc{Non-deducible experiential truths}} — the epistemic conclusion whose further, ontological reading is the separable sub-issue noted above.
\end{seealso}

\end{multicols}
\bookgroup{Z}
\begin{multicols}{2}
\entryhead{argument-zombie-conceivability}{Zombie argument}{Argument \,\textperiodcentered\, after Chalmers}{Zombie argument}{The zombie / conceivability argument against physicalism}{$\vdash$}
The \emph{zombie argument} (also the \emph{conceivability argument}) tries to show that consciousness is not physical by way of an imagined creature: a \hyperref[entry:concept-philosophical-zombie]{\textsc{philosophical zombie}}, a being microphysically and functionally identical to a conscious person but with no inner experience at all (see \hyperref[source:source-chalmers-1996]{\textsc{Chalmers — The Conscious Mind (1996)}}). If such a duplicate is even possible, the complete physical story about a person leaves the facts about experience unsettled — so those facts must be something beyond the physical. It is a \emph{modal} argument: it reasons from what is possible, not from observed evidence, and arrives at a conclusion about the fundamental nature of mind.

\webheading{The inference}

\begin{enumerate}
  \item \emph{Scenario and intuition} (see \hyperref[entry:claim-zombie-conceivable]{\textsc{Zombie conceivability}}): a microphysical duplicate wholly lacking phenomenal consciousness is \emph{ideally conceivable} — coherently and stably imaginable, with no hidden contradiction surfacing on careful rational reflection.
  \item \emph{Bridge} (see \hyperref[entry:claim-conceivability-entails-possibility]{\textsc{Conceivability entails possibility}}): ideal conceivability entails metaphysical possibility.
\end{enumerate}

From these it follows that a zombie is metaphysically possible; so the complete physical facts do not necessitate the facts about experience, and phenomenal consciousness is \emph{over and above} the physical (see \hyperref[entry:claim-consciousness-fundamental-nonphysical]{\textsc{Fundamental non-physical consciousness}}). The conclusion does double duty against physicalism: it \texttt{supports} the weaker thesis that \hyperref[entry:claim-phenomenal-residue]{\textsc{phenomenal states are not exhausted by their functional role}} (a fortiori — if the physical facts do not even necessitate experience, they certainly do not exhaust it), and it \texttt{attacks} both the physicalist thesis that \textsc{mental events are physical} (the central commitment of \hyperref[entry:position-physicalism]{\textsc{Physicalism}}) and the uniform, phenomenal-covering reading of \textsc{substrate independence}.

\webheading{The weakest premise}

The pressure point is the \textbf{bridge} (see \hyperref[entry:claim-conceivability-entails-possibility]{\textsc{Conceivability entails possibility}}). The most common reply, from the \emph{type-B physicalist} (see \hyperref[entry:concept-materialism-types]{\textsc{Type-A / Type-B / Type-C materialism}}), grants that a zombie is conceivable but denies that it is possible, on the model of necessary \emph{a posteriori} identities: just as ``water is not H₂O'' is conceivable yet impossible, a zombie may be imaginable yet metaphysically ruled out. This is the \emph{bridge} critical question that the \emph{thought-experiment pattern} standardly invites. A related move, the \textsc{phenomenal-concept strategy}, explains the conceivability away as an artefact of the special concepts we use to think about experience, rather than a window onto possibility.

A second line of reply attacks the first premise instead. The \emph{type-A physicalist} (see \hyperref[entry:concept-materialism-types]{\textsc{Type-A / Type-B / Type-C materialism}}) and the illusionist press the pattern's \emph{coherence} critical question, denying \hyperref[entry:claim-zombie-conceivable]{\textsc{that a zombie is conceivable}} at all: on full reflection there is no further phenomenal fact left to subtract, so the description only \emph{seems} coherent. Where the argument lands is itself disputed. The very same premises are read by some not as dualism but as \hyperref[entry:position-russellian-monism]{\textsc{Russellian monism}} — the view that the intrinsic nature of the physical is already experiential — which secures the conclusion while avoiding the epiphenomenalism that dogs \textsc{Dualism}.

The argument is distinct from the \textsc{explanatory-gap argument}: that one runs an \emph{evidential} comparison on how stubbornly the gap persists and only challenges the \emph{uniform scope} of substrate independence, whereas this is a \emph{modal} argument for the stronger conclusion that the phenomenal is non-physical. (The explanatory-gap argument's premises already lean on zombie conceivability; this node factors that move out.)

\begin{plainterms}
Picture an atom-for-atom copy of a conscious person that has no inner experience at all — a ``zombie.'' Step one: that seems coherently imaginable. Step two (the controversial bit): if something is truly, coherently imaginable, it's genuinely possible. Put them together and a zombie is possible — which means the physical facts alone don't force experience to exist, so consciousness must be something extra beyond the physical. The argument lives or dies on step two: opponents say you can imagine plenty of things that turn out impossible (like water \emph{not} being H₂O), so imaginability isn't proof of possibility.
\end{plainterms}

\webheading{See also}
\begin{seealso}
  \item \hyperref[entry:concept-philosophical-zombie]{\textsc{Philosophical zombie}} — the imagined duplicate the argument runs on.
  \item \hyperref[entry:concept-hard-problem]{\textsc{Hard problem}} — the puzzle of why physical processing is felt at all; the intuition this argument formalises.
  \item \textsc{Functionalism} — the theory of mind whose phenomenal-covering ambitions the argument's conclusion targets.
  \item \hyperref[entry:concept-materialism-types]{\textsc{Type-A / Type-B / Type-C materialism}} — the type-A / type-B taxonomy organising the two main replies.
  \item \textsc{Explanatory-gap argument} — the evidential cousin, contrasted above.
  \item \hyperref[entry:position-russellian-monism]{\textsc{Russellian monism}} — the non-dualist landing spot some give the same premises.
  \item \textsc{Dualism} — the standard reading, with its attendant epiphenomenalism worry.
\end{seealso}

\entryhead{claim-zombie-conceivable}{Zombie conceivability}{Claim \,\textperiodcentered\, after Chalmers}{Zombie conceivability}{A microphysical duplicate of a conscious being that wholly lacks phenomenal consciousness is conceivable}{$\odot$}
A being identical to a conscious person in every microphysical and functional respect, yet with no phenomenal consciousness — a \hyperref[entry:concept-philosophical-zombie]{\textsc{philosophical zombie}} — is \emph{ideally conceivable}: coherently and stably imaginable, with no hidden contradiction revealed on ideal rational reflection (see \hyperref[source:source-chalmers-1996]{\textsc{Chalmers — The Conscious Mind (1996)}}). This is the scenario-and-intuition premise of the zombie argument.

The claim is about \emph{conceivability}, not yet possibility. It does not by itself assert that zombies are metaphysically possible; that further step needs the modal bridge, \hyperref[entry:claim-conceivability-entails-possibility]{\textsc{Conceivability entails possibility}}. The qualifier ``ideal'' is doing real work: the claim is not that a zombie is easy to picture, but that the description survives idealised reflection — the kind of careful scrutiny that would expose a contradiction if there were one.

Those who deny it are chiefly type-A physicalists (see \hyperref[entry:concept-materialism-types]{\textsc{Type-A / Type-B / Type-C materialism}}) and illusionists, who hold that on full reflection the zombie description \emph{is} incoherent — there is no further phenomenal fact left over to subtract — so it is not genuinely conceivable, only apparently so.

This is distinct from the modal bridge \hyperref[entry:claim-conceivability-entails-possibility]{\textsc{Conceivability entails possibility}} and from \hyperref[entry:claim-phenomenal-residue]{\textsc{the thesis that phenomenal states are not exhausted by their role}}: it asserts only the coherent imaginability of the absent-consciousness duplicate.

\end{multicols}
\sourcesband
\begin{sourcelist}
\sourceitem{source-armstrong-1968}{Armstrong, D. M. (1968). \emph{A materialist theory of the mind}. Routledge \& Kegan Paul.}
\sourceitem{source-block-1978}{Block, N. (1978). Troubles with functionalism. In C. W. Savage (Ed.), \emph{Perception and cognition: Issues in the foundations of psychology} (Minnesota Studies in the Philosophy of Science, Vol. 9, pp. 261–325). University of Minnesota Press. (Introduces the Chinese Nation / homunculi-headed robot.)}
\sourceitem{source-breuer-1995}{Breuer, T. (1995). The impossibility of accurate state self-measurements. \emph{Philosophy of Science}, \emph{62}(2), 197–214. \url{https://www.jstor.org/stable/188430}}
\sourceitem{}{Cambridge: \url{https://www.cambridge.org/core/journals/philosophy-of-science/article/abs/impossibility-of-accurate-state-selfmeasurements/80B368D210379DA587D41603B551B95D}}
\sourceitem{}{PDF: \url{https://cqi.inf.usi.ch/qic/Breuer95.pdf} · JSTOR: \url{https://www.jstor.org/stable/188430}}
\sourceitem{source-chalmers-1995}{Chalmers, D. J. (1995). Facing up to the problem of consciousness. \emph{Journal of Consciousness Studies, 2}(3), 200–219.}
\sourceitem{source-chalmers-1996}{Chalmers, D. J. (1996). \emph{The conscious mind: In search of a fundamental theory}. Oxford University Press. (The zombie and the two-dimensional conceivability→possibility apparatus (ch. 3–4); the failure of logical supervenience of the phenomenal on the physical; naturalistic dualism and the epiphenomenalism worry that pushes toward Russellian monism.)}
\sourceitem{source-chalmers-2002}{Chalmers, D. J. (2002). On sense and intension. \emph{Philosophical Perspectives}, \emph{16}, 135–182. (The two-dimensional framework as an account of the \emph{epistemic} dimension of meaning: primary intensions defined over worlds considered as actual, secondary intensions over worlds considered as counterfactual; the bearing on the necessary a posteriori and on the conceivability–possibility link.)}
\sourceitem{source-chalmers-2003}{Chalmers, D. J. (2003). Consciousness and its place in nature. In S. P. Stich \& T. A. Warfield (Eds.), \emph{The Blackwell guide to philosophy of mind} (pp. 102–142). Blackwell.}
\sourceitem{source-chalmers-2017}{Chalmers, D. J. (2017). Panpsychism and panprotopsychism. In G. Brüntrup \& L. Jaskolla (Eds.), \emph{Panpsychism: Contemporary perspectives} (pp. 19–47). Oxford University Press.}
\sourceitem{source-chalmers-2018}{Chalmers, D. J. (2018). The meta-problem of consciousness. \emph{Journal of Consciousness Studies}, \emph{25}(9–10), 6–61.}
\sourceitem{source-coleman-2014}{Coleman, S. (2014). The real combination problem: Panpsychism, micro-subjects, and emergence. \emph{Erkenntnis}, \emph{79}(1), 19–44. \url{https://doi.org/10.1007/s10670-013-9431-x}}
\sourceitem{source-conee-1994}{Conee, E. (1994). Phenomenal knowledge. \emph{Australasian Journal of Philosophy}, \emph{72}(2), 136–150. \url{https://doi.org/10.1080/00048409412345971}}
\sourceitem{source-dennett-1988}{Dennett, D. C. (1988). Quining qualia. In A. J. Marcel \& E. Bisiach (Eds.), \emph{Consciousness in contemporary science} (pp. 42–77). Oxford University Press.}
\sourceitem{source-dennett-1991}{Dennett, D. C. (1991). \emph{Consciousness explained}. Little, Brown and Company.}
\sourceitem{source-dennett-2005}{Dennett, D. C. (2005). What RoboMary knows. In \emph{Sweet dreams: Philosophical obstacles to a science of consciousness} (pp. 103–130). MIT Press.}
\sourceitem{source-descartes-meditations}{Descartes, R. (1996). \emph{Meditations on first philosophy} (J. Cottingham, Trans.). Cambridge University Press. (Original work published 1641)}
\sourceitem{source-dominus-2011}{Dominus, S. (2011, May 29). Could conjoined twins share a mind? \emph{The New York Times Magazine}. \url{https://www.nytimes.com/2011/05/29/magazine/could-conjoined-twins-share-a-mind.html}}
\sourceitem{source-goff-2017}{Goff, P. (2017). \emph{Consciousness and fundamental reality}. Oxford University Press.}
\sourceitem{source-goff-2019}{Goff, P. (2019). \emph{Galileo's error: Foundations for a new science of consciousness}. Pantheon Books.}
\sourceitem{source-horgan-1984}{Horgan, T. (1984). Jackson on physical information and qualia. \emph{The Philosophical Quarterly}, \emph{34}(135), 147–152. \url{https://doi.org/10.2307/2219508}}
\sourceitem{source-jackson-1982}{Jackson, F. (1982). Epiphenomenal qualia. \emph{The Philosophical Quarterly}, \emph{32}(127), 127–136. (Mary, who knows all the physical facts about colour vision from inside a black-and-white room, learns something on first seeing red. See also Jackson, F. (1986). What Mary didn't know. \emph{The Journal of Philosophy}, \emph{83}(5), 291–295.)}
\sourceitem{}{Notable: Jackson \textbf{later recanted}, coming to reject the argument and accept physicalism (the ``knowledge'' being a new \emph{acquaintance/representation}, not a new fact) — recorded so the web does not attribute to him a view he abandoned; the argument is imported on its merits, not as his settled view.}
\sourceitem{source-jackson-1998}{Jackson, F. (1998). \emph{From metaphysics to ethics: A defence of conceptual analysis}. Clarendon Press. (Reference-fixing descriptions and the two-dimensional treatment of the necessary a posteriori; the ``location problem'' and a priori entailment of the macro from the micro. See also Chalmers, D. J., \& Jackson, F. (2001). Conceptual analysis and reductive explanation. \emph{The Philosophical Review}, \emph{110}(3), 315–360.)}
\sourceitem{source-james-1890}{James, W. (1890). \emph{The principles of psychology} (Vol. 1). Henry Holt. (See ch. 6, ``The mind-stuff theory.'')}
\sourceitem{source-kim-2005}{Kim, J. (2005). \emph{Physicalism, or something near enough}. Princeton University Press.}
\sourceitem{source-kripke-1980}{Kripke, S. A. (1980). \emph{Naming and necessity}. Harvard University Press. (A proper name or natural-kind term is a \emph{rigid designator}, picking out the same object in every possible world in which it exists; reference is fixed causally/by stipulation rather than by a cluster of descriptions; hence identity statements between rigid designators, if true, are necessarily true, though they may be discovered only a posteriori. Originally delivered as lectures at Princeton in 1970 and first published in D. Davidson \& G. Harman (Eds.), \emph{Semantics of natural language}, Reidel, 1972.)}
\sourceitem{source-levine-1983}{Levine — Materialism and Qualia: The Explanatory Gap (1983)}
\sourceitem{source-ci-lewis-1929}{Lewis, C. I. (1929). \emph{Mind and the world-order: Outline of a theory of knowledge}. Charles Scribner's Sons.}
\sourceitem{source-lewis-1979}{Lewis, D. (1979). Attitudes de dicto and de se. \emph{The Philosophical Review}, \emph{88}(4), 513–543. (Reprinted in D. Lewis, \emph{Philosophical papers, Volume I}, pp. 133–159, Oxford University Press, 1983.)}
\sourceitem{source-lewis-1988}{Lewis, D. (1988). What experience teaches. \emph{Proceedings of the Russellian Society}, \emph{13}, 29–35. (Reprinted in W. G. Lycan (Ed.), \emph{Mind and cognition}, Blackwell, 1990; and in D. Lewis, \emph{Papers in metaphysics and epistemology}, Cambridge University Press, 1999.)}
\sourceitem{source-liu-cixin}{Liu, C. (2014). \emph{The three-body problem} (K. Liu, Trans.). Tor Books. (Original work published 2008) (The ``human-formation computer'' episode — the Qin Shi Huang / von Neumann chapter of the in-game civilisation; \textbf{exact chapter and page refs to be confirmed} before any treatise cites it.)}
\sourceitem{source-loar-1990}{Loar, B. (1990). Phenomenal states. \emph{Philosophical Perspectives}, \emph{4}, 81–108. (A revised ``second version'' appears in N. Block, O. Flanagan, \& G. Güzeldere (Eds.), \emph{The nature of consciousness: Philosophical debates}, pp. 597–616, MIT Press, 1997.)}
\sourceitem{source-nagasawa-wager-2017}{Nagasawa, Y., \& Wager, K. (2017). Panpsychism and priority cosmopsychism. In G. Brüntrup \& L. Jaskolla (Eds.), \emph{Panpsychism: Contemporary perspectives} (pp. 113–129). Oxford University Press.}
\sourceitem{source-nagel-1971}{Nagel, T. (1971). Brain bisection and the unity of consciousness. \emph{Synthese}, \emph{22}(3–4), 396–413. \url{https://doi.org/10.1007/BF00413435}}
\sourceitem{source-nagel-1974}{Nagel, T. (1974). What is it like to be a bat? \emph{The Philosophical Review}, \emph{83}(4), 435–450.}
\sourceitem{source-nagel-1979}{Nagel, T. (1979). Panpsychism. In \emph{Mortal questions} (pp. 181–195). Cambridge University Press.}
\sourceitem{source-nemirow-1990}{Nemirow, L. (1990). Physicalism and the cognitive role of acquaintance. In W. G. Lycan (Ed.), \emph{Mind and cognition}. Blackwell.}
\sourceitem{source-papineau-2002}{Papineau, D. (2002). \emph{Thinking about consciousness}. Clarendon Press.}
\sourceitem{source-perry-2001}{Perry, J. (2001). \emph{Knowledge, possibility, and consciousness}. MIT Press. (Based on the 1999 Jean Nicod Lectures.)}
\sourceitem{source-roelofs-2019}{Roelofs, L. (2019). \emph{Combining minds: How to think about composite subjectivity}. Oxford University Press.}
\sourceitem{source-russell-1911}{Russell, B. (1911). Knowledge by acquaintance and knowledge by description. \emph{Proceedings of the Aristotelian Society}, \emph{11}, 108–128. \url{https://doi.org/10.1093/aristotelian/11.1.108}}
\sourceitem{source-russell-1927}{Russell, B. (1927). \emph{The analysis of matter}. Kegan Paul, Trench, Trübner \& Co.}
\sourceitem{source-searle-1980}{Searle, J. R. (1980). Minds, brains, and programs. \emph{Behavioral and Brain Sciences}, \emph{3}(3), 417–457. (The room; the slogan ``syntax is not sufficient for semantics''; and Searle's canvassed replies with his rejoinders — the \textbf{Systems Reply}, the \textbf{Robot Reply}, the \textbf{Brain Simulator Reply}, the \textbf{Combination Reply} — printed with peer commentary.)}
\sourceitem{source-shani-2015}{Shani, I. (2015). Cosmopsychism: A holistic approach to the metaphysics of experience. \emph{Philosophical Papers}, \emph{44}(3), 389–437. \url{https://doi.org/10.1080/05568641.2015.1106709}}
\sourceitem{source-stanley-williamson-2001}{Stanley, J., \& Williamson, T. (2001). Knowing how. \emph{The Journal of Philosophy}, \emph{98}(8), 411–444. \url{https://doi.org/10.2307/2678403}}
\sourceitem{source-strawson-2006}{Strawson, G. (2006). Realistic monism: Why physicalism entails panpsychism. \emph{Journal of Consciousness Studies}, \emph{13}(10–11), 3–31.}
\sourceitem{source-tye-2000}{Tye, M. (2000). \emph{Consciousness, color, and content}. MIT Press.}
\end{sourcelist}
\end{document}
